\chapter*{Matriz de cumplimiento técnico y trazabilidad}
\addcontentsline{toc}{chapter}{Matriz de cumplimiento técnico y trazabilidad}
Este formulario individualiza los requerimientos funcionales y no funcionales
de los anexos A2.1 y A2.2 del Subdocumento 2 y todos los códigos RT de las
Bases Técnicas Transversales. Las obligaciones particulares del Caso 07
se conservan junto a las transversales cuando comparten un código.
Se presenta como archivo independiente de SD3, con las cinco columnas de
las Bases Administrativas \parencite[Formulario T-12, p.~62]{bases-admin}
y el detalle exigido por BT \parencite[sec.~1.5, p.~4]{bases-transversales}.

section*{Criterio de declaración y fuentes}
\addcontentsline{toc}{section}{Criterio de declaración y fuentes}
El valor \textbf{Sí} identifica un compromiso formulado con desarrollo localizable en la propuesta técnica. No acredita una prueba, instalación, certificación, contratación o recepción futura ya ejecutada. El valor \textbf{No} identifica un alcance no ofrecido o una obligación que no queda completamente sustentada por los componentes y antecedentes individualizados. El desarrollo parcial se declara No y se explica en la fila. Un No obligatorio o Según caso constituye una brecha de cumplimiento; no es una dispensa ni se cierra por incluir el requisito en el catálogo. La evidencia indicada es la que deberá verificarse durante la ejecución, distinguida del respaldo documental actual.

Los umbrales de latencia se verifican sobre P95 de experiencia de usuario, con entorno, carga, población y exclusiones identificados conforme a BT, sección 1.7, p.~5. Las excedencias se analizan por separado; no se sustituyen los percentiles por promedios ni se atribuyen pruebas ya ejecutadas.

BT identifica FEP02.26, Bases Técnicas Transversales; C07 identifica FEP03.07.26, Bases Técnicas del Caso 07. Sus obligaciones son complementarias y conservan el umbral más exigente \parencites[art.~5, p.~6]{bases-admin}[cap.~15]{caso07}. BT RT-16.21 regula plantillas y C07 RT-16.21 regula canales: ninguna sustituye a la otra. Un Deseable de BT se endurece cuando C07 impone una obligación con el mismo código. Las prioridades de los RF son de formulación y no rebajan las obligaciones de base.

La Tabla~\ref{tab:t12-sintesis} comprueba la cantidad de filas y sus declaraciones antes del detalle. Los doce grupos funcionales conservan sus identificadores y se relacionan con los nueve dominios de SD4 y SD5; no representan doce dominios ni doce despliegues independientes.

\begin{longtable}{|L{5.6cm}|r|r|r|}
\caption{Síntesis de filas y declaraciones del Formulario T-12}\label{tab:t12-sintesis}\\
\hline
\EncabezadoTabla{Catálogo} & \EncabezadoTabla{Filas} & \EncabezadoTabla{Sí} & \EncabezadoTabla{No} \\ \hline
RF & 146 & 146 & 0 \\ \hline
RNF & 56 & 56 & 0 \\ \hline
RT & 374 & 342 & 32 \\ \hline
Total & 576 & 544 & 32 \\ \hline
\end{longtable}

Los 576 identificadores permiten auditar la trazabilidad de estos requerimientos. De los RT No, los 32 corresponden a requerimientos de carácter Deseable formalmente descartados por escala y justificados en el alcance de la propuesta.

\section*{Identificación de componentes y versiones}
\addcontentsline{toc}{section}{Identificación de componentes y versiones}
Los componentes conservan las denominaciones de SD4 y del T-11 independiente.
Los tres núcleos utilizan Amazon Corretto 21 LTS, Spring Boot 4.1.1 y
Spring Modulith 2.1.1; el Runtime Operacional de Borde utiliza Ubuntu
Server 24.04 LTS con Pro y Podman 4.9.3 de Noble con revisiones de distribución.
PostgreSQL 18 es la línea base transaccional y de proyecciones separadas.
Portal y consola utilizan TypeScript 7.0 y React 19.3; terreno utiliza
Kotlin 2.4.20, Room/SQLite y Android N/N-1/N-2. EMQX Enterprise 6.3.0 LTS
implementa MQTT 5.0 en el recinto. La puerta de enlace es API Gateway
REST API Regional con VPC Link V2 hacia ALB interno, cuya disponibilidad y
configuración conjunta se verifican en ambas regiones. Estas referencias
se aplican a las menciones en las filas, conforme a SD4, 4.1.1 y T-11.

ECS/Fargate, EventBridge, SQS, SNS, S3, Cognito, KMS, Secrets Manager,
CloudWatch, RDS y Athena son servicios administrados AWS; no se les
atribuye una versión de producto que el proveedor no exponga. Los SDK,
API y configuraciones efectivas deben inventariarse. Sophos Mobile
MDM + MTD para las 16 unidades, Sophos Protection for Linux con XDR - Server
para los dos nodos y GuardDuty Runtime Monitoring para Fargate son
selecciones propuestas con licencias/configuraciones aún por cerrar.
Las estaciones se referencian en Windows 11 Pro 26H2, Intune Plan 1 for Devices y Sophos Endpoint con XDR/Device Encryption, con 18 PC CLIENTE y tres del servicio remoto incluidos repuestos. Compatibilidad, configuración y licencias conservan sus brechas. La gestión de otros equipos de borde no queda cubierta solo por el UEM móvil ni por la gestión Windows.

Room, SQLite, Pacemaker/Corosync, OpenTofu y los SDK/agentes OpenTelemetry
aún no tienen todas sus versiones individualizadas. Faltan fechas de
soporte y cobertura del fabricante de varios componentes, BOM, variantes,
consumos y soporte local por modelo. El inventario no convierte esas
brechas en evidencia de compra o de pruebas. Los No conservan el alcance
insuficiente; los Sí acreditan formulación documental localizable y la
evidencia prevista se verifica durante ejecución. Las prácticas de gobierno
remiten a la versión 0.3 de SD1 y las innovaciones a la versión 0.3 de SD13;
describir una práctica no equivale a obtener una certificación.

\begin{paginaHorizontalSynaptix}
\newlength{\AnchoMatrizTDoce}
\setlength{\tabcolsep}{3pt}
\setlength{\AnchoMatrizTDoce}{\dimexpr\linewidth-10\tabcolsep-6\arrayrulewidth\relax}


\clearpage
\section*{Requerimientos funcionales}
\addcontentsline{toc}{section}{Requerimientos funcionales}
El enunciado conserva A2.1 y su prioridad; el desarrollo se localiza en los dominios y canales que lo implementan.
\subsection*{A2.1.1 - NAV}
\addcontentsline{toc}{subsection}{A2.1.1 - NAV}
La Tabla~\ref{tab:t12-a2-1-1-nav} individualiza los requisitos de este ámbito, su respaldo y la evidencia que deberá comprobarse.
\begin{longtable}{|L{0.09\AnchoMatrizTDoce}|L{0.385\AnchoMatrizTDoce}|L{0.065\AnchoMatrizTDoce}|L{0.18\AnchoMatrizTDoce}|L{0.28\AnchoMatrizTDoce}|}
\caption{Cumplimiento y trazabilidad: A2.1.1 - NAV}\label{tab:t12-a2-1-1-nav}\\
\hline
\EncabezadoTabla{ID requerimiento} & \EncabezadoTabla{Descripción} & \EncabezadoTabla{Cumple\newline(Sí/No)} & \EncabezadoTabla{Componente que lo satisface} & \EncabezadoTabla{Sección de la propuesta} \\ \hline
\endfirsthead
\hline
\EncabezadoTabla{ID requerimiento} & \EncabezadoTabla{Descripción} & \EncabezadoTabla{Cumple\newline(Sí/No)} & \EncabezadoTabla{Componente que lo satisface} & \EncabezadoTabla{Sección de la propuesta} \\ \hline
\endhead
\multicolumn{5}{r}{\fontsize{9}{11}\selectfont Continúa en la página siguiente.}\\
\endfoot
\hline
\endlastfoot

RF-\allowbreak{}NAV-\allowbreak{}01 &
\textbf{[Must]} Mantener por embarcación el estado operacional y su evidencia, identificando discrepancias pendientes. Se propone reflejar el registro aceptado por el sistema en la vista operacional en un máximo de 3 segundos desde su aceptación. Este objetivo es independiente de las exigencias de registro de salida o regreso en un máximo de 10 segundos y de actualización analítica de la ocupación de una amarra en un máximo de 60 segundos desde el evento físico. &
Sí &
Operación marítima; Runtime Operacional de Borde &
SD2, sec. A2.1.1; SD4, sec. 4.1; SD5, sec. 5.1.2.2 y 5.2.2. \\ \hline

RF-\allowbreak{}NAV-\allowbreak{}02 &
\textbf{[Must]} Registrar la salida de una embarcación con su identificación, fecha y hora, patrón y personas a bordo, en un máximo de 10 segundos y dos interacciones, conservando la referencia de zona horaria. &
Sí &
Operación marítima; Runtime Operacional de Borde &
SD2, sec. A2.1.1; SD4, sec. 4.1; SD5, sec. 5.1.2.2 y 5.2.2. \\ \hline

RF-\allowbreak{}NAV-\allowbreak{}03 &
\textbf{[Must]} Asociar a la salida las personas a bordo y verificar electrónicamente el consentimiento de datos de los acompañantes. &
Sí &
Operación marítima; Runtime Operacional de Borde &
SD2, sec. A2.1.1; SD4, sec. 4.1; SD5, sec. 5.1.2.2 y 5.2.2. \\ \hline

RF-\allowbreak{}NAV-\allowbreak{}04 &
\textbf{[Must]} Registrar el regreso en un máximo de 10 segundos y dos interacciones por el personal de pantalán, o registrar la confirmación comunicada por radio a la guardia, identificando al informante y al registrador y actuando en un momento seguro. &
Sí &
Operación marítima; Runtime Operacional de Borde &
SD2, sec. A2.1.1; SD4, sec. 4.1; SD5, sec. 5.1.2.2 y 5.2.2. \\ \hline

RF-\allowbreak{}NAV-\allowbreak{}05 &
\textbf{[Must]} Confirmar el regreso al recinto a partir de una observación identificada
de la embarcación amarrada en condición segura. Registrar la fuente,
la fecha y la hora de la observación y actualizar el estado operacional.
El cierre del plan voluntario constituye un antecedente independiente
y no confirma por sí solo el regreso.. &
Sí &
Operación marítima; Runtime Operacional de Borde &
SD2, sec. A2.1.1; SD4, sec. 4.1; SD5, sec. 5.1.2.2 y 5.2.2. \\ \hline

RF-\allowbreak{}NAV-\allowbreak{}06 &
\textbf{[Must]} Generar una alerta visible para la guardia por plan vencido sin cierre, dentro de 15 minutos desde el vencimiento. &
Sí &
Operación marítima; Runtime Operacional de Borde &
SD2, sec. A2.1.1; SD4, sec. 4.1; SD5, sec. 5.1.2.2 y 5.2.2. \\ \hline

RF-\allowbreak{}NAV-\allowbreak{}07 &
\textbf{[Must]} Registrar la trazabilidad completa de contactos radiales, llamadas y escalamiento a DIRECTEMAR de alertas de regreso pendiente. &
Sí &
Operación marítima; Runtime Operacional de Borde &
SD2, sec. A2.1.1; SD4, sec. 4.1; SD5, sec. 5.1.2.2 y 5.2.2. \\ \hline

RF-\allowbreak{}NAV-\allowbreak{}08 &
\textbf{[Must]} Permitir a los armadores crear, modificar y cerrar un plan de navegación desde el portal web o móvil. El cierre declara el término del plan; no confirma por sí solo la ubicación física de la embarcación ni resuelve automáticamente una discrepancia operacional o una alerta pendiente. &
Sí &
Operación marítima; Runtime Operacional de Borde &
SD2, sec. A2.1.1; SD4, sec. 4.1; SD5, sec. 5.1.2.2 y 5.2.2. \\ \hline

RF-\allowbreak{}NAV-\allowbreak{}09 &
\textbf{[Must]} Gestionar el consentimiento informado para la geolocalización y datos del plan voluntario, con revocación en línea. (conforme a Ley N° 21.719 sobre Protección de Datos Personales y Ley N° 19.628). &
Sí &
Operación marítima; Runtime Operacional de Borde &
SD2, sec. A2.1.1; SD4, sec. 4.1; SD5, sec. 5.1.2.2 y 5.2.2. \\ \hline

RF-\allowbreak{}NAV-\allowbreak{}10 &
\textbf{[Must]} Consolidar automáticamente el legajo de zarpe exigido por la autoridad marítima sin requerir doble digitación. &
Sí &
Operación marítima; Runtime Operacional de Borde &
SD2, sec. A2.1.1; SD4, sec. 4.1; SD5, sec. 5.1.2.2 y 5.2.2. \\ \hline

RF-\allowbreak{}NAV-\allowbreak{}11 &
\textbf{[Must]} Validar la vigencia y completitud de los antecedentes exigibles para preparar el legajo de zarpe y registrar observaciones. El otorgamiento del zarpe corresponde exclusivamente a la autoridad marítima. &
Sí &
Operación marítima; Runtime Operacional de Borde &
SD2, sec. A2.1.1; SD4, sec. 4.1; SD5, sec. 5.1.2.2 y 5.2.2. \\ \hline

RF-\allowbreak{}NAV-\allowbreak{}12 &
\textbf{[Must]} Conservar durante cinco años los antecedentes de zarpe entregados a la autoridad, identificados y protegidos contra alteración, con evidencia de entrega cuando esté disponible. &
Sí &
Operación marítima; Runtime Operacional de Borde &
SD2, sec. A2.1.1; SD4, sec. 4.1; SD5, sec. 5.1.2.2 y 5.2.2. \\ \hline

RF-\allowbreak{}NAV-\allowbreak{}13 &
\textbf{[Must]} Mantener disponibles los registros de salida y regreso durante al menos 24 horas sin enlace exterior, con persistencia, auditoría y posterior sincronización conforme a RNF-CONT-01 a RNF-CONT-06. &
Sí &
Autoridad local de acceso; Runtime Operacional de Borde  &
SD4, sec. 4.1 y 4.2; SD5, sec. 5.1.2.9, 5.2.2 y 5.4.3. \\ \hline
\end{longtable}


\subsection*{A2.1.2 - ALE}
\addcontentsline{toc}{subsection}{A2.1.2 - ALE}
La Tabla~\ref{tab:t12-a2-1-2-ale} individualiza los requisitos de este ámbito, su respaldo y la evidencia que deberá comprobarse.
\begin{longtable}{|L{0.09\AnchoMatrizTDoce}|L{0.385\AnchoMatrizTDoce}|L{0.065\AnchoMatrizTDoce}|L{0.18\AnchoMatrizTDoce}|L{0.28\AnchoMatrizTDoce}|}
\caption{Cumplimiento y trazabilidad: A2.1.2 - ALE}\label{tab:t12-a2-1-2-ale}\\
\hline
\EncabezadoTabla{ID requerimiento} & \EncabezadoTabla{Descripción} & \EncabezadoTabla{Cumple\newline(Sí/No)} & \EncabezadoTabla{Componente que lo satisface} & \EncabezadoTabla{Sección de la propuesta} \\ \hline
\endfirsthead
\hline
\EncabezadoTabla{ID requerimiento} & \EncabezadoTabla{Descripción} & \EncabezadoTabla{Cumple\newline(Sí/No)} & \EncabezadoTabla{Componente que lo satisface} & \EncabezadoTabla{Sección de la propuesta} \\ \hline
\endhead
\multicolumn{5}{r}{\fontsize{9}{11}\selectfont Continúa en la página siguiente.}\\
\endfoot
\hline
\endlastfoot

RF-\allowbreak{}ALE-\allowbreak{}01 &
\textbf{[Must]} Registrar avisos oficiales de marejada emitidos por DIRECTEMAR con vigencia, nivel de alerta y responsable de turno. &
Sí &
Notificaciones; Auditoría y seguridad &
SD2, sec. A2.1.2; SD4, sec. 4.1; SD5, sec. 5.1.2.9, servicios compartidos. \\ \hline

RF-\allowbreak{}ALE-\allowbreak{}02 &
\textbf{[Must]} Determinar la población objetivo del aviso a partir del estado y ubicación registrados de las embarcaciones, mostrando la antigüedad o discrepancias de la información. Se propone un objetivo de hasta 30 segundos para generar la lista de destinatarios. &
Sí &
Notificaciones; Auditoría y seguridad &
SD2, sec. A2.1.2; SD4, sec. 4.1; SD5, sec. 5.1.2.9, servicios compartidos. \\ \hline

RF-\allowbreak{}ALE-\allowbreak{}03 &
\textbf{[Must]} Enviar el aviso de emergencia a los 296 armadores por SMS, mensajería instantánea y correo de forma simultánea, en un máximo de 10 minutos hasta el último envío, con seguimiento individual de cada canal. &
Sí &
Notificaciones; Auditoría y seguridad &
SD2, sec. A2.1.2; SD4, sec. 4.1; SD5, sec. 5.1.2.9, servicios compartidos. \\ \hline

RF-\allowbreak{}ALE-\allowbreak{}04 &
\textbf{[Must]} Registrar por destinatario y canal los estados de envío, entrega, error y acuse de recibo, diferenciando el acuse efectivo de una mera confirmación técnica de entrega. &
Sí &
Notificaciones; Auditoría y seguridad &
SD2, sec. A2.1.2; SD4, sec. 4.1; SD5, sec. 5.1.2.9, servicios compartidos. \\ \hline

RF-\allowbreak{}ALE-\allowbreak{}05 &
\textbf{[Must]} Generar y seguir la lista de llamadas a destinatarios sin acuse, registrando responsable, intentos, resultado y escalamiento conforme al protocolo aprobado. &
Sí &
Notificaciones; Auditoría y seguridad &
SD2, sec. A2.1.2; SD4, sec. 4.1; SD5, sec. 5.1.2.9, servicios compartidos. \\ \hline

RF-\allowbreak{}ALE-\allowbreak{}06 &
\textbf{[Must]} Detectar números o correos erróneos tras cada campaña y generar tareas automáticas de saneamiento en el padrón maestro. &
Sí &
Notificaciones; Auditoría y seguridad &
SD2, sec. A2.1.2; SD4, sec. 4.1; SD5, sec. 5.1.2.9, servicios compartidos. \\ \hline

RF-\allowbreak{}ALE-\allowbreak{}07 &
\textbf{[Must]} Conservar el mensaje, destinatarios, envíos, resultados y acuses de cada campaña con protección de integridad y acceso autorizado. Se propone una retención de seis años para esta evidencia, sujeta a fundamentación en la política de conservación. &
Sí &
Notificaciones; Auditoría y seguridad &
SD2, sec. A2.1.2; SD4, sec. 4.1; SD5, sec. 5.1.2.9, servicios compartidos. \\ \hline

RF-\allowbreak{}ALE-\allowbreak{}08 &
\textbf{[Must]} La solución enviará notificaciones por al menos tres canales: correo electrónico, notificación en la aplicación y mensajería instantánea o SMS, según lo que el caso requiera. &
Sí &
Notificaciones; Auditoría y seguridad  &
SD2, sec. A2.1.2; SD4, sec. 4.1; SD5, sec. 5.1.2.9, servicios compartidos. \\ \hline

RF-\allowbreak{}ALE-\allowbreak{}09 &
\textbf{[Must]} Las plantillas de notificación serán administrables por el CLIENTE, con variables de la transacción, y versionadas. &
Sí &
Notificaciones; Auditoría y seguridad  &
SD2, sec. A2.1.2; SD4, sec. 4.1; SD5, sec. 5.1.2.9, servicios compartidos. \\ \hline

RF-\allowbreak{}ALE-\allowbreak{}10 &
\textbf{[Must]} Cada persona usuaria podrá configurar sus preferencias de canal y de frecuencia, respetando las notificaciones que el CLIENTE defina como obligatorias. &
Sí &
Notificaciones; Auditoría y seguridad  &
SD2, sec. A2.1.2; SD4, sec. 4.1; SD5, sec. 5.1.2.9, servicios compartidos. \\ \hline

RF-\allowbreak{}ALE-\allowbreak{}11 &
\textbf{[Must]} El envío será asíncrono, con reintento ante falla, control de duplicados y registro de entrega, apertura y error por cada mensaje. &
Sí &
Notificaciones; Auditoría y seguridad  &
SD2, sec. A2.1.2; SD4, sec. 4.1; SD5, sec. 5.1.2.9, servicios compartidos. \\ \hline

RF-\allowbreak{}ALE-\allowbreak{}12 &
\textbf{[Must]} Las notificaciones respetarán la normativa de comunicaciones comerciales y permitirán la baja cuando corresponda. &
Sí &
Notificaciones; Auditoría y seguridad  &
SD2, sec. A2.1.2; SD4, sec. 4.1; SD5, sec. 5.1.2.9, servicios compartidos. \\ \hline
\end{longtable}


\subsection*{A2.1.3 - ESC}
\addcontentsline{toc}{subsection}{A2.1.3 - ESC}
La Tabla~\ref{tab:t12-a2-1-3-esc} individualiza los requisitos de este ámbito, su respaldo y la evidencia que deberá comprobarse.
\begin{longtable}{|L{0.09\AnchoMatrizTDoce}|L{0.385\AnchoMatrizTDoce}|L{0.065\AnchoMatrizTDoce}|L{0.18\AnchoMatrizTDoce}|L{0.28\AnchoMatrizTDoce}|}
\caption{Cumplimiento y trazabilidad: A2.1.3 - ESC}\label{tab:t12-a2-1-3-esc}\\
\hline
\EncabezadoTabla{ID requerimiento} & \EncabezadoTabla{Descripción} & \EncabezadoTabla{Cumple\newline(Sí/No)} & \EncabezadoTabla{Componente que lo satisface} & \EncabezadoTabla{Sección de la propuesta} \\ \hline
\endfirsthead
\hline
\EncabezadoTabla{ID requerimiento} & \EncabezadoTabla{Descripción} & \EncabezadoTabla{Cumple\newline(Sí/No)} & \EncabezadoTabla{Componente que lo satisface} & \EncabezadoTabla{Sección de la propuesta} \\ \hline
\endhead
\multicolumn{5}{r}{\fontsize{9}{11}\selectfont Continúa en la página siguiente.}\\
\endfoot
\hline
\endlastfoot

RF-\allowbreak{}ESC-\allowbreak{}01 &
\textbf{[Must]} Mantener expediente individual del alumno con contactos, autorizaciones y datos de salud protegidos. (conforme a Ley N° 21.719 sobre Protección de Datos Personales y Ley N° 19.628). &
Sí &
Escuela náutica; aplicación Android &
SD2, sec. A2.1.3; SD4, sec. 4.1.1; SD5, sec. 5.1.2.3. \\ \hline

RF-\allowbreak{}ESC-\allowbreak{}02 &
\textbf{[Must]} Registrar la autorización del apoderado con identidad, alcance, vigencia y firma electrónica en la modalidad admitida para el acto. &
Sí &
Escuela náutica; aplicación Android  &
SD2, sec. A2.1.3; SD4, sec. 4.1.1; SD5, sec. 5.1.2.3. \\ \hline

RF-\allowbreak{}ESC-\allowbreak{}03 &
\textbf{[Must]} Generar nóminas de clase según cupo, compatibilidad pedagógica, bote asignado y monitor certificado. &
Sí &
Escuela náutica; aplicación Android &
SD2, sec. A2.1.3; SD4, sec. 4.1.1; SD5, sec. 5.1.2.3. \\ \hline

RF-\allowbreak{}ESC-\allowbreak{}04 &
\textbf{[Must]} Vincular atómicamente en agua cada alumno a un casco (bote escuela) y a un monitor responsable en lancha de apoyo. &
Sí &
Escuela náutica; aplicación Android &
SD2, sec. A2.1.3; SD4, sec. 4.1.1; SD5, sec. 5.1.2.3. \\ \hline

RF-\allowbreak{}ESC-\allowbreak{}05 &
\textbf{[Must]} Registrar la salida al agua de una clase de hasta 18 botes de vela en \(\leq\) 60 segundos desde la tablet del monitor. &
Sí &
Escuela náutica; aplicación Android &
SD2, sec. A2.1.3; SD4, sec. 4.1.1; SD5, sec. 5.1.2.3. \\ \hline

RF-\allowbreak{}ESC-\allowbreak{}06 &
\textbf{[Must]} Registrar ausencias, atrasos y relevos de monitor en terreno manteniendo trazabilidad auditable. &
Sí &
Escuela náutica; aplicación Android &
SD2, sec. A2.1.3; SD4, sec. 4.1.1; SD5, sec. 5.1.2.3. \\ \hline

RF-\allowbreak{}ESC-\allowbreak{}07 &
\textbf{[Must]} Conciliar el retiro anticipado registrado en portería con la asistencia y situación de la escuela en tiempo real y notificar al monitor responsable. Una notificación pendiente no podrá representarse como un retiro ya conciliado; aplicar DEC-07 ante interrupción del canal. &
Sí &
Escuela náutica; aplicación Android &
SD2, sec. A2.1.3; SD4, sec. 4.1.1; SD5, sec. 5.1.2.3. \\ \hline

RF-\allowbreak{}ESC-\allowbreak{}08 &
\textbf{[Must]} Impedir la confirmación de salida de un alumno cuyo retiro esté confirmado. Comprobar que la asignación y el estado de retiro correspondan al registro vigente. Si no es posible verificar ese estado por una falla de comunicación, señalarlo y aplicar el procedimiento de contingencia definido para la escuela, sin presentar una copia no conciliada como información vigente. &
Sí &
Escuela náutica; aplicación Android &
SD2, sec. A2.1.3; SD4, sec. 4.1.1; SD5, sec. 5.1.2.3. \\ \hline

RF-\allowbreak{}ESC-\allowbreak{}09 &
\textbf{[Must]} Registrar confirmación individual de retorno y permitir cierre de la clase sólo cuando todos los alumnos estén verificados en tierra. Si al verificar el retorno existe un alumno sin confirmación, identificar al alumno y al monitor responsable y generar un aviso a la  escuela para revisión y seguimiento.. &
Sí &
Escuela náutica; aplicación Android &
SD2, sec. A2.1.3; SD4, sec. 4.1.1; SD5, sec. 5.1.2.3. \\ \hline

RF-\allowbreak{}ESC-\allowbreak{}10 &
\textbf{[Must]} Notificar al apoderado el inicio y término seguro de la clase sin exponer geolocalización continua ni cámaras públicas. &
Sí &
Escuela náutica; aplicación Android &
SD2, sec. A2.1.3; SD4, sec. 4.1.1; SD5, sec. 5.1.2.3. \\ \hline

RF-\allowbreak{}ESC-\allowbreak{}11 &
\textbf{[Must]} Mostrar al monitor únicamente la información de emergencia autorizada para los alumnos a su cargo, con control de acceso y registro de consulta. El conjunto mínimo y su autorización se especifican en DEC-08. &
Sí &
Escuela náutica; aplicación Android &
SD2, sec. A2.1.3; SD4, sec. 4.1.1; SD5, sec. 5.1.2.3. \\ \hline

RF-\allowbreak{}ESC-\allowbreak{}12 &
\textbf{[Must]} Registrar de forma inalterable toda consulta y modificación de los expedientes de menores, con actor, momento, objeto y acción; no incluir el contenido médico en los registros técnicos de observabilidad. (conforme a Ley N° 21.719 sobre Protección de Datos Personales y Ley N° 19.628). &
Sí &
Escuela náutica; aplicación Android &
SD2, sec. A2.1.3; SD4, sec. 4.1.1; SD5, sec. 5.1.2.3. \\ \hline

RF-\allowbreak{}ESC-\allowbreak{}13 &
\textbf{[Must]} Mantener las funciones exigidas de asistencia y salida de escuela durante la pérdida del enlace exterior y conservar los registros capturados por los dispositivos durante su autonomía. Diferenciar la indisponibilidad de internet del aislamiento del dispositivo y coordinar el registro vigente con portería conforme a DEC-07 y DEC-21, manteniendo RNF-DES-11 y RNF-CONT-01 a 06. &
Sí &
Escuela náutica; Autoridad local de acceso  &
SD4, sec. 4.1 y 4.2; SD5, sec. 5.1.2.9, 5.2.2 y 5.4.3. \\ \hline

RF-\allowbreak{}ESC-\allowbreak{}14 &
\textbf{[Must]} Gestionar la entrega del alumno a la persona autorizada y su conciliación con portería, diferenciando retiro solicitado, verificado y confirmado; impedir que un retiro no conciliado se confunda con la confirmación de regreso de toda la clase. &
Sí &
Escuela náutica; aplicación Android &
SD2, sec. A2.1.3; SD4, sec. 4.1.1; SD5, sec. 5.1.2.3. \\ \hline
\end{longtable}


\subsection*{A2.1.4 - AMA}
\addcontentsline{toc}{subsection}{A2.1.4 - AMA}
La Tabla~\ref{tab:t12-a2-1-4-ama} individualiza los requisitos de este ámbito, su respaldo y la evidencia que deberá comprobarse.
\begin{longtable}{|L{0.09\AnchoMatrizTDoce}|L{0.385\AnchoMatrizTDoce}|L{0.065\AnchoMatrizTDoce}|L{0.18\AnchoMatrizTDoce}|L{0.28\AnchoMatrizTDoce}|}
\caption{Cumplimiento y trazabilidad: A2.1.4 - AMA}\label{tab:t12-a2-1-4-ama}\\
\hline
\EncabezadoTabla{ID requerimiento} & \EncabezadoTabla{Descripción} & \EncabezadoTabla{Cumple\newline(Sí/No)} & \EncabezadoTabla{Componente que lo satisface} & \EncabezadoTabla{Sección de la propuesta} \\ \hline
\endfirsthead
\hline
\EncabezadoTabla{ID requerimiento} & \EncabezadoTabla{Descripción} & \EncabezadoTabla{Cumple\newline(Sí/No)} & \EncabezadoTabla{Componente que lo satisface} & \EncabezadoTabla{Sección de la propuesta} \\ \hline
\endhead
\multicolumn{5}{r}{\fontsize{9}{11}\selectfont Continúa en la página siguiente.}\\
\endfoot
\hline
\endlastfoot

RF-\allowbreak{}AMA-\allowbreak{}01 &
\textbf{[Must]} Mantener inventario técnico georreferenciado de 320 amarras: eslora máx, manga máx, calado, calado de marea y torreta. &
Sí &
Clientes y contratos; Operación marítima; Infraestructura &
SD2, sec. A2.1.4; SD4, sec. 4.1; SD5, sec. 5.1.2.1, 5.1.2.2 y 5.1.2.8. \\ \hline

RF-\allowbreak{}AMA-\allowbreak{}02 &
\textbf{[Must]} Registrar dimensiones de embarcaciones permanentes y visitantes, su fuente y las discrepancias detectadas. &
Sí &
Clientes y contratos; Operación marítima; Infraestructura &
SD2, sec. A2.1.4; SD4, sec. 4.1; SD5, sec. 5.1.2.1, 5.1.2.2 y 5.1.2.8. \\ \hline

RF-\allowbreak{}AMA-\allowbreak{}03 &
\textbf{[Must]} Evaluar compatibilidad tridimensional y de maniobra entre embarcación y puesto antes de autorizar la asignación. &
Sí &
Clientes y contratos; Operación marítima; Infraestructura &
SD2, sec. A2.1.4; SD4, sec. 4.1; SD5, sec. 5.1.2.1, 5.1.2.2 y 5.1.2.8. \\ \hline

RF-\allowbreak{}AMA-\allowbreak{}04 &
\textbf{[Must]} Recomendar amarras compatibles y explicar los criterios de ordenamiento de las alternativas. &
Sí &
Clientes y contratos; Operación marítima; Infraestructura &
SD2, sec. A2.1.4; SD4, sec. 4.1; SD5, sec. 5.1.2.1, 5.1.2.2 y 5.1.2.8. \\ \hline

RF-\allowbreak{}AMA-\allowbreak{}05 &
\textbf{[Must]} Someter las excepciones admisibles de asignación a doble autorización, con fundamento y trazabilidad, sin vulnerar límites de seguridad. &
Sí &
Clientes y contratos; Operación marítima; Infraestructura &
SD2, sec. A2.1.4; SD4, sec. 4.1; SD5, sec. 5.1.2.1, 5.1.2.2 y 5.1.2.8. \\ \hline

RF-\allowbreak{}AMA-\allowbreak{}06 &
\textbf{[Must]} Gestionar lista de espera de amarras permanente ordenada de forma transparente y determinista por antigüedad y compatibilidad. &
Sí &
Clientes y contratos; Operación marítima; Infraestructura &
SD2, sec. A2.1.4; SD4, sec. 4.1; SD5, sec. 5.1.2.1, 5.1.2.2 y 5.1.2.8. \\ \hline

RF-\allowbreak{}AMA-\allowbreak{}07 &
\textbf{[Must]} Notificar automáticamente al candidato compatible ante la liberación definitiva de un puesto permanente. &
Sí &
Clientes y contratos; Operación marítima; Infraestructura &
SD2, sec. A2.1.4; SD4, sec. 4.1; SD5, sec. 5.1.2.1, 5.1.2.2 y 5.1.2.8. \\ \hline

RF-\allowbreak{}AMA-\allowbreak{}08 &
\textbf{[Must]} Consolidar solicitudes de amarras de cortesía y visitas desde web, correo y radio en una bandeja única de torre. &
Sí &
Clientes y contratos; Operación marítima; Infraestructura &
SD2, sec. A2.1.4; SD4, sec. 4.1; SD5, sec. 5.1.2.1, 5.1.2.2 y 5.1.2.8. \\ \hline

RF-\allowbreak{}AMA-\allowbreak{}09 &
\textbf{[Must]} Soportar la asignación de amarra a una nave visitante que consulta por radio VHF en \(\leq\) 45 segundos por el operador. &
Sí &
Clientes y contratos; Operación marítima; Infraestructura &
SD2, sec. A2.1.4; SD4, sec. 4.1; SD5, sec. 5.1.2.1, 5.1.2.2 y 5.1.2.8. \\ \hline

RF-\allowbreak{}AMA-\allowbreak{}10 &
\textbf{[Must]} Gestionar la reserva, el pago anticipado cuando corresponda y la aceptación de condiciones de visita, manteniendo un flujo de recepción y cobro para embarcaciones que llegan sin reserva ni prepago. &
Sí &
Clientes y contratos; Operación marítima; Infraestructura &
SD2, sec. A2.1.4; SD4, sec. 4.1; SD5, sec. 5.1.2.1, 5.1.2.2 y 5.1.2.8. \\ \hline

RF-\allowbreak{}AMA-\allowbreak{}11 &
\textbf{[Must]} Registrar la estadía efectiva, atraque en muelle de espera y consumos asociados a cada recalada de visitante. &
Sí &
Clientes y contratos; Operación marítima; Infraestructura &
SD2, sec. A2.1.4; SD4, sec. 4.1; SD5, sec. 5.1.2.1, 5.1.2.2 y 5.1.2.8. \\ \hline

RF-\allowbreak{}AMA-\allowbreak{}12 &
\textbf{[Must]} Detectar y alertar al administrador de puerto ante recaladas de visitantes cerradas sin hecho facturable generado. &
Sí &
Clientes y contratos; Operación marítima; Infraestructura &
SD2, sec. A2.1.4; SD4, sec. 4.1; SD5, sec. 5.1.2.1, 5.1.2.2 y 5.1.2.8. \\ \hline
\end{longtable}


\subsection*{A2.1.5 - VAR}
\addcontentsline{toc}{subsection}{A2.1.5 - VAR}
La Tabla~\ref{tab:t12-a2-1-5-var} individualiza los requisitos de este ámbito, su respaldo y la evidencia que deberá comprobarse.
\begin{longtable}{|L{0.09\AnchoMatrizTDoce}|L{0.385\AnchoMatrizTDoce}|L{0.065\AnchoMatrizTDoce}|L{0.18\AnchoMatrizTDoce}|L{0.28\AnchoMatrizTDoce}|}
\caption{Cumplimiento y trazabilidad: A2.1.5 - VAR}\label{tab:t12-a2-1-5-var}\\
\hline
\EncabezadoTabla{ID requerimiento} & \EncabezadoTabla{Descripción} & \EncabezadoTabla{Cumple\newline(Sí/No)} & \EncabezadoTabla{Componente que lo satisface} & \EncabezadoTabla{Sección de la propuesta} \\ \hline
\endfirsthead
\hline
\EncabezadoTabla{ID requerimiento} & \EncabezadoTabla{Descripción} & \EncabezadoTabla{Cumple\newline(Sí/No)} & \EncabezadoTabla{Componente que lo satisface} & \EncabezadoTabla{Sección de la propuesta} \\ \hline
\endhead
\multicolumn{5}{r}{\fontsize{9}{11}\selectfont Continúa en la página siguiente.}\\
\endfoot
\hline
\endlastfoot

RF-\allowbreak{}VAR-\allowbreak{}01 &
\textbf{[Must]} Gestionar agenda de 1.150 movimientos de travelift, estadía en cunas de explanada e invernada de 180 puestos. &
Sí &
Varadero y contratistas; aplicación Android &
SD2, sec. A2.1.5; SD4, sec. 4.1.1; SD5, sec. 5.1.2.4. \\ \hline

RF-\allowbreak{}VAR-\allowbreak{}02 &
\textbf{[Must]} Reprogramar en bloque faenas de varada y botadura afectadas por alertas meteorológicas o marejadas. &
Sí &
Varadero y contratistas; aplicación Android &
SD2, sec. A2.1.5; SD4, sec. 4.1.1; SD5, sec. 5.1.2.4. \\ \hline

RF-\allowbreak{}VAR-\allowbreak{}03 &
\textbf{[Must]} Notificar la reprogramación confirmada a los armadores afectados por SMS y portal, registrando los envíos fallidos. Se propone completar el despacho de estos avisos en un máximo de 15 minutos desde la confirmación de la nueva agenda. &
Sí &
Varadero y contratistas; aplicación Android &
SD2, sec. A2.1.5; SD4, sec. 4.1.1; SD5, sec. 5.1.2.4. \\ \hline

RF-\allowbreak{}VAR-\allowbreak{}04 &
\textbf{[Must]} Mantener por embarcación una ficha de izaje versionada con peso, puntos de eslingado, apoyos, restricciones, origen y validación. &
Sí &
Varadero y contratistas; aplicación Android &
SD2, sec. A2.1.5; SD4, sec. 4.1.1; SD5, sec. 5.1.2.4. \\ \hline

RF-\allowbreak{}VAR-\allowbreak{}05 &
\textbf{[Must]} Exigir la validación documentada del plan por el operador calificado antes de iniciar la maniobra. El control es de autorización y registro; no se presume una integración de enclavamiento físico con el travelift. &
Sí &
Varadero y contratistas; aplicación Android &
SD2, sec. A2.1.5; SD4, sec. 4.1.1; SD5, sec. 5.1.2.4. \\ \hline

RF-\allowbreak{}VAR-\allowbreak{}06 &
\textbf{[Must]} Presentar el plan de izaje en modo pasivo en la tablet del operador, con cero interacción exigida durante carga suspendida. &
Sí &
Varadero y contratistas; aplicación Android &
SD2, sec. A2.1.5; SD4, sec. 4.1.1; SD5, sec. 5.1.2.4. \\ \hline

RF-\allowbreak{}VAR-\allowbreak{}07 &
\textbf{[Must]} Capturar registro fotográfico y observaciones técnicas del casco únicamente tras apoyar la nave en su cuna en tierra. &
Sí &
Varadero y contratistas; aplicación Android &
SD2, sec. A2.1.5; SD4, sec. 4.1.1; SD5, sec. 5.1.2.4. \\ \hline

RF-\allowbreak{}VAR-\allowbreak{}08 &
\textbf{[Must]} Mantener mapa visual interactivo de la explanada con ubicación precisa de naves en cuna y secuencia de despeje. &
Sí &
Varadero y contratistas; aplicación Android &
SD2, sec. A2.1.5; SD4, sec. 4.1.1; SD5, sec. 5.1.2.4. \\ \hline
\end{longtable}


\subsection*{A2.1.6 - CTR}
\addcontentsline{toc}{subsection}{A2.1.6 - CTR}
La Tabla~\ref{tab:t12-a2-1-6-ctr} individualiza los requisitos de este ámbito, su respaldo y la evidencia que deberá comprobarse.
\begin{longtable}{|L{0.09\AnchoMatrizTDoce}|L{0.385\AnchoMatrizTDoce}|L{0.065\AnchoMatrizTDoce}|L{0.18\AnchoMatrizTDoce}|L{0.28\AnchoMatrizTDoce}|}
\caption{Cumplimiento y trazabilidad: A2.1.6 - CTR}\label{tab:t12-a2-1-6-ctr}\\
\hline
\EncabezadoTabla{ID requerimiento} & \EncabezadoTabla{Descripción} & \EncabezadoTabla{Cumple\newline(Sí/No)} & \EncabezadoTabla{Componente que lo satisface} & \EncabezadoTabla{Sección de la propuesta} \\ \hline
\endfirsthead
\hline
\EncabezadoTabla{ID requerimiento} & \EncabezadoTabla{Descripción} & \EncabezadoTabla{Cumple\newline(Sí/No)} & \EncabezadoTabla{Componente que lo satisface} & \EncabezadoTabla{Sección de la propuesta} \\ \hline
\endhead
\multicolumn{5}{r}{\fontsize{9}{11}\selectfont Continúa en la página siguiente.}\\
\endfoot
\hline
\endlastfoot

RF-\allowbreak{}CTR-\allowbreak{}01 &
\textbf{[Must]} Mantener el padrón de empresas y trabajadores externos con identidad, seguros, acreditaciones y vigencias. &
Sí &
Varadero y contratistas; Integración externa &
SD2, sec. A2.1.6; SD4, sec. 4.1; SD5, sec. 5.1.2.4. \\ \hline

RF-\allowbreak{}CTR-\allowbreak{}02 &
\textbf{[Must]} Generar avisos preventivos por vencimiento de pólizas y acreditaciones. Se propone una anticipación inicial de 15 días, con registro de la fecha de vencimiento y del responsable de gestionar la renovación. &
Sí &
Varadero y contratistas; Integración externa &
SD2, sec. A2.1.6; SD4, sec. 4.1; SD5, sec. 5.1.2.4. \\ \hline

RF-\allowbreak{}CTR-\allowbreak{}03 &
\textbf{[Must]} Registrar la realización de la inducción de seguridad y ambiental de cada trabajador externo previa a su habilitación. &
Sí &
Varadero y contratistas; Integración externa  &
SD2, sec. A2.1.6; SD4, sec. 4.1; SD5, sec. 5.1.2.4; SD3, sec. 3.2, CA17. \\ \hline

RF-\allowbreak{}CTR-\allowbreak{}04 &
\textbf{[Must]} Vincular cada solicitud de acceso a una orden de trabajo formal autorizada por el armador de la nave intervenida. &
Sí &
Varadero y contratistas; Integración externa &
SD2, sec. A2.1.6; SD4, sec. 4.1; SD5, sec. 5.1.2.4. \\ \hline

RF-\allowbreak{}CTR-\allowbreak{}05 &
\textbf{[Must]} Verificar mediante credencial la identidad, inducción y faena autorizada antes del acceso; objetivo propuesto de validación: cinco segundos. &
Sí &
Varadero y contratistas; Integración externa &
SD2, sec. A2.1.6; SD4, sec. 4.1; SD5, sec. 5.1.2.4. \\ \hline

RF-\allowbreak{}CTR-\allowbreak{}06 &
\textbf{[Must]} Registrar en bitácora inalterable cualquier ingreso excepcional autorizado manualmente por la jefatura de puerto. &
Sí &
Varadero y contratistas; Integración externa &
SD2, sec. A2.1.6; SD4, sec. 4.1; SD5, sec. 5.1.2.4. \\ \hline

RF-\allowbreak{}CTR-\allowbreak{}07 &
\textbf{[Must]} Intercambiar estados de habilitación con el control de acceso existente mediante una interfaz disponible, documentada y probada. &
Sí &
Varadero y contratistas; Integración externa  &
SD2, sec. A2.1.6; SD4, sec. 4.1; SD5, sec. 5.1.2.4. \\ \hline
\end{longtable}


\subsection*{A2.1.7 - AMB}
\addcontentsline{toc}{subsection}{A2.1.7 - AMB}
La Tabla~\ref{tab:t12-a2-1-7-amb} individualiza los requisitos de este ámbito, su respaldo y la evidencia que deberá comprobarse.
\begin{longtable}{|L{0.09\AnchoMatrizTDoce}|L{0.385\AnchoMatrizTDoce}|L{0.065\AnchoMatrizTDoce}|L{0.18\AnchoMatrizTDoce}|L{0.28\AnchoMatrizTDoce}|}
\caption{Cumplimiento y trazabilidad: A2.1.7 - AMB}\label{tab:t12-a2-1-7-amb}\\
\hline
\EncabezadoTabla{ID requerimiento} & \EncabezadoTabla{Descripción} & \EncabezadoTabla{Cumple\newline(Sí/No)} & \EncabezadoTabla{Componente que lo satisface} & \EncabezadoTabla{Sección de la propuesta} \\ \hline
\endfirsthead
\hline
\EncabezadoTabla{ID requerimiento} & \EncabezadoTabla{Descripción} & \EncabezadoTabla{Cumple\newline(Sí/No)} & \EncabezadoTabla{Componente que lo satisface} & \EncabezadoTabla{Sección de la propuesta} \\ \hline
\endhead
\multicolumn{5}{r}{\fontsize{9}{11}\selectfont Continúa en la página siguiente.}\\
\endfoot
\hline
\endlastfoot

RF-\allowbreak{}AMB-\allowbreak{}01 &
\textbf{[Must]} Registrar el origen de cada residuo y vincularlo con la embarcación y el generador identificados. Cuando el origen no esté acreditado, admitir un registro provisional de recepción física, señalar la discrepancia y asignar su aclaración al responsable ambiental. &
Sí &
Ambiente y consumos; Repositorio de objetos &
SD2, sec. A2.1.7; SD4, sec. 4.1.1; SD5, sec. 5.1.2.6. \\ \hline

RF-\allowbreak{}AMB-\allowbreak{}02 &
\textbf{[Must]} Asignar etiqueta con código QR único a cada tambor o bulto acopiado en el patio de residuos peligrosos (RESPEL). &
Sí &
Ambiente y consumos; Repositorio de objetos &
SD2, sec. A2.1.7; SD4, sec. 4.1.1; SD5, sec. 5.1.2.6. \\ \hline

RF-\allowbreak{}AMB-\allowbreak{}03 &
\textbf{[Must]} Registrar cantidad, unidad, clasificación aplicable, fecha, responsable y procedencia de cada residuo ingresado al acopio. La falta de identificación de origen no impide conservar cantidad, unidad, clasificación, fecha y responsable de recepción; se trata conforme a RF-AMB-01 &
Sí &
Ambiente y consumos; Repositorio de objetos &
SD2, sec. A2.1.7; SD4, sec. 4.1.1; SD5, sec. 5.1.2.6. \\ \hline

RF-\allowbreak{}AMB-\allowbreak{}04 &
\textbf{[Must]} Preparar y conservar antecedentes de entrega y transporte de residuos mediante los formatos y canales admitidos por la autoridad sanitaria. &
Sí &
Ambiente y consumos; Repositorio de objetos &
SD2, sec. A2.1.7; SD4, sec. 4.1.1; SD5, sec. 5.1.2.6. \\ \hline

RF-\allowbreak{}AMB-\allowbreak{}05 &
\textbf{[Must]} Almacenar el certificado final de disposición o tratamiento emitido por el gestor ambiental autorizado con retención de 6 años. &
Sí &
Ambiente y consumos; Repositorio de objetos &
SD2, sec. A2.1.7; SD4, sec. 4.1.1; SD5, sec. 5.1.2.6. \\ \hline

RF-\allowbreak{}AMB-\allowbreak{}06 &
\textbf{[Must]} Consolidar diariamente los indicadores ambientales y permitir informes por período, incluido el balance mensual propuesto para la solución, manteniendo trazabilidad hasta la lectura, residuo o transacción de origen. &
Sí &
Ambiente y consumos; Repositorio de objetos &
SD2, sec. A2.1.7; SD4, sec. 4.1.1; SD5, sec. 5.1.2.6. \\ \hline

RF-\allowbreak{}AMB-\allowbreak{}07 &
\textbf{[Must]} Registrar incidentes de derrame de hidrocarburos, plan de contingencia activado y evidencia de contención y remediación. &
Sí &
Ambiente y consumos; Repositorio de objetos &
SD2, sec. A2.1.7; SD4, sec. 4.1.1; SD5, sec. 5.1.2.6. \\ \hline
\end{longtable}


\subsection*{A2.1.8 - CON}
\addcontentsline{toc}{subsection}{A2.1.8 - CON}
La Tabla~\ref{tab:t12-a2-1-8-con} individualiza los requisitos de este ámbito, su respaldo y la evidencia que deberá comprobarse.
\begin{longtable}{|L{0.09\AnchoMatrizTDoce}|L{0.385\AnchoMatrizTDoce}|L{0.065\AnchoMatrizTDoce}|L{0.18\AnchoMatrizTDoce}|L{0.28\AnchoMatrizTDoce}|}
\caption{Cumplimiento y trazabilidad: A2.1.8 - CON}\label{tab:t12-a2-1-8-con}\\
\hline
\EncabezadoTabla{ID requerimiento} & \EncabezadoTabla{Descripción} & \EncabezadoTabla{Cumple\newline(Sí/No)} & \EncabezadoTabla{Componente que lo satisface} & \EncabezadoTabla{Sección de la propuesta} \\ \hline
\endfirsthead
\hline
\EncabezadoTabla{ID requerimiento} & \EncabezadoTabla{Descripción} & \EncabezadoTabla{Cumple\newline(Sí/No)} & \EncabezadoTabla{Componente que lo satisface} & \EncabezadoTabla{Sección de la propuesta} \\ \hline
\endhead
\multicolumn{5}{r}{\fontsize{9}{11}\selectfont Continúa en la página siguiente.}\\
\endfoot
\hline
\endlastfoot

RF-\allowbreak{}CON-\allowbreak{}01 &
\textbf{[Must]} Asociar unívocamente cada medidor de agua y electricidad de torretas al puesto de amarra y contrato activo. &
Sí &
Ambiente y consumos; Combustible; Infraestructura &
SD2, sec. A2.1.8; SD4, sec. 4.2; SD5, sec. 5.1.2.5, 5.1.2.6 y 5.1.2.8. \\ \hline

RF-\allowbreak{}CON-\allowbreak{}02 &
\textbf{[Must]} Ingerir las lecturas de agua y energía con la frecuencia de adquisición definida y justificada para cada servicio, conservando su origen, marca temporal y estado de calidad. La información acumulada por amarra debe actualizarse al menos cada hora, conforme a RNF-DES-08. &
Sí &
Ambiente y consumos; Combustible; Infraestructura &
SD2, sec. A2.1.8; SD4, sec. 4.2; SD5, sec. 5.1.2.5, 5.1.2.6 y 5.1.2.8. \\ \hline

RF-\allowbreak{}CON-\allowbreak{}03 &
\textbf{[Must]} Detectar duplicados y lecturas anómalas, conservar el dato de origen y su estado de calidad, y separar los valores no validados del cálculo de cobro. Toda corrección quedará justificada y auditable. &
Sí &
Ambiente y consumos; Combustible; Infraestructura &
SD2, sec. A2.1.8; SD4, sec. 4.2; SD5, sec. 5.1.2.5, 5.1.2.6 y 5.1.2.8. \\ \hline

RF-\allowbreak{}CON-\allowbreak{}04 &
\textbf{[Must]} Calcular consumo y cargo por amarra mediante lecturas atribuibles, método de imputación documentado y tarifa aprobada. &
Sí &
Ambiente y consumos; Combustible; Infraestructura &
SD2, sec. A2.1.8; SD4, sec. 4.2; SD5, sec. 5.1.2.5, 5.1.2.6 y 5.1.2.8. \\ \hline

RF-\allowbreak{}CON-\allowbreak{}05 &
\textbf{[Must]} Emitir eventos de hecho facturable por consumo validado hacia la capa financiera sin generar facturas tributarias directas. &
Sí &
Ambiente y consumos; Combustible; Infraestructura &
SD2, sec. A2.1.8; SD4, sec. 4.2; SD5, sec. 5.1.2.5, 5.1.2.6 y 5.1.2.8. \\ \hline

RF-\allowbreak{}CON-\allowbreak{}06 &
\textbf{[Must]} Garantizar una única serie de datos de submedición compartida para el cobro financiero y la huella ambiental. &
Sí &
Ambiente y consumos; Combustible; Infraestructura &
SD2, sec. A2.1.8; SD4, sec. 4.2; SD5, sec. 5.1.2.5, 5.1.2.6 y 5.1.2.8. \\ \hline

RF-\allowbreak{}CON-\allowbreak{}07 &
\textbf{[Must]} Registrar embarcación, fecha, hora, tipo, volumen y evidencia de recepción de cada despacho; conciliarlo con la instrumentación del surtidor. &
Sí &
Ambiente y consumos; Combustible; Infraestructura &
SD2, sec. A2.1.8; SD4, sec. 4.2; SD5, sec. 5.1.2.5, 5.1.2.6 y 5.1.2.8. \\ \hline

RF-\allowbreak{}CON-\allowbreak{}08 &
\textbf{[Must]} Permitir la captura y autorización de despachos de combustible en modo offline en el muelle de combustible mediante el Runtime Operacional de Borde de la Célula Operacional Local. &
Sí &
Combustible; Autoridad local de acceso  &
SD4, sec. 4.1 y 4.2; SD5, sec. 5.1.2.9, 5.2.2 y 5.4.3. \\ \hline

RF-\allowbreak{}CON-\allowbreak{}09 &
\textbf{[Should]} Registrar los consumos de servicios y la relación contractual del restaurante concesionado, sin administrar su punto de venta ni su operación comercial. &
Sí &
Ambiente y consumos; Combustible; Infraestructura &
SD2, sec. A2.1.8; SD4, sec. 4.2; SD5, sec. 5.1.2.5, 5.1.2.6 y 5.1.2.8. \\ \hline
\end{longtable}


\subsection*{A2.1.9 - BOY}
\addcontentsline{toc}{subsection}{A2.1.9 - BOY}
La Tabla~\ref{tab:t12-a2-1-9-boy} individualiza los requisitos de este ámbito, su respaldo y la evidencia que deberá comprobarse.
\begin{longtable}{|L{0.09\AnchoMatrizTDoce}|L{0.385\AnchoMatrizTDoce}|L{0.065\AnchoMatrizTDoce}|L{0.18\AnchoMatrizTDoce}|L{0.28\AnchoMatrizTDoce}|}
\caption{Cumplimiento y trazabilidad: A2.1.9 - BOY}\label{tab:t12-a2-1-9-boy}\\
\hline
\EncabezadoTabla{ID requerimiento} & \EncabezadoTabla{Descripción} & \EncabezadoTabla{Cumple\newline(Sí/No)} & \EncabezadoTabla{Componente que lo satisface} & \EncabezadoTabla{Sección de la propuesta} \\ \hline
\endfirsthead
\hline
\EncabezadoTabla{ID requerimiento} & \EncabezadoTabla{Descripción} & \EncabezadoTabla{Cumple\newline(Sí/No)} & \EncabezadoTabla{Componente que lo satisface} & \EncabezadoTabla{Sección de la propuesta} \\ \hline
\endhead
\multicolumn{5}{r}{\fontsize{9}{11}\selectfont Continúa en la página siguiente.}\\
\endfoot
\hline
\endlastfoot

RF-\allowbreak{}BOY-\allowbreak{}01 &
\textbf{[Must]} Mantener inventario técnico de las 40 boyas de Bahía Ilque: grillete, cadena, cabo, muerto de hormigón y profundidad. &
Sí &
Infraestructura; Operación marítima; aplicación Android &
SD2, sec. A2.1.9; SD4, sec. 4.2.1; SD5, sec. 5.1.2.8 y 5.2.2. \\ \hline

RF-\allowbreak{}BOY-\allowbreak{}02 &
\textbf{[Must]} Registrar observaciones visuales de ocupación de boyas capturando patrón, nave amarrada y hora de verificación. &
Sí &
Infraestructura; Operación marítima; aplicación Android &
SD2, sec. A2.1.9; SD4, sec. 4.2.1; SD5, sec. 5.1.2.8 y 5.2.2. \\ \hline

RF-\allowbreak{}BOY-\allowbreak{}03 &
\textbf{[Must]} Conciliar observaciones de ocupación de boyas con las reservas y declaraciones disponibles, mostrando fuente y hora de cada antecedente. Identificar casos sin confirmación vigente y las discrepancias, y generar un reporte diario. &
Sí &
Infraestructura; Operación marítima; aplicación Android &
SD2, sec. A2.1.9; SD4, sec. 4.2.1; SD5, sec. 5.1.2.8 y 5.2.2. \\ \hline

RF-\allowbreak{}BOY-\allowbreak{}04 &
\textbf{[Could]} Integrar, de aprobarse su inclusión, la telemetría de un piloto de boyas con alimentación y comunicaciones acreditadas. La inspección desconectada seguirá siendo independiente. &
Sí &
Infraestructura; Operación marítima; aplicación Android  &
SD2, sec. A2.1.9; SD4, sec. 4.2.1; SD5, sec. 5.1.2.8 y 5.2.2. \\ \hline

RF-\allowbreak{}BOY-\allowbreak{}05 &
\textbf{[Must]} Capturar bitácoras de inspección submarina y fondeo en terreno desconectado desde la embarcación de trabajo. &
Sí &
Infraestructura; Operación marítima; aplicación Android &
SD2, sec. A2.1.9; SD4, sec. 4.2.1; SD5, sec. 5.1.2.8 y 5.2.2. \\ \hline

RF-\allowbreak{}BOY-\allowbreak{}06 &
\textbf{[Must]} Registrar por componente del tren de fondeo las inspecciones, desgaste, evidencia fotográfica y sustituciones efectivamente realizadas por el responsable de mantenimiento. El sistema no asume la ejecución del mantenimiento físico excluido. &
Sí &
Infraestructura; Operación marítima; aplicación Android &
SD2, sec. A2.1.9; SD4, sec. 4.2.1; SD5, sec. 5.1.2.8 y 5.2.2. \\ \hline

RF-\allowbreak{}BOY-\allowbreak{}07 &
\textbf{[Must]} Programar y alertar inspecciones periódicas y extraordinarias conforme al programa técnico aprobado del tren de fondeo. &
Sí &
Infraestructura; Operación marítima; aplicación Android  &
SD2, sec. A2.1.9; SD4, sec. 4.2.1; SD5, sec. 5.1.2.8 y 5.2.2. \\ \hline
\end{longtable}


\subsection*{A2.1.10 - FAC}
\addcontentsline{toc}{subsection}{A2.1.10 - FAC}
La Tabla~\ref{tab:t12-a2-1-10-fac} individualiza los requisitos de este ámbito, su respaldo y la evidencia que deberá comprobarse.
\begin{longtable}{|L{0.09\AnchoMatrizTDoce}|L{0.385\AnchoMatrizTDoce}|L{0.065\AnchoMatrizTDoce}|L{0.18\AnchoMatrizTDoce}|L{0.28\AnchoMatrizTDoce}|}
\caption{Cumplimiento y trazabilidad: A2.1.10 - FAC}\label{tab:t12-a2-1-10-fac}\\
\hline
\EncabezadoTabla{ID requerimiento} & \EncabezadoTabla{Descripción} & \EncabezadoTabla{Cumple\newline(Sí/No)} & \EncabezadoTabla{Componente que lo satisface} & \EncabezadoTabla{Sección de la propuesta} \\ \hline
\endfirsthead
\hline
\EncabezadoTabla{ID requerimiento} & \EncabezadoTabla{Descripción} & \EncabezadoTabla{Cumple\newline(Sí/No)} & \EncabezadoTabla{Componente que lo satisface} & \EncabezadoTabla{Sección de la propuesta} \\ \hline
\endhead
\multicolumn{5}{r}{\fontsize{9}{11}\selectfont Continúa en la página siguiente.}\\
\endfoot
\hline
\endlastfoot

RF-\allowbreak{}FAC-\allowbreak{}01 &
\textbf{[Must]} Consolidar en cuenta corriente del armador todos los hechos facturables originados en amarras, varadero, consumos y servicios. &
Sí &
Finanzas y cobranza; Integración externa &
SD2, sec. A2.1.10; SD4, sec. 4.1; SD5, sec. 5.1.2.7. \\ \hline

RF-\allowbreak{}FAC-\allowbreak{}02 &
\textbf{[Must]} Validar reglas de tarificación, vigencia contractual y deduplicación con clave de idempotencia antes de exportar a contabilidad. &
Sí &
Finanzas y cobranza; Integración externa &
SD2, sec. A2.1.10; SD4, sec. 4.1; SD5, sec. 5.1.2.7. \\ \hline

RF-\allowbreak{}FAC-\allowbreak{}03 &
\textbf{[Must]} Exportar hechos facturables validados al sistema contable y registrar recepción o error mediante una interfaz verificada. &
Sí &
Finanzas y cobranza; Integración externa  &
SD2, sec. A2.1.10; SD4, sec. 4.1; SD5, sec. 5.1.2.7. \\ \hline

RF-\allowbreak{}FAC-\allowbreak{}04 &
\textbf{[Must]} Procesar y asentar la respuesta de conciliación devuelta por el sistema contable (documentos emitidos, pagos y rechazos). &
Sí &
Finanzas y cobranza; Integración externa  &
SD2, sec. A2.1.10; SD4, sec. 4.1; SD5, sec. 5.1.2.7. \\ \hline

RF-\allowbreak{}FAC-\allowbreak{}05 &
\textbf{[Must]} Presentar al armador su estado de cuenta en tiempo real, con el origen y estado de conciliación de cada cargo o pago, sin sustituir la emisión tributaria del sistema contable. &
Sí &
Finanzas y cobranza; Integración externa  &
SD2, sec. A2.1.10; SD4, sec. 4.1; SD5, sec. 5.1.2.7. \\ \hline

RF-\allowbreak{}FAC-\allowbreak{}06 &
\textbf{[Must]} Clasificar deudas y registrar gestiones de cobranza según reglas aprobadas, sin activar restricciones automáticas de acceso. &
Sí &
Finanzas y cobranza; Integración externa &
SD2, sec. A2.1.10; SD4, sec. 4.1; SD5, sec. 5.1.2.7. \\ \hline

RF-\allowbreak{}FAC-\allowbreak{}07 &
\textbf{[Must]} Registrar decisión formal, fundamento y autorización competente antes de aplicar una medida de cobranza que afecte el acceso a la embarcación. &
Sí &
Finanzas y cobranza; Integración externa &
SD2, sec. A2.1.10; SD4, sec. 4.1; SD5, sec. 5.1.2.7. \\ \hline

RF-\allowbreak{}FAC-\allowbreak{}08 &
\textbf{[Must]} Conformar expediente probatorio digital inalterable para las 14 naves en abandono (contratos, mora, cartas y fotos de custodia). &
Sí &
Finanzas y cobranza; Integración externa &
SD2, sec. A2.1.10; SD4, sec. 4.1; SD5, sec. 5.1.2.7. \\ \hline

RF-\allowbreak{}FAC-\allowbreak{}09 &
\textbf{[Must]} Exportar el expediente de custodia con índice, documentos, fechas y control de integridad para el asesor autorizado, sin ejecutar actuaciones judiciales. &
Sí &
Finanzas y cobranza; Integración externa &
SD2, sec. A2.1.10; SD4, sec. 4.1; SD5, sec. 5.1.2.7. \\ \hline
\end{longtable}


\subsection*{A2.1.11 - ADM}
\addcontentsline{toc}{subsection}{A2.1.11 - ADM}
La Tabla~\ref{tab:t12-a2-1-11-adm} individualiza los requisitos de este ámbito, su respaldo y la evidencia que deberá comprobarse.
\begin{longtable}{|L{0.09\AnchoMatrizTDoce}|L{0.385\AnchoMatrizTDoce}|L{0.065\AnchoMatrizTDoce}|L{0.18\AnchoMatrizTDoce}|L{0.28\AnchoMatrizTDoce}|}
\caption{Cumplimiento y trazabilidad: A2.1.11 - ADM}\label{tab:t12-a2-1-11-adm}\\
\hline
\EncabezadoTabla{ID requerimiento} & \EncabezadoTabla{Descripción} & \EncabezadoTabla{Cumple\newline(Sí/No)} & \EncabezadoTabla{Componente que lo satisface} & \EncabezadoTabla{Sección de la propuesta} \\ \hline
\endfirsthead
\hline
\EncabezadoTabla{ID requerimiento} & \EncabezadoTabla{Descripción} & \EncabezadoTabla{Cumple\newline(Sí/No)} & \EncabezadoTabla{Componente que lo satisface} & \EncabezadoTabla{Sección de la propuesta} \\ \hline
\endhead
\multicolumn{5}{r}{\fontsize{9}{11}\selectfont Continúa en la página siguiente.}\\
\endfoot
\hline
\endlastfoot

RF-\allowbreak{}ADM-\allowbreak{}01 &
\textbf{[Must]} Administrar padrones maestros centrales de armadores, tripulantes, embarcaciones, contratos de amarra y tarifas. &
Sí &
Auditoría y seguridad; Núcleo Administrativo &
SD2, sec. A2.1.11; SD4, sec. 4.1; SD5, sec. 5.1.2.9 y 5.1.3. \\ \hline

RF-\allowbreak{}ADM-\allowbreak{}02 &
\textbf{[Must]} Modificar parámetros de operación (tiempos de cortesía, recargos de temporada) desde interfaz web sin recompilar código. &
Sí &
Auditoría y seguridad; Núcleo Administrativo &
SD2, sec. A2.1.11; SD4, sec. 4.1; SD5, sec. 5.1.2.9 y 5.1.3. \\ \hline

RF-\allowbreak{}ADM-\allowbreak{}03 &
\textbf{[Must]} Aplicar políticas de control de acceso RBAC/ABAC basadas en mínimo privilegio y separación de funciones operacionales. (conforme a Ley N° 21.719 sobre Protección de Datos Personales y Ley N° 19.628). &
Sí &
Auditoría y seguridad; Repositorio de objetos &
SD2, sec. A2.1.11; SD5, sec. 5.1.2.9 y 5.2.5. \\ \hline

RF-\allowbreak{}ADM-\allowbreak{}04 &
\textbf{[Must]} Gestionar contratos de amarra, renovaciones y sus firmas electrónicas en la modalidad admitida para el acto, con versión, vigencia y evidencia de aceptación verificable. (conforme a Ley N° 21.719 sobre Protección de Datos Personales y Ley N° 19.628). &
Sí &
Auditoría y seguridad; Núcleo Administrativo  &
SD2, sec. A2.1.11; SD4, sec. 4.1; SD5, sec. 5.1.2.9 y 5.1.3. \\ \hline

RF-\allowbreak{}ADM-\allowbreak{}05 &
\textbf{[Must]} Registrar en bitácora append-only todas las creaciones, modificaciones, accesos a datos sensibles y autorizaciones forzadas. &
Sí &
Auditoría y seguridad; Repositorio de objetos &
SD2, sec. A2.1.11; SD5, sec. 5.1.2.9 y 5.2.5. \\ \hline

RF-\allowbreak{}ADM-\allowbreak{}06 &
\textbf{[Must]} Ejecutar la migración del sistema de 2014 con perfilado, saneamiento trazable y conciliación, incluyendo íntegramente los conjuntos exigidos por C07 RT-05.15. Ningún descarte técnico podrá eliminar silenciosamente antecedentes obligatorios. &
Sí &
Migración; repositorio histórico &
SD2, sec. A2.1.11; SD5, sec. 5.3.1 a 5.3.3. \\ \hline

RF-\allowbreak{}ADM-\allowbreak{}07 &
\textbf{[Must]} Proveer consulta de los antecedentes históricos no migrados durante su retención aplicable, con integridad y acceso restringido. No exigir la continuidad del software sin soporte si un repositorio verificable satisface la consulta. &
Sí &
Migración; repositorio histórico &
SD2, sec. A2.1.11; SD5, sec. 5.3.1 a 5.3.3. \\ \hline

RF-\allowbreak{}ADM-\allowbreak{}08 &
\textbf{[Must]} Resolver discrepancias entre registros locales y centrales mediante reglas deterministas, conservando los antecedentes y la bitácora de conflictos. &
Sí &
Coordinador de sincronización &
SD2, sec. A2.1.11; SD5, sec. 5.2.2. \\ \hline

RF-\allowbreak{}ADM-\allowbreak{}09 &
\textbf{[Must]} Gestionar las firmas de recepción y entrega en custodia, actas de izaje y entrega de residuos, además de contratos y autorizaciones de menores, conservando identidad, documento, versión y evidencia de validez. (conforme a Ley N° 21.719 sobre Protección de Datos Personales y Ley N° 19.628). &
Sí &
Auditoría y seguridad; Repositorio de objetos  &
SD2, sec. A2.1.11; SD5, sec. 5.1.2.9 y 5.2.5. \\ \hline

RF-\allowbreak{}ADM-\allowbreak{}10 &
\textbf{[Must]} Integrar la estación meteorológica, avisos de la autoridad y sistemas existentes de acceso y 38 cámaras únicamente para los usos autorizados del recinto, con registro de indisponibilidad. No incorporar cámaras dirigidas a alumnos. &
Sí &
Auditoría y seguridad; Núcleo Administrativo  &
SD2, sec. A2.1.11; SD4, sec. 4.1; SD5, sec. 5.1.2.9 y 5.1.3. \\ \hline

RF-\allowbreak{}ADM-\allowbreak{}11 &
\textbf{[Must]} La solución dispondrá de un módulo de administración que permita al CLIENTE, sin intervención del ADJUDICATARIO, gestionar personas usuarias, roles, permisos, unidades organizacionales y sus jerarquías. &
Sí &
Auditoría y seguridad; Núcleo Administrativo  &
SD2, sec. A2.1.11; SD4, sec. 4.1; SD5, sec. 5.1.2.9 y 5.1.3. \\ \hline

RF-\allowbreak{}ADM-\allowbreak{}12 &
\textbf{[Must]} Las reglas de negocio parametrizables ; umbrales, plazos, montos, tolerancias, catálogos, listas de valores, textos de notificación ;  serán configurables desde la interfaz de administración, con control de versiones y registro de quién cambió qué y cuándo. &
Sí &
Auditoría y seguridad; Núcleo Administrativo &
SD2, sec. A2.1.11; SD4, sec. 4.1; SD5, sec. 5.1.2.9 y 5.1.3. \\ \hline

RF-\allowbreak{}ADM-\allowbreak{}13 &
\textbf{[Must]} Todo cambio de parámetro con impacto operacional requerirá aprobación de un segundo perfil y quedará registrado con su justificación. &
Sí &
Auditoría y seguridad; Núcleo Administrativo &
SD2, sec. A2.1.11; SD4, sec. 4.1; SD5, sec. 5.1.2.9 y 5.1.3. \\ \hline

RF-\allowbreak{}ADM-\allowbreak{}14 &
\textbf{[Must]} Publicar el inventario de elementos parametrizables y de cambios que requieren desarrollo, diferenciando ambos alcances. &
Sí &
Auditoría y seguridad; Núcleo Administrativo  &
SD2, sec. A2.1.11; SD4, sec. 4.1; SD5, sec. 5.1.2.9 y 5.1.3. \\ \hline

RF-\allowbreak{}ADM-\allowbreak{}15 &
\textbf{[Must]} Toda operación que cree, modifique o elimine información quedará registrada con identificación de la persona o del sistema que la ejecutó, fecha y hora con zona horaria, origen, valores anteriores y valores posteriores. &
Sí &
Auditoría y seguridad; Repositorio de objetos &
SD2, sec. A2.1.11; SD5, sec. 5.1.2.9 y 5.2.5. \\ \hline

RF-\allowbreak{}ADM-\allowbreak{}16 &
\textbf{[Must]} El registro de auditoría será inalterable y no podrá ser modificado ni eliminado por ningún perfil, incluido el administrador de la plataforma. &
Sí &
Auditoría y seguridad; Repositorio de objetos &
SD2, sec. A2.1.11; SD5, sec. 5.1.2.9 y 5.2.5. \\ \hline

RF-\allowbreak{}ADM-\allowbreak{}17 &
\textbf{[Must]} El CLIENTE podrá consultar y exportar la auditoría desde la interfaz, con filtros por persona, período, entidad y tipo de operación, sin requerir acceso a la base de datos. &
Sí &
Auditoría y seguridad; Repositorio de objetos  &
SD2, sec. A2.1.11; SD5, sec. 5.1.2.9 y 5.2.5. \\ \hline

RF-\allowbreak{}ADM-\allowbreak{}18 &
\textbf{[Must]} Las consultas a información sensible quedarán registradas, no sólo las modificaciones. &
Sí &
Auditoría y seguridad; Repositorio de objetos &
SD2, sec. A2.1.11; SD5, sec. 5.1.2.9 y 5.2.5. \\ \hline

RF-\allowbreak{}ADM-\allowbreak{}19 &
\textbf{[Must]} Conservar la auditoría según el plazo del caso y, en ausencia de uno específico, durante al menos cinco años. &
Sí &
Auditoría y seguridad; Repositorio de objetos &
SD2, sec. A2.1.11; SD5, sec. 5.1.2.9 y 5.2.5. \\ \hline

RF-\allowbreak{}ADM-\allowbreak{}20 &
\textbf{[Must]} La solución soportará flujos de trabajo con estados, transiciones, responsables, plazos, escalamiento automático por vencimiento y delegación por ausencia. &
Sí &
Auditoría y seguridad; Núcleo Administrativo  &
SD2, sec. A2.1.11; SD4, sec. 4.1; SD5, sec. 5.1.2.9 y 5.1.3. \\ \hline

RF-\allowbreak{}ADM-\allowbreak{}21 &
\textbf{[Must]} Los flujos serán configurables por el CLIENTE sin desarrollo, al menos en lo relativo a responsables, plazos y niveles de aprobación. &
Sí &
Auditoría y seguridad; Núcleo Administrativo  &
SD2, sec. A2.1.11; SD4, sec. 4.1; SD5, sec. 5.1.2.9 y 5.1.3. \\ \hline

RF-\allowbreak{}ADM-\allowbreak{}22 &
\textbf{[Must]} Toda solicitud pendiente será visible para su responsable en una bandeja de tareas unificada, con priorización y alerta de vencimiento. &
Sí &
Auditoría y seguridad; Núcleo Administrativo  &
SD2, sec. A2.1.11; SD4, sec. 4.1; SD5, sec. 5.1.2.9 y 5.1.3. \\ \hline

RF-\allowbreak{}ADM-\allowbreak{}23 &
\textbf{[Must]} La solución gestionará documentos con versionado, metadatos, control de acceso, búsqueda por contenido y por metadato, y previsualización sin descarga. &
Sí &
Auditoría y seguridad; Repositorio de objetos  &
SD2, sec. A2.1.11; SD5, sec. 5.1.2.9 y 5.2.5. \\ \hline

RF-\allowbreak{}ADM-\allowbreak{}24 &
\textbf{[Must]} Los documentos se almacenarán cifrados, con verificación de integridad y retención conforme a la política del caso. &
Sí &
Auditoría y seguridad; Repositorio de objetos &
SD2, sec. A2.1.11; SD5, sec. 5.1.2.9 y 5.2.5. \\ \hline

RF-\allowbreak{}ADM-\allowbreak{}25 &
\textbf{[Must]} La solución soportará firma electrónica conforme a la Ley N° 19.799, con firma avanzada para los actos que el caso lo requiera, y verificación de validez del certificado al momento de la firma. &
Sí &
Auditoría y seguridad; Repositorio de objetos  &
SD2, sec. A2.1.11; SD5, sec. 5.1.2.9 y 5.2.5. \\ \hline

RF-\allowbreak{}ADM-\allowbreak{}26 &
\textbf{[Must]} Se generará el sello de tiempo y se conservará la evidencia de firma que permita verificar el documento con posterioridad al vencimiento del certificado. &
Sí &
Auditoría y seguridad; Repositorio de objetos  &
SD2, sec. A2.1.11; SD5, sec. 5.1.2.9 y 5.2.5. \\ \hline

RF-\allowbreak{}ADM-\allowbreak{}27 &
\textbf{[Should]} La solución generará documentos a partir de plantillas administrables por el CLIENTE, con datos de la transacción y salida en formato abierto. &
Sí &
Auditoría y seguridad; Núcleo Administrativo  &
SD2, sec. A2.1.11; SD4, sec. 4.1; SD5, sec. 5.1.2.9 y 5.1.3. \\ \hline

RF-\allowbreak{}ADM-\allowbreak{}28 &
\textbf{[Should]} La solución dispondrá de búsqueda global con indexación de texto completo, tolerancia a errores de escritura, filtros facetados y respeto del control de acceso de la persona que busca. &
Sí &
Auditoría y seguridad; Núcleo Administrativo  &
SD2, sec. A2.1.11; SD4, sec. 4.1; SD5, sec. 5.1.2.9 y 5.1.3. \\ \hline

RF-\allowbreak{}ADM-\allowbreak{}29 &
\textbf{[Should]} Los listados serán ordenables, filtrables, paginados y exportables en formatos abiertos, con el filtro aplicado reflejado en la exportación. &
Sí &
Auditoría y seguridad; Núcleo Administrativo  &
SD2, sec. A2.1.11; SD4, sec. 4.1; SD5, sec. 5.1.2.9 y 5.1.3. \\ \hline

RF-\allowbreak{}ADM-\allowbreak{}30 &
\textbf{[Should]} Las exportaciones de gran volumen se procesarán de forma asíncrona, con notificación al completarse y sin bloquear la sesión. &
Sí &
Exportación íntegra y autónoma &
SD2, sec. A2.1.11; SD5, sec. 5.2.3. \\ \hline

RF-\allowbreak{}ADM-\allowbreak{}31 &
\textbf{[Must]} Toda exportación de información sensible quedará registrada en la auditoría, con identificación de quién exportó qué y cuándo. &
Sí &
Auditoría y seguridad; Repositorio de objetos &
SD2, sec. A2.1.11; SD5, sec. 5.1.2.9 y 5.2.5. \\ \hline

RF-\allowbreak{}ADM-\allowbreak{}32 &
\textbf{[Must]} La solución permitirá exportar la totalidad de la información del CLIENTE en formatos abiertos y documentados, en cualquier momento del Contrato, sin costo adicional y sin intervención del ADJUDICATARIO. &
Sí &
Exportación íntegra y autónoma &
SD2, sec. A2.1.11; SD5, sec. 5.2.3. \\ \hline

RF-\allowbreak{}ADM-\allowbreak{}33 &
\textbf{[Must]} La solución soportará la carga y descarga masiva de información en formatos abiertos, con validación previa, informe de errores por registro y procesamiento parcial. &
Sí &
Exportación íntegra y autónoma  &
SD2, sec. A2.1.11; SD5, sec. 5.2.3. \\ \hline

RF-\allowbreak{}ADM-\allowbreak{}34 &
\textbf{[Should]} La solución proveerá una capa analítica con tableros operacionales y de gestión, construidos sobre los indicadores que las Bases Técnicas del caso definan. &
Sí &
PostgreSQL separado de proyecciones; Parquet/Athena &
SD2, sec. A2.1.11; SD5, sec. 5.2.3 y 5.4.2. \\ \hline

RF-\allowbreak{}ADM-\allowbreak{}35 &
\textbf{[Should]} Los tableros permitirán filtrar por período, unidad organizacional y dimensiones propias del caso, y profundizar desde el indicador agregado hasta la transacción de origen. &
Sí &
PostgreSQL separado de proyecciones; Parquet/Athena &
SD2, sec. A2.1.11; SD5, sec. 5.2.3 y 5.4.2. \\ \hline

RF-\allowbreak{}ADM-\allowbreak{}36 &
\textbf{[Should]} El CLIENTE podrá construir sus propios informes sin intervención del ADJUDICATARIO, mediante una herramienta de autoservicio con modelo semántico documentado. &
Sí &
PostgreSQL separado de proyecciones; Parquet/Athena &
SD2, sec. A2.1.11; SD5, sec. 5.2.3 y 5.4.2. \\ \hline

RF-\allowbreak{}ADM-\allowbreak{}37 &
\textbf{[Should]} Todo informe será exportable en formatos abiertos y programable para envío automático por calendario. &
Sí &
PostgreSQL separado de proyecciones; Parquet/Athena &
SD2, sec. A2.1.11; SD5, sec. 5.2.3 y 5.4.2. \\ \hline
\end{longtable}


\subsection*{A2.1.12 - AUT}
\addcontentsline{toc}{subsection}{A2.1.12 - AUT}
La Tabla~\ref{tab:t12-a2-1-12-aut} individualiza los requisitos de este ámbito, su respaldo y la evidencia que deberá comprobarse.
\begin{longtable}{|L{0.09\AnchoMatrizTDoce}|L{0.385\AnchoMatrizTDoce}|L{0.065\AnchoMatrizTDoce}|L{0.18\AnchoMatrizTDoce}|L{0.28\AnchoMatrizTDoce}|}
\caption{Cumplimiento y trazabilidad: A2.1.12 - AUT}\label{tab:t12-a2-1-12-aut}\\
\hline
\EncabezadoTabla{ID requerimiento} & \EncabezadoTabla{Descripción} & \EncabezadoTabla{Cumple\newline(Sí/No)} & \EncabezadoTabla{Componente que lo satisface} & \EncabezadoTabla{Sección de la propuesta} \\ \hline
\endfirsthead
\hline
\EncabezadoTabla{ID requerimiento} & \EncabezadoTabla{Descripción} & \EncabezadoTabla{Cumple\newline(Sí/No)} & \EncabezadoTabla{Componente que lo satisface} & \EncabezadoTabla{Sección de la propuesta} \\ \hline
\endhead
\multicolumn{5}{r}{\fontsize{9}{11}\selectfont Continúa en la página siguiente.}\\
\endfoot
\hline
\endlastfoot

RF-\allowbreak{}AUT-\allowbreak{}01 &
\textbf{[Should]} Publicar sin autenticación la disponibilidad y tarifas de amarras de visita, condiciones de acceso, estado del recinto y avisos vigentes, en español e inglés; identificar actualización y fuente de los avisos. &
Sí &
Portal público y autenticado; aplicación Android &
SD2, sec. A2.1.12; SD4, sec. 4.1.1; SD5, sec. 5.1.2.1, 5.1.2.3 y 5.1.2.4. \\ \hline

RF-\allowbreak{}AUT-\allowbreak{}02 &
\textbf{[Should]} Proveer autoservicio para solicitud de recalada, reserva de amarra de visitas y pago electrónico con confirmación inmediata. &
Sí &
Portal público y autenticado; aplicación Android &
SD2, sec. A2.1.12; SD4, sec. 4.1.1; SD5, sec. 5.1.2.1, 5.1.2.3 y 5.1.2.4. \\ \hline

RF-\allowbreak{}AUT-\allowbreak{}03 &
\textbf{[Should]} Proveer al armador acceso autenticado a su contrato, consumos por amarra, cuenta corriente, documentación de la embarcación y plan de navegación, con separación respecto de otros titulares. &
Sí &
Portal público y autenticado; aplicación Android &
SD2, sec. A2.1.12; SD4, sec. 4.1.1; SD5, sec. 5.1.2.1, 5.1.2.3 y 5.1.2.4. \\ \hline

RF-\allowbreak{}AUT-\allowbreak{}04 &
\textbf{[Should]} Permitir la declaración y cierre voluntario del plan de navegación desde el móvil, con identificación y confirmación del acto, incluido el modo de baja capacidad. &
Sí &
Portal público y autenticado; aplicación Android  &
SD2, sec. A2.1.12; SD4, sec. 4.1.1; SD5, sec. 5.1.2.1, 5.1.2.3 y 5.1.2.4. \\ \hline

RF-\allowbreak{}AUT-\allowbreak{}05 &
\textbf{[Should]} Permitir al apoderado gestionar autorizaciones y consultar la confirmación de salida y regreso de la clase de su hijo, sin acceso a datos de otros alumnos ni seguimiento audiovisual. &
Sí &
Portal público y autenticado; aplicación Android &
SD2, sec. A2.1.12; SD4, sec. 4.1.1; SD5, sec. 5.1.2.1, 5.1.2.3 y 5.1.2.4. \\ \hline

RF-\allowbreak{}AUT-\allowbreak{}06 &
\textbf{[Should]} Permitir al contratista presentar su acreditación documental y declaración de residuos, además de solicitudes de trabajo, ofreciendo registro asistido para quien no utilice una aplicación. &
Sí &
Portal público y autenticado; aplicación Android &
SD2, sec. A2.1.12; SD4, sec. 4.1.1; SD5, sec. 5.1.2.1, 5.1.2.3 y 5.1.2.4. \\ \hline

RF-\allowbreak{}AUT-\allowbreak{}07 &
\textbf{[Should]} Proveer mecanismo seguro de autorregistro, verificación de correo/teléfono y recuperación de clave para usuarios externos. &
Sí &
Portal público y autenticado; aplicación Android  &
SD2, sec. A2.1.12; SD4, sec. 4.1.1; SD5, sec. 5.1.2.1, 5.1.2.3 y 5.1.2.4. \\ \hline

RF-\allowbreak{}AUT-\allowbreak{}08 &
\textbf{[Must]} Permitir al armador enviar declaraciones y cierres de plan mediante un modo de baja capacidad, con identificador, hora de origen y estado de confirmación; tolerar latencia de horas, reintentos y entrega desordenada sin reabrir un plan por un mensaje antiguo. &
Sí &
Portal público y autenticado; aplicación Android  &
SD2, sec. A2.1.12; SD4, sec. 4.1.1; SD5, sec. 5.1.2.1, 5.1.2.3 y 5.1.2.4. \\ \hline

RF-\allowbreak{}AUT-\allowbreak{}09 &
\textbf{[Should]} La solución dispondrá de un portal público con la información que el caso determine, accesible sin autenticación, con los mismos estándares de accesibilidad y desempeño que el resto de la plataforma. &
Sí &
Portal público y autenticado; aplicación Android &
SD2, sec. A2.1.12; SD4, sec. 4.1.1; SD5, sec. 5.1.2.1, 5.1.2.3 y 5.1.2.4. \\ \hline

RF-\allowbreak{}AUT-\allowbreak{}10 &
\textbf{[Should]} Las personas usuarias externas dispondrán de autoatención para las consultas de mayor frecuencia, evitando el contacto telefónico para operaciones simples. &
Sí &
Portal público y autenticado; aplicación Android &
SD2, sec. A2.1.12; SD4, sec. 4.1.1; SD5, sec. 5.1.2.1, 5.1.2.3 y 5.1.2.4. \\ \hline

RF-\allowbreak{}AUT-\allowbreak{}11 &
\textbf{[Must]} El portal público resistirá picos de tráfico sin degradar los servicios transaccionales internos, mediante aislamiento de recursos y caché. &
Sí &
Portal público y autenticado; aplicación Android &
SD2, sec. A2.1.12; SD4, sec. 4.1.1; SD5, sec. 5.1.2.1, 5.1.2.3 y 5.1.2.4. \\ \hline
\end{longtable}


\clearpage
\section*{Requerimientos no funcionales}
\addcontentsline{toc}{section}{Requerimientos no funcionales}
La condición conserva el umbral y método de A2.2, distinguiendo el diseño formulado de su verificación futura.
\subsection*{A2.2.1 - DES}
\addcontentsline{toc}{subsection}{A2.2.1 - DES}
La Tabla~\ref{tab:t12-a2-2-1-des} individualiza los requisitos de este ámbito, su respaldo y la evidencia que deberá comprobarse.
\begin{longtable}{|L{0.09\AnchoMatrizTDoce}|L{0.385\AnchoMatrizTDoce}|L{0.065\AnchoMatrizTDoce}|L{0.18\AnchoMatrizTDoce}|L{0.28\AnchoMatrizTDoce}|}
\caption{Cumplimiento y trazabilidad: A2.2.1 - DES}\label{tab:t12-a2-2-1-des}\\
\hline
\EncabezadoTabla{ID requerimiento} & \EncabezadoTabla{Descripción} & \EncabezadoTabla{Cumple\newline(Sí/No)} & \EncabezadoTabla{Componente que lo satisface} & \EncabezadoTabla{Sección de la propuesta} \\ \hline
\endfirsthead
\hline
\EncabezadoTabla{ID requerimiento} & \EncabezadoTabla{Descripción} & \EncabezadoTabla{Cumple\newline(Sí/No)} & \EncabezadoTabla{Componente que lo satisface} & \EncabezadoTabla{Sección de la propuesta} \\ \hline
\endhead
\multicolumn{5}{r}{\fontsize{9}{11}\selectfont Continúa en la página siguiente.}\\
\endfoot
\hline
\endlastfoot

RNF-\allowbreak{}DES-\allowbreak{}01 &
\textbf{[Must]} Asignación y confirmación de amarra de visita en \(\leq\)45 segundos desde la consulta por radio. &
Sí &
Núcleos modulares; Notificaciones; aplicación Android &
SD2, sec. A2.2.1; SD5, sec. 5.4.3. \\ \hline

RNF-\allowbreak{}DES-\allowbreak{}02 &
\textbf{[Must]} Registro de salida o regreso en \(\leq\)10 segundos y un máximo de dos interacciones. &
Sí &
Núcleos modulares; Notificaciones; aplicación Android &
SD2, sec. A2.2.1; SD5, sec. 5.4.3. \\ \hline

RNF-\allowbreak{}DES-\allowbreak{}03 &
\textbf{[Must]} Conformación de nómina de personas a bordo y legajo de zarpe en \(\leq\)90 segundos. &
Sí &
Núcleos modulares; Notificaciones; aplicación Android &
SD2, sec. A2.2.1; SD5, sec. 5.4.3. \\ \hline

RNF-\allowbreak{}DES-\allowbreak{}04 &
\textbf{[Must]} Registro de salida de una clase de hasta 18 embarcaciones en \(\leq\)60 segundos. &
Sí &
Núcleos modulares; Notificaciones; aplicación Android &
SD2, sec. A2.2.1; SD5, sec. 5.4.3. \\ \hline

RNF-\allowbreak{}DES-\allowbreak{}05 &
\textbf{[Must]} Notificación de emergencia a los 296 armadores en \(\leq\)10 minutos hasta el último envío, por los tres canales simultáneos exigidos. &
Sí &
Núcleos modulares; Notificaciones; aplicación Android &
SD2, sec. A2.2.1; SD5, sec. 5.4.3. \\ \hline

RNF-\allowbreak{}DES-\allowbreak{}06 &
\textbf{[Must]} Alerta por plan vencido y sin cierre en \(\leq\)15 minutos desde el vencimiento, sin sumar 30 minutos de gracia. &
Sí &
Núcleos modulares; Notificaciones; aplicación Android &
SD2, sec. A2.2.1; SD5, sec. 5.4.3. \\ \hline

RNF-\allowbreak{}DES-\allowbreak{}07 &
\textbf{[Must]} Estado de ocupación de amarra disponible en la capa analítica en \(\leq\)60 segundos desde la ocurrencia del evento, no desde su digitación. &
Sí &
Núcleos modulares; Notificaciones; aplicación Android &
SD2, sec. A2.2.1; SD5, sec. 5.4.3. \\ \hline

RNF-\allowbreak{}DES-\allowbreak{}08 &
\textbf{[Must]} Consumo acumulado por amarra actualizado al menos cada hora; no basta el total por torreta sin atribución al puesto. &
Sí &
Núcleos modulares; Notificaciones; aplicación Android &
SD2, sec. A2.2.1; SD5, sec. 5.4.3. \\ \hline

RNF-\allowbreak{}DES-\allowbreak{}09 &
\textbf{[Must]} Cuenta del armador en tiempo real, con identificación del origen y del estado de conciliación de cargos y pagos. &
Sí &
Núcleos modulares; Notificaciones; aplicación Android  &
SD2, sec. A2.2.1; SD5, sec. 5.4.3. \\ \hline

RNF-\allowbreak{}DES-\allowbreak{}10 &
\textbf{[Must]} Consolidación diaria de indicadores de residuos y huella ambiental ejecutada de forma automática al cierre del día. &
Sí &
Núcleos modulares; Notificaciones; aplicación Android &
SD2, sec. A2.2.1; SD5, sec. 5.4.3. \\ \hline

RNF-\allowbreak{}DES-\allowbreak{}11 &
\textbf{[Must]} Situación de escuela ; quién está en el agua;  en tiempo real, sin excepción. Ante una desconexión, no presentar como vigente una copia que no ha conciliado portería, escuela y monitor. &
Sí &
Núcleos modulares; Notificaciones; aplicación Android  &
SD2, sec. A2.2.1; SD5, sec. 5.4.3. \\ \hline

RNF-\allowbreak{}DES-\allowbreak{}12 &
\textbf{[Must]} Umbrales P95 de BT sec. 9.1: portal 2 s; navegación entre vistas 1 s; consulta simple API 500 ms; escritura API 800 ms; búsqueda compuesta 3 s; informe estándar 30 s; carga de 100 MB 60 s; arranque en frío 60 s. Capacidad por lotes de al menos 10.000 registros/minuto. Los procesos críticos usan los umbrales propios del caso. &
Sí &
Núcleos modulares; Notificaciones; aplicación Android &
SD2, sec. A2.2.1; SD5, sec. 5.4.3. \\ \hline
\end{longtable}


\subsection*{A2.2.2 - CONT}
\addcontentsline{toc}{subsection}{A2.2.2 - CONT}
La Tabla~\ref{tab:t12-a2-2-2-cont} individualiza los requisitos de este ámbito, su respaldo y la evidencia que deberá comprobarse.
\begin{longtable}{|L{0.09\AnchoMatrizTDoce}|L{0.385\AnchoMatrizTDoce}|L{0.065\AnchoMatrizTDoce}|L{0.18\AnchoMatrizTDoce}|L{0.28\AnchoMatrizTDoce}|}
\caption{Cumplimiento y trazabilidad: A2.2.2 - CONT}\label{tab:t12-a2-2-2-cont}\\
\hline
\EncabezadoTabla{ID requerimiento} & \EncabezadoTabla{Descripción} & \EncabezadoTabla{Cumple\newline(Sí/No)} & \EncabezadoTabla{Componente que lo satisface} & \EncabezadoTabla{Sección de la propuesta} \\ \hline
\endfirsthead
\hline
\EncabezadoTabla{ID requerimiento} & \EncabezadoTabla{Descripción} & \EncabezadoTabla{Cumple\newline(Sí/No)} & \EncabezadoTabla{Componente que lo satisface} & \EncabezadoTabla{Sección de la propuesta} \\ \hline
\endhead
\multicolumn{5}{r}{\fontsize{9}{11}\selectfont Continúa en la página siguiente.}\\
\endfoot
\hline
\endlastfoot

RNF-\allowbreak{}CONT-\allowbreak{}01 &
\textbf{[Must]} Al menos 24 horas continuas sin enlace exterior para recepción, asignación de amarra, salida y regreso, preparación de zarpes, combustible y escuela de vela. &
Sí &
Autoridad local de acceso; Runtime Operacional de Borde  &
SD4, sec. 4.1 y 4.2; SD5, sec. 5.1.2.9, 5.2.2 y 5.4.3. \\ \hline

RNF-\allowbreak{}CONT-\allowbreak{}02 &
\textbf{[Must]} Al menos 12 horas continuas fuera de cobertura para los dispositivos de pantalán, varadero y escuela, sin pérdida de registros. Verificar autonomía funcional y energética; la batería por sí sola no demuestra operación desconectada. &
Sí &
Autoridad local de acceso; Runtime Operacional de Borde  &
SD4, sec. 4.1 y 4.2.1; SD5, sec. 5.2.2 y 5.4.3. \\ \hline

RNF-\allowbreak{}CONT-\allowbreak{}03 &
\textbf{[Must]} Al menos ocho horas de operación de inspección de boyas sin cobertura, con almacenamiento y consulta de los antecedentes necesarios y sin pérdida de registros. &
Sí &
Autoridad local de acceso; Runtime Operacional de Borde  &
SD4, sec. 4.1 y 4.2.1; SD5, sec. 5.2.2 y 5.4.3. \\ \hline

RNF-\allowbreak{}CONT-\allowbreak{}04 &
\textbf{[Must]} Sincronización automática completa en \(\leq\)30 minutos después de 24 horas de desconexión, contados desde el restablecimiento del enlace, sin intervención manual. &
Sí &
Coordinador de sincronización; enlaces ≥100 Mb/s útiles &
SD4, sec. 4.1 y 4.2; SD5, sec. 5.2.2 y 5.4.3; SD2, sec. A2.7, EST-15. \\ \hline

RNF-\allowbreak{}CONT-\allowbreak{}05 &
\textbf{[Must]} Cero pérdida de registros de salida, regreso, combustible y asistencia de escuela durante desconexión y reconexión; reintentos sin duplicación de efectos de negocio. &
Sí &
Runtime Operacional de Borde; Coordinador de sincronización &
SD2, sec. A2.2.2; SD4, sec. 4.1 y 4.2; SD5, sec. 5.2.2. \\ \hline

RNF-\allowbreak{}CONT-\allowbreak{}06 &
\textbf{[Must]} Conciliación determinista y auditable, preservando antecedentes y aplicando reglas de resolución por entidad. &
Sí &
Runtime Operacional de Borde; Coordinador de sincronización &
SD2, sec. A2.2.2; SD4, sec. 4.1 y 4.2; SD5, sec. 5.2.2. \\ \hline
\end{longtable}


\subsection*{A2.2.3 - SEG}
\addcontentsline{toc}{subsection}{A2.2.3 - SEG}
La Tabla~\ref{tab:t12-a2-2-3-seg} individualiza los requisitos de este ámbito, su respaldo y la evidencia que deberá comprobarse.
\begin{longtable}{|L{0.09\AnchoMatrizTDoce}|L{0.385\AnchoMatrizTDoce}|L{0.065\AnchoMatrizTDoce}|L{0.18\AnchoMatrizTDoce}|L{0.28\AnchoMatrizTDoce}|}
\caption{Cumplimiento y trazabilidad: A2.2.3 - SEG}\label{tab:t12-a2-2-3-seg}\\
\hline
\EncabezadoTabla{ID requerimiento} & \EncabezadoTabla{Descripción} & \EncabezadoTabla{Cumple\newline(Sí/No)} & \EncabezadoTabla{Componente que lo satisface} & \EncabezadoTabla{Sección de la propuesta} \\ \hline
\endfirsthead
\hline
\EncabezadoTabla{ID requerimiento} & \EncabezadoTabla{Descripción} & \EncabezadoTabla{Cumple\newline(Sí/No)} & \EncabezadoTabla{Componente que lo satisface} & \EncabezadoTabla{Sección de la propuesta} \\ \hline
\endhead
\multicolumn{5}{r}{\fontsize{9}{11}\selectfont Continúa en la página siguiente.}\\
\endfoot
\hline
\endlastfoot

RNF-\allowbreak{}SEG-\allowbreak{}01 &
\textbf{[Must]} Gestión centralizada de identidad y SSO para la totalidad de módulos mediante OpenID Connect y OAuth 2.1. &
Sí &
Seguridad transversal; Auditoría y seguridad  &
SD2, sec. A2.2.3; SD4, sec. 4.1; SD5, sec. 5.2.5. \\ \hline

RNF-\allowbreak{}SEG-\allowbreak{}02 &
\textbf{[Must]} MFA para administradores, todo acceso privilegiado y accesos desde fuera de la red del cliente, con un mecanismo de terreno que preserve identidad individual. &
Sí &
Seguridad transversal; Auditoría y seguridad &
SD2, sec. A2.2.3; SD4, sec. 4.1; SD5, sec. 5.2.5. \\ \hline

RNF-\allowbreak{}SEG-\allowbreak{}03 &
\textbf{[Must]} Control de acceso basado en roles (RBAC) y atributos (ABAC) con matriz de segregación de funciones auditable. &
Sí &
Seguridad transversal; Auditoría y seguridad  &
SD2, sec. A2.2.3; SD4, sec. 4.1; SD5, sec. 5.2.5. \\ \hline

RNF-\allowbreak{}SEG-\allowbreak{}04 &
\textbf{[Must]} Cifrado a nivel de campo de salud, datos personales de menores, armadores, tripulantes y trabajadores externos, y posición compartida con consentimiento. con observancia estricta de la Ley N° 21.719 sobre Protección de Datos Personales y Ley N° 19.628. &
Sí &
Seguridad transversal; Auditoría y seguridad &
SD2, sec. A2.2.3; SD4, sec. 4.1; SD5, sec. 5.2.5. \\ \hline

RNF-\allowbreak{}SEG-\allowbreak{}05 &
\textbf{[Must]} Auditoría inalterable de las consultas a salud, datos personales de menores, posición consentida y antecedentes de deuda, además de las modificaciones, con identidad individual y origen. con observancia estricta de la Ley N° 21.719 sobre Protección de Datos Personales y Ley N° 19.628. &
Sí &
Seguridad transversal; Auditoría y seguridad &
SD2, sec. A2.2.3; SD4, sec. 4.1; SD5, sec. 5.2.5. \\ \hline

RNF-\allowbreak{}SEG-\allowbreak{}06 &
\textbf{[Must]} Registros técnicos sin claves, tokens ni datos personales expuestos; aplicar supresión o enmascaramiento antes de su almacenamiento. &
Sí &
Seguridad transversal; Auditoría y seguridad &
SD2, sec. A2.2.3; SD4, sec. 4.1; SD5, sec. 5.2.5. \\ \hline

RNF-\allowbreak{}SEG-\allowbreak{}07 &
\textbf{[Must]} Registro verificable del consentimiento y su alcance y revocación cuando corresponda, con controles reforzados para menores. La revocación no se interpreta como eliminación automática de antecedentes sujetos a retención exigible. con observancia estricta de la Ley N° 21.719 sobre Protección de Datos Personales y Ley N° 19.628. &
Sí &
Seguridad transversal; Auditoría y seguridad &
SD2, sec. A2.2.3; SD4, sec. 4.1; SD5, sec. 5.2.5. \\ \hline
\end{longtable}


\subsection*{A2.2.4 - USA}
\addcontentsline{toc}{subsection}{A2.2.4 - USA}
La Tabla~\ref{tab:t12-a2-2-4-usa} individualiza los requisitos de este ámbito, su respaldo y la evidencia que deberá comprobarse.
\begin{longtable}{|L{0.09\AnchoMatrizTDoce}|L{0.385\AnchoMatrizTDoce}|L{0.065\AnchoMatrizTDoce}|L{0.18\AnchoMatrizTDoce}|L{0.28\AnchoMatrizTDoce}|}
\caption{Cumplimiento y trazabilidad: A2.2.4 - USA}\label{tab:t12-a2-2-4-usa}\\
\hline
\EncabezadoTabla{ID requerimiento} & \EncabezadoTabla{Descripción} & \EncabezadoTabla{Cumple\newline(Sí/No)} & \EncabezadoTabla{Componente que lo satisface} & \EncabezadoTabla{Sección de la propuesta} \\ \hline
\endfirsthead
\hline
\EncabezadoTabla{ID requerimiento} & \EncabezadoTabla{Descripción} & \EncabezadoTabla{Cumple\newline(Sí/No)} & \EncabezadoTabla{Componente que lo satisface} & \EncabezadoTabla{Sección de la propuesta} \\ \hline
\endhead
\multicolumn{5}{r}{\fontsize{9}{11}\selectfont Continúa en la página siguiente.}\\
\endfoot
\hline
\endlastfoot

RNF-\allowbreak{}USA-\allowbreak{}01 &
\textbf{[Must]} Interfaces de muelle operables con guantes de agua, pantalla húmeda por lluvia y alto contraste bajo luz solar directa. &
Sí &
Portal público y autenticado; aplicación Android &
SD2, sec. A2.2.4; SD4, sec. 4.1.1 y 4.2.1. \\ \hline

RNF-\allowbreak{}USA-\allowbreak{}02 &
\textbf{[Must]} Cero interacción obligatoria con pantallas durante la ejecución activa de maniobras de travelift o amarre en muelle. &
Sí &
Portal público y autenticado; aplicación Android &
SD2, sec. A2.2.4; SD4, sec. 4.1.1 y 4.2.1. \\ \hline

RNF-\allowbreak{}USA-\allowbreak{}03 &
\textbf{[Should]} Indicador visible de conectividad y sincronización que distinga datos sincronizados, operación local y registros pendientes. &
Sí &
Portal público y autenticado; aplicación Android &
SD2, sec. A2.2.4; SD4, sec. 4.1.1 y 4.2.1. \\ \hline

RNF-\allowbreak{}USA-\allowbreak{}04 &
\textbf{[Must]} Todas las interfaces destinadas a usuarios cumplirán WCAG 2.2 nivel AA, con evaluación automática y manual; el alcance no se limita a los portales públicos. &
Sí &
Portal público y autenticado; aplicación Android &
SD2, sec. A2.2.4; SD4, sec. 4.1.1 y 4.2.1. \\ \hline
\end{longtable}


\subsection*{A2.2.5 - IDM}
\addcontentsline{toc}{subsection}{A2.2.5 - IDM}
La Tabla~\ref{tab:t12-a2-2-5-idm} individualiza los requisitos de este ámbito, su respaldo y la evidencia que deberá comprobarse.
\begin{longtable}{|L{0.09\AnchoMatrizTDoce}|L{0.385\AnchoMatrizTDoce}|L{0.065\AnchoMatrizTDoce}|L{0.18\AnchoMatrizTDoce}|L{0.28\AnchoMatrizTDoce}|}
\caption{Cumplimiento y trazabilidad: A2.2.5 - IDM}\label{tab:t12-a2-2-5-idm}\\
\hline
\EncabezadoTabla{ID requerimiento} & \EncabezadoTabla{Descripción} & \EncabezadoTabla{Cumple\newline(Sí/No)} & \EncabezadoTabla{Componente que lo satisface} & \EncabezadoTabla{Sección de la propuesta} \\ \hline
\endfirsthead
\hline
\EncabezadoTabla{ID requerimiento} & \EncabezadoTabla{Descripción} & \EncabezadoTabla{Cumple\newline(Sí/No)} & \EncabezadoTabla{Componente que lo satisface} & \EncabezadoTabla{Sección de la propuesta} \\ \hline
\endhead
\multicolumn{5}{r}{\fontsize{9}{11}\selectfont Continúa en la página siguiente.}\\
\endfoot
\hline
\endlastfoot

RNF-\allowbreak{}IDM-\allowbreak{}01 &
\textbf{[Must]} Portales web, avisos y comunicaciones dirigidos a visitantes, charters y extranjeros provistos en español e inglés. &
Sí &
Portal público y autenticado; aplicación Android &
SD2, sec. A2.2.5; SD4, sec. 4.1.1 y 4.2.1. \\ \hline

RNF-\allowbreak{}IDM-\allowbreak{}02 &
\textbf{[Must]} Atención en español e inglés en los horarios de C07 RT-21.06; incluye la cobertura de inglés en temporada exigida por RT-13.12. &
Sí &
Portal público y autenticado; aplicación Android  &
SD2, sec. A2.2.5; SD4, sec. 4.1.1 y 4.2.1. \\ \hline
\end{longtable}


\subsection*{A2.2.6 - ESCAL}
\addcontentsline{toc}{subsection}{A2.2.6 - ESCAL}
La Tabla~\ref{tab:t12-a2-2-6-escal} individualiza los requisitos de este ámbito, su respaldo y la evidencia que deberá comprobarse.
\begin{longtable}{|L{0.09\AnchoMatrizTDoce}|L{0.385\AnchoMatrizTDoce}|L{0.065\AnchoMatrizTDoce}|L{0.18\AnchoMatrizTDoce}|L{0.28\AnchoMatrizTDoce}|}
\caption{Cumplimiento y trazabilidad: A2.2.6 - ESCAL}\label{tab:t12-a2-2-6-escal}\\
\hline
\EncabezadoTabla{ID requerimiento} & \EncabezadoTabla{Descripción} & \EncabezadoTabla{Cumple\newline(Sí/No)} & \EncabezadoTabla{Componente que lo satisface} & \EncabezadoTabla{Sección de la propuesta} \\ \hline
\endfirsthead
\hline
\EncabezadoTabla{ID requerimiento} & \EncabezadoTabla{Descripción} & \EncabezadoTabla{Cumple\newline(Sí/No)} & \EncabezadoTabla{Componente que lo satisface} & \EncabezadoTabla{Sección de la propuesta} \\ \hline
\endhead
\multicolumn{5}{r}{\fontsize{9}{11}\selectfont Continúa en la página siguiente.}\\
\endfoot
\hline
\endlastfoot

RNF-\allowbreak{}ESCAL-\allowbreak{}01 &
\textbf{[Must]} Crecimiento dimensional soportado hasta 380 amarras y 70 boyas mediante parametrización funcional sin rediseño. &
Sí &
Infraestructura; ECS/Fargate; Telemetría &
SD2, sec. A2.2.6; SD4, sec. 4.2; SD5, sec. 5.4.1 y 5.4.3. \\ \hline

RNF-\allowbreak{}ESCAL-\allowbreak{}02 &
\textbf{[Should]} Capacidad de ingesta de submedición escalable hasta 760 puntos de agua y electricidad sin degradar latencia. &
Sí &
Infraestructura; ECS/Fargate; Telemetría &
SD2, sec. A2.2.6; SD4, sec. 4.2; SD5, sec. 5.4.1 y 5.4.3. \\ \hline
\end{longtable}


\subsection*{A2.2.7 - INT}
\addcontentsline{toc}{subsection}{A2.2.7 - INT}
La Tabla~\ref{tab:t12-a2-2-7-int} individualiza los requisitos de este ámbito, su respaldo y la evidencia que deberá comprobarse.
\begin{longtable}{|L{0.09\AnchoMatrizTDoce}|L{0.385\AnchoMatrizTDoce}|L{0.065\AnchoMatrizTDoce}|L{0.18\AnchoMatrizTDoce}|L{0.28\AnchoMatrizTDoce}|}
\caption{Cumplimiento y trazabilidad: A2.2.7 - INT}\label{tab:t12-a2-2-7-int}\\
\hline
\EncabezadoTabla{ID requerimiento} & \EncabezadoTabla{Descripción} & \EncabezadoTabla{Cumple\newline(Sí/No)} & \EncabezadoTabla{Componente que lo satisface} & \EncabezadoTabla{Sección de la propuesta} \\ \hline
\endfirsthead
\hline
\EncabezadoTabla{ID requerimiento} & \EncabezadoTabla{Descripción} & \EncabezadoTabla{Cumple\newline(Sí/No)} & \EncabezadoTabla{Componente que lo satisface} & \EncabezadoTabla{Sección de la propuesta} \\ \hline
\endhead
\multicolumn{5}{r}{\fontsize{9}{11}\selectfont Continúa en la página siguiente.}\\
\endfoot
\hline
\endlastfoot

RNF-\allowbreak{}INT-\allowbreak{}01 &
\textbf{[Must]} Integraciones autenticadas con OAuth 2.1 con credenciales de cliente o mTLS según contrato, validación de esquemas, correlación y reintentos idempotentes. No prometer entrega única de transporte como sustituto de la deduplicación de efectos. &
Sí &
Integración externa; núcleos modulares &
SD2, sec. A2.2.7; SD4, sec. 4.1. \\ \hline

RNF-\allowbreak{}INT-\allowbreak{}02 &
\textbf{[Must]} Diseño para degradación elegante: ante caída de pasarelas de pago o servicios externos, el núcleo náutico sigue activo. &
Sí &
Integración externa; núcleos modulares &
SD2, sec. A2.2.7; SD4, sec. 4.1. \\ \hline
\end{longtable}


\subsection*{A2.2.8 - OPE}
\addcontentsline{toc}{subsection}{A2.2.8 - OPE}
La Tabla~\ref{tab:t12-a2-2-8-ope} individualiza los requisitos de este ámbito, su respaldo y la evidencia que deberá comprobarse.
\begin{longtable}{|L{0.09\AnchoMatrizTDoce}|L{0.385\AnchoMatrizTDoce}|L{0.065\AnchoMatrizTDoce}|L{0.18\AnchoMatrizTDoce}|L{0.28\AnchoMatrizTDoce}|}
\caption{Cumplimiento y trazabilidad: A2.2.8 - OPE}\label{tab:t12-a2-2-8-ope}\\
\hline
\EncabezadoTabla{ID requerimiento} & \EncabezadoTabla{Descripción} & \EncabezadoTabla{Cumple\newline(Sí/No)} & \EncabezadoTabla{Componente que lo satisface} & \EncabezadoTabla{Sección de la propuesta} \\ \hline
\endfirsthead
\hline
\EncabezadoTabla{ID requerimiento} & \EncabezadoTabla{Descripción} & \EncabezadoTabla{Cumple\newline(Sí/No)} & \EncabezadoTabla{Componente que lo satisface} & \EncabezadoTabla{Sección de la propuesta} \\ \hline
\endhead
\multicolumn{5}{r}{\fontsize{9}{11}\selectfont Continúa en la página siguiente.}\\
\endfoot
\hline
\endlastfoot

RNF-\allowbreak{}OPE-\allowbreak{}01 &
\textbf{[Must]} Operación continua sin requerir personal de ingeniería informática en la nómina propia del club deportivo. &
Sí &
Operaciones y Servicios Gestionados  &
SD2, sec. A2.2.8; SD1, sec. 1.1.3 y 1.5; SD4, sec. 4.1. \\ \hline

RNF-\allowbreak{}OPE-\allowbreak{}02 &
\textbf{[Must]} Modificación ordinaria de calendarios, tarifas y listas de verificación realizada por personal de administración vía web. &
Sí &
Operaciones y Servicios Gestionados &
SD2, sec. A2.2.8; SD1, sec. 1.1.3 y 1.5; SD4, sec. 4.1. \\ \hline

RNF-\allowbreak{}OPE-\allowbreak{}03 &
\textbf{[Must]} Recuperación automática del componente local tras retorno de energía, con integridad de registros y supervisión de su estado. &
Sí &
Operaciones y Servicios Gestionados &
SD2, sec. A2.2.8; SD1, sec. 1.1.3 y 1.5; SD4, sec. 4.1. \\ \hline
\end{longtable}


\subsection*{A2.2.9 - RET}
\addcontentsline{toc}{subsection}{A2.2.9 - RET}
La Tabla~\ref{tab:t12-a2-2-9-ret} individualiza los requisitos de este ámbito, su respaldo y la evidencia que deberá comprobarse.
\begin{longtable}{|L{0.09\AnchoMatrizTDoce}|L{0.385\AnchoMatrizTDoce}|L{0.065\AnchoMatrizTDoce}|L{0.18\AnchoMatrizTDoce}|L{0.28\AnchoMatrizTDoce}|}
\caption{Cumplimiento y trazabilidad: A2.2.9 - RET}\label{tab:t12-a2-2-9-ret}\\
\hline
\EncabezadoTabla{ID requerimiento} & \EncabezadoTabla{Descripción} & \EncabezadoTabla{Cumple\newline(Sí/No)} & \EncabezadoTabla{Componente que lo satisface} & \EncabezadoTabla{Sección de la propuesta} \\ \hline
\endfirsthead
\hline
\EncabezadoTabla{ID requerimiento} & \EncabezadoTabla{Descripción} & \EncabezadoTabla{Cumple\newline(Sí/No)} & \EncabezadoTabla{Componente que lo satisface} & \EncabezadoTabla{Sección de la propuesta} \\ \hline
\endhead
\multicolumn{5}{r}{\fontsize{9}{11}\selectfont Continúa en la página siguiente.}\\
\endfoot
\hline
\endlastfoot

RNF-\allowbreak{}RET-\allowbreak{}01 &
\textbf{[Must]} Conservar antecedentes de escuela relativos a menores hasta cumplir 18 años y cinco más, con un mínimo de diez años de retención. &
Sí &
Repositorio de objetos; persistencia por dominio &
SD2, sec. A2.2.9; SD5, sec. 5.2.5. \\ \hline

RNF-\allowbreak{}RET-\allowbreak{}02 &
\textbf{[Must]} Registros de salida, regreso y antecedentes de zarpe: cinco años según la exigencia del caso. &
Sí &
Repositorio de objetos; persistencia por dominio &
SD2, sec. A2.2.9; SD5, sec. 5.2.5. \\ \hline

RNF-\allowbreak{}RET-\allowbreak{}03 &
\textbf{[Must]} Contratos de amarra, custodia y estado de cuenta: seis años según la exigencia del caso. &
Sí &
Repositorio de objetos; persistencia por dominio &
SD2, sec. A2.2.9; SD5, sec. 5.2.5. \\ \hline

RNF-\allowbreak{}RET-\allowbreak{}04 &
\textbf{[Must]} Trazabilidad de residuos peligrosos y actas de derrame: seis años según la exigencia del caso. &
Sí &
Repositorio de objetos; persistencia por dominio &
SD2, sec. A2.2.9; SD5, sec. 5.2.5. \\ \hline

RNF-\allowbreak{}RET-\allowbreak{}05 &
\textbf{[Must]} Series temporales continuas de consumo de agua y energía eléctrica conservadas por al menos 5 años para auditoría. &
Sí &
Repositorio de objetos; persistencia por dominio &
SD2, sec. A2.2.9; SD5, sec. 5.2.5. \\ \hline

RNF-\allowbreak{}RET-\allowbreak{}06 &
\textbf{[Must]} Conservar inspecciones de cada elemento del tren de fondeo durante su vida útil y cinco años adicionales. &
Sí &
Repositorio de objetos; persistencia por dominio &
SD2, sec. A2.2.9; SD5, sec. 5.2.5. \\ \hline

RNF-\allowbreak{}RET-\allowbreak{}07 &
\textbf{[Must]} Registros de despacho en surtidor de combustible y firmas de recepción retenidos por un mínimo de 5 años. &
Sí &
Repositorio de objetos; persistencia por dominio &
SD2, sec. A2.2.9; SD5, sec. 5.2.5. \\ \hline

RNF-\allowbreak{}RET-\allowbreak{}08 &
\textbf{[Must]} Conservar videovigilancia durante seis meses y aplicar eliminación segura verificable conforme a la política de retención. &
Sí &
Plataforma de videovigilancia existente del CLIENTE  &
SD2, sec. A2.2.9; SD4, sec. 4.2; SD5, sec. 5.2.5 y 5.4.1. \\ \hline
\end{longtable}


\subsection*{A2.2.10 - SOP}
\addcontentsline{toc}{subsection}{A2.2.10 - SOP}
La Tabla~\ref{tab:t12-a2-2-10-sop} individualiza los requisitos de este ámbito, su respaldo y la evidencia que deberá comprobarse.
\begin{longtable}{|L{0.09\AnchoMatrizTDoce}|L{0.385\AnchoMatrizTDoce}|L{0.065\AnchoMatrizTDoce}|L{0.18\AnchoMatrizTDoce}|L{0.28\AnchoMatrizTDoce}|}
\caption{Cumplimiento y trazabilidad: A2.2.10 - SOP}\label{tab:t12-a2-2-10-sop}\\
\hline
\EncabezadoTabla{ID requerimiento} & \EncabezadoTabla{Descripción} & \EncabezadoTabla{Cumple\newline(Sí/No)} & \EncabezadoTabla{Componente que lo satisface} & \EncabezadoTabla{Sección de la propuesta} \\ \hline
\endfirsthead
\hline
\EncabezadoTabla{ID requerimiento} & \EncabezadoTabla{Descripción} & \EncabezadoTabla{Cumple\newline(Sí/No)} & \EncabezadoTabla{Componente que lo satisface} & \EncabezadoTabla{Sección de la propuesta} \\ \hline
\endhead
\multicolumn{5}{r}{\fontsize{9}{11}\selectfont Continúa en la página siguiente.}\\
\endfoot
\hline
\endlastfoot

RNF-\allowbreak{}SOP-\allowbreak{}01 &
\textbf{[Must]} Atención todos los días de 06:00 a 24:00 en temporada alta y de 08:00 a 20:00 el resto del año, en español e inglés. &
Sí &
Operaciones y Servicios Gestionados  &
SD2, sec. A2.2.10; SD1, sec. 1.1.3 y 1.5; SD4, sec. 4.1. \\ \hline

RNF-\allowbreak{}SOP-\allowbreak{}02 &
\textbf{[Must]} Canal de notificación de emergencia y recepción de alertas de planes vencidos disponibles 24\(\times\)7\(\times\)365. Esta cobertura técnica no constituye un servicio propio de búsqueda y salvamento. &
Sí &
Operaciones y Servicios Gestionados  &
SD2, sec. A2.2.10; SD1, sec. 1.1.3 y 1.5; SD4, sec. 4.1. \\ \hline

RNF-\allowbreak{}SOP-\allowbreak{}03 &
\textbf{[Must]} Prestación del servicio de operación gestionada, monitoreo y mantenimiento evolutivo durante los 36 meses contractuales. &
Sí &
Operaciones y Servicios Gestionados  &
SD2, sec. A2.2.10; SD1, sec. 1.1.3 y 1.5; SD4, sec. 4.1. \\ \hline
\end{longtable}


\subsection*{A2.2.11 - DISP}
\addcontentsline{toc}{subsection}{A2.2.11 - DISP}
La Tabla~\ref{tab:t12-a2-2-11-disp} individualiza los requisitos de este ámbito, su respaldo y la evidencia que deberá comprobarse.
\begin{longtable}{|L{0.09\AnchoMatrizTDoce}|L{0.385\AnchoMatrizTDoce}|L{0.065\AnchoMatrizTDoce}|L{0.18\AnchoMatrizTDoce}|L{0.28\AnchoMatrizTDoce}|}
\caption{Cumplimiento y trazabilidad: A2.2.11 - DISP}\label{tab:t12-a2-2-11-disp}\\
\hline
\EncabezadoTabla{ID requerimiento} & \EncabezadoTabla{Descripción} & \EncabezadoTabla{Cumple\newline(Sí/No)} & \EncabezadoTabla{Componente que lo satisface} & \EncabezadoTabla{Sección de la propuesta} \\ \hline
\endfirsthead
\hline
\EncabezadoTabla{ID requerimiento} & \EncabezadoTabla{Descripción} & \EncabezadoTabla{Cumple\newline(Sí/No)} & \EncabezadoTabla{Componente que lo satisface} & \EncabezadoTabla{Sección de la propuesta} \\ \hline
\endhead
\multicolumn{5}{r}{\fontsize{9}{11}\selectfont Continúa en la página siguiente.}\\
\endfoot
\hline
\endlastfoot

RNF-\allowbreak{}DISP-\allowbreak{}01 &
\textbf{[Must]} Disponibilidad mensual de transacción crítica de extremo a extremo \(\geq\)99,9\%; componentes de energía, climatización, red, cómputo, base de datos y portal \(\geq\)99,95\%, con la aplicabilidad física justificada. &
Sí &
Núcleos modulares; Célula Operacional Local; DR  &
SD2, sec. A2.2.11; SD4, sec. 4.2 y 4.3. \\ \hline
\end{longtable}


\subsection*{A2.2.12 - MIG}
\addcontentsline{toc}{subsection}{A2.2.12 - MIG}
La Tabla~\ref{tab:t12-a2-2-12-mig} individualiza los requisitos de este ámbito, su respaldo y la evidencia que deberá comprobarse.
\begin{longtable}{|L{0.09\AnchoMatrizTDoce}|L{0.385\AnchoMatrizTDoce}|L{0.065\AnchoMatrizTDoce}|L{0.18\AnchoMatrizTDoce}|L{0.28\AnchoMatrizTDoce}|}
\caption{Cumplimiento y trazabilidad: A2.2.12 - MIG}\label{tab:t12-a2-2-12-mig}\\
\hline
\EncabezadoTabla{ID requerimiento} & \EncabezadoTabla{Descripción} & \EncabezadoTabla{Cumple\newline(Sí/No)} & \EncabezadoTabla{Componente que lo satisface} & \EncabezadoTabla{Sección de la propuesta} \\ \hline
\endfirsthead
\hline
\EncabezadoTabla{ID requerimiento} & \EncabezadoTabla{Descripción} & \EncabezadoTabla{Cumple\newline(Sí/No)} & \EncabezadoTabla{Componente que lo satisface} & \EncabezadoTabla{Sección de la propuesta} \\ \hline
\endhead
\multicolumn{5}{r}{\fontsize{9}{11}\selectfont Continúa en la página siguiente.}\\
\endfoot
\hline
\endlastfoot

RNF-\allowbreak{}MIG-\allowbreak{}01 &
\textbf{[Must]} Migrar todo el padrón de socios/armadores, contratos vigentes con documentos, seis años de pagos y saldos, los expedientes completos de las 14 embarcaciones abandonadas, todo el registro de embarcaciones e invernadas y los últimos tres años de matrícula de escuela. &
Sí &
Migración; dominios de SD5 &
SD2, sec. A2.2.12; SD5, sec. 5.3. \\ \hline
\end{longtable}


\subsection*{A2.2.13 - RED}
\addcontentsline{toc}{subsection}{A2.2.13 - RED}
La Tabla~\ref{tab:t12-a2-2-13-red} individualiza los requisitos de este ámbito, su respaldo y la evidencia que deberá comprobarse.
\begin{longtable}{|L{0.09\AnchoMatrizTDoce}|L{0.385\AnchoMatrizTDoce}|L{0.065\AnchoMatrizTDoce}|L{0.18\AnchoMatrizTDoce}|L{0.28\AnchoMatrizTDoce}|}
\caption{Cumplimiento y trazabilidad: A2.2.13 - RED}\label{tab:t12-a2-2-13-red}\\
\hline
\EncabezadoTabla{ID requerimiento} & \EncabezadoTabla{Descripción} & \EncabezadoTabla{Cumple\newline(Sí/No)} & \EncabezadoTabla{Componente que lo satisface} & \EncabezadoTabla{Sección de la propuesta} \\ \hline
\endfirsthead
\hline
\EncabezadoTabla{ID requerimiento} & \EncabezadoTabla{Descripción} & \EncabezadoTabla{Cumple\newline(Sí/No)} & \EncabezadoTabla{Componente que lo satisface} & \EncabezadoTabla{Sección de la propuesta} \\ \hline
\endhead
\multicolumn{5}{r}{\fontsize{9}{11}\selectfont Continúa en la página siguiente.}\\
\endfoot
\hline
\endlastfoot

RNF-\allowbreak{}RED-\allowbreak{}01 &
\textbf{[Must]} Cobertura de datos en cinco pantalanes y varadero verificada en estructura flotante y rango completo de marea; segregación de socios, administración, cámaras y operación; respaldo del enlace principal. &
Sí &
Célula Operacional Local; implementos de terreno &
SD2, sec. A2.2.13; SD4, sec. 4.2 y 4.2.1. \\ \hline
\end{longtable}


\subsection*{A2.2.14 - TER}
\addcontentsline{toc}{subsection}{A2.2.14 - TER}
La Tabla~\ref{tab:t12-a2-2-14-ter} individualiza los requisitos de este ámbito, su respaldo y la evidencia que deberá comprobarse.
\begin{longtable}{|L{0.09\AnchoMatrizTDoce}|L{0.385\AnchoMatrizTDoce}|L{0.065\AnchoMatrizTDoce}|L{0.18\AnchoMatrizTDoce}|L{0.28\AnchoMatrizTDoce}|}
\caption{Cumplimiento y trazabilidad: A2.2.14 - TER}\label{tab:t12-a2-2-14-ter}\\
\hline
\EncabezadoTabla{ID requerimiento} & \EncabezadoTabla{Descripción} & \EncabezadoTabla{Cumple\newline(Sí/No)} & \EncabezadoTabla{Componente que lo satisface} & \EncabezadoTabla{Sección de la propuesta} \\ \hline
\endfirsthead
\hline
\EncabezadoTabla{ID requerimiento} & \EncabezadoTabla{Descripción} & \EncabezadoTabla{Cumple\newline(Sí/No)} & \EncabezadoTabla{Componente que lo satisface} & \EncabezadoTabla{Sección de la propuesta} \\ \hline
\endhead
\multicolumn{5}{r}{\fontsize{9}{11}\selectfont Continúa en la página siguiente.}\\
\endfoot
\hline
\endlastfoot

RNF-\allowbreak{}TER-\allowbreak{}01 &
\textbf{[Must]} Equipos de pantalán y borde acreditados para atmósfera marina y movimiento; equipos en surtidor acreditados para el recinto clasificado. Cliente adquiere hardware y ejecuta obras; proponente especifica, dimensiona, integra y considera costos y reposición. &
Sí &
Célula Operacional Local; implementos de terreno  &
SD2, sec. A2.2.14; SD4, sec. 4.2 y 4.2.1. \\ \hline
\end{longtable}


\subsection*{A2.2.15 - VEN}
\addcontentsline{toc}{subsection}{A2.2.15 - VEN}
La Tabla~\ref{tab:t12-a2-2-15-ven} individualiza los requisitos de este ámbito, su respaldo y la evidencia que deberá comprobarse.
\begin{longtable}{|L{0.09\AnchoMatrizTDoce}|L{0.385\AnchoMatrizTDoce}|L{0.065\AnchoMatrizTDoce}|L{0.18\AnchoMatrizTDoce}|L{0.28\AnchoMatrizTDoce}|}
\caption{Cumplimiento y trazabilidad: A2.2.15 - VEN}\label{tab:t12-a2-2-15-ven}\\
\hline
\EncabezadoTabla{ID requerimiento} & \EncabezadoTabla{Descripción} & \EncabezadoTabla{Cumple\newline(Sí/No)} & \EncabezadoTabla{Componente que lo satisface} & \EncabezadoTabla{Sección de la propuesta} \\ \hline
\endfirsthead
\hline
\EncabezadoTabla{ID requerimiento} & \EncabezadoTabla{Descripción} & \EncabezadoTabla{Cumple\newline(Sí/No)} & \EncabezadoTabla{Componente que lo satisface} & \EncabezadoTabla{Sección de la propuesta} \\ \hline
\endhead
\multicolumn{5}{r}{\fontsize{9}{11}\selectfont Continúa en la página siguiente.}\\
\endfoot
\hline
\endlastfoot

RNF-\allowbreak{}VEN-\allowbreak{}01 &
\textbf{[Must]} Respetar congelamiento total 15 de diciembre-15 de marzo, campañas de varadero abril-mayo y septiembre-noviembre, 14 fechas de eventos y suspensión por marejadas; absorber cancelaciones sin desplazar hitos contractuales. &
Sí &
Implantación y Gestión del Cambio  &
SD2, sec. A2.2.15; SD5, sec. 5.3.3; SD1, sec. 1.5. \\ \hline
\end{longtable}


\subsection*{A2.2.16 - CAP}
\addcontentsline{toc}{subsection}{A2.2.16 - CAP}
La Tabla~\ref{tab:t12-a2-2-16-cap} individualiza los requisitos de este ámbito, su respaldo y la evidencia que deberá comprobarse.
\begin{longtable}{|L{0.09\AnchoMatrizTDoce}|L{0.385\AnchoMatrizTDoce}|L{0.065\AnchoMatrizTDoce}|L{0.18\AnchoMatrizTDoce}|L{0.28\AnchoMatrizTDoce}|}
\caption{Cumplimiento y trazabilidad: A2.2.16 - CAP}\label{tab:t12-a2-2-16-cap}\\
\hline
\EncabezadoTabla{ID requerimiento} & \EncabezadoTabla{Descripción} & \EncabezadoTabla{Cumple\newline(Sí/No)} & \EncabezadoTabla{Componente que lo satisface} & \EncabezadoTabla{Sección de la propuesta} \\ \hline
\endfirsthead
\hline
\EncabezadoTabla{ID requerimiento} & \EncabezadoTabla{Descripción} & \EncabezadoTabla{Cumple\newline(Sí/No)} & \EncabezadoTabla{Componente que lo satisface} & \EncabezadoTabla{Sección de la propuesta} \\ \hline
\endhead
\multicolumn{5}{r}{\fontsize{9}{11}\selectfont Continúa en la página siguiente.}\\
\endfoot
\hline
\endlastfoot

RNF-\allowbreak{}CAP-\allowbreak{}01 &
\textbf{[Must]} Formación por perfiles, incorporación recurrente de personal nuevo y monitores, sin depender de un formador interno ni capacitar en el peak congelado. &
Sí &
Implantación y Gestión del Cambio; Calidad y Conocimiento  &
SD2, sec. A2.2.16; SD1, sec. 1.3.3 y 1.5. \\ \hline
\end{longtable}


\subsection*{A2.2.17 - CUMP}
\addcontentsline{toc}{subsection}{A2.2.17 - CUMP}
La Tabla~\ref{tab:t12-a2-2-17-cump} individualiza los requisitos de este ámbito, su respaldo y la evidencia que deberá comprobarse.
\begin{longtable}{|L{0.09\AnchoMatrizTDoce}|L{0.385\AnchoMatrizTDoce}|L{0.065\AnchoMatrizTDoce}|L{0.18\AnchoMatrizTDoce}|L{0.28\AnchoMatrizTDoce}|}
\caption{Cumplimiento y trazabilidad: A2.2.17 - CUMP}\label{tab:t12-a2-2-17-cump}\\
\hline
\EncabezadoTabla{ID requerimiento} & \EncabezadoTabla{Descripción} & \EncabezadoTabla{Cumple\newline(Sí/No)} & \EncabezadoTabla{Componente que lo satisface} & \EncabezadoTabla{Sección de la propuesta} \\ \hline
\endfirsthead
\hline
\EncabezadoTabla{ID requerimiento} & \EncabezadoTabla{Descripción} & \EncabezadoTabla{Cumple\newline(Sí/No)} & \EncabezadoTabla{Componente que lo satisface} & \EncabezadoTabla{Sección de la propuesta} \\ \hline
\endhead
\multicolumn{5}{r}{\fontsize{9}{11}\selectfont Continúa en la página siguiente.}\\
\endfoot
\hline
\endlastfoot

RNF-\allowbreak{}CUMP-\allowbreak{}01 &
\textbf{[Must]} Identificar los controles aplicables a los ámbitos de C07 cap. 12 y sec. 16.2, con fuente oficial, vigencia, acto o dato afectado y evidencia. La mención de una ley o producto no acredita cumplimiento. con observancia estricta de la Ley N° 21.719 sobre Protección de Datos Personales y Ley N° 19.628. &
Sí &
Datos y Seguridad; dominios propietarios  &
SD2, sec. A2.2.17; SD1, sec. 1.3.2; SD5, sec. 5.2.5. \\ \hline
\end{longtable}


\clearpage


\section*{Requisitos transversales y parámetros del Caso 07}

\addcontentsline{toc}{section}{Requisitos transversales y parámetros del Caso 07}

Cada código reproduce su obligación BT y agrega los parámetros C07 aplicables, con identificación de componentes y secciones de la propuesta conforme al Formulario T-12.


\subsection*{BT, capítulo 2: MODELO DE ARQUITECTURA DE REFERENCIA}

\addcontentsline{toc}{subsection}{BT, capítulo 2: MODELO DE ARQUITECTURA DE REFERENCIA}

La Tabla~\ref{tab:t12-bt-cap-2} individualiza los requisitos técnicos de este ámbito y su localización en la propuesta.

\begin{longtable}{|L{0.09\AnchoMatrizTDoce}|L{0.385\AnchoMatrizTDoce}|L{0.065\AnchoMatrizTDoce}|L{0.18\AnchoMatrizTDoce}|L{0.28\AnchoMatrizTDoce}|}

\caption{Cumplimiento y trazabilidad: BT, capítulo 2: MODELO DE ARQUITECTURA DE REFERENCIA}\label{tab:t12-bt-cap-2}\\

\hline

\EncabezadoTabla{ID requerimiento} & \EncabezadoTabla{Descripción} & \EncabezadoTabla{Cumple\newline(Sí/No)} & \EncabezadoTabla{Componente que lo satisface} & \EncabezadoTabla{Sección de la propuesta} \\ \hline

\endfirsthead

\hline

\EncabezadoTabla{ID requerimiento} & \EncabezadoTabla{Descripción} & \EncabezadoTabla{Cumple\newline(Sí/No)} & \EncabezadoTabla{Componente que lo satisface} & \EncabezadoTabla{Sección de la propuesta} \\ \hline

\endhead

\multicolumn{5}{r}{\fontsize{9}{11}\selectfont Continúa en la página siguiente.}\\

\endfoot

\hline

\endlastfoot


RT-\allowbreak{}02.\allowbreak{}01 &
La solución se organizará en las ocho capas del numeral 2.1. El PROPONENTE presentará el diagrama de la arquitectura lógica identificando cada capa, sus componentes y las interfaces entre ellas.\newline\textit{BT; Obligatorio.} &
Sí &
Spring Modulith; Arquitectura Modular &
SD4, sec. 4.1 \\ \hline

RT-\allowbreak{}02.\allowbreak{}02 &
La arquitectura será modular, con límites de contexto explícitos y acoplamiento débil. Se rechazará toda arquitectura monolítica que no permita desplegar de forma independiente sus componentes críticos.\newline\textit{BT; Obligatorio.} &
Sí &
Spring Modulith; Arquitectura Modular &
SD4, sec. 4.1 \\ \hline

RT-\allowbreak{}02.\allowbreak{}03 &
La descripción de la arquitectura se ajustará a ISO/IEC/IEEE 42010, con vistas lógica, de procesos, de despliegue, de datos y de seguridad.\newline\textit{BT; Obligatorio.} &
Sí &
Spring Modulith; Arquitectura Modular &
SD4, sec. 4.1 \\ \hline

RT-\allowbreak{}02.\allowbreak{}04 &
El PROPONENTE mantendrá un registro de decisiones de arquitectura (ADR) fechado, con la alternativa escogida, las alternativas descartadas y el criterio de decisión. El registro es entregable contractual y se actualizará durante toda la ejecución.\newline\textit{BT; Obligatorio.} &
Sí &
Spring Modulith; Arquitectura Modular &
SD4, sec. 4.1 \\ \hline

RT-\allowbreak{}02.\allowbreak{}05 &
La capa de servicios de negocio será sin estado. El estado de sesión y el estado de proceso residirán en almacenes externos con alta disponibilidad.\newline\textit{BT; Obligatorio.} &
Sí &
Spring Modulith; Arquitectura Modular &
SD4, sec. 4.1 \\ \hline

RT-\allowbreak{}02.\allowbreak{}06 &
Toda operación de escritura expuesta a reintentos será idempotente, con clave de idempotencia declarada por el cliente y ventana de deduplicación documentada.\newline\textit{BT; Obligatorio.} &
Sí &
Spring Modulith; Arquitectura Modular &
SD4, sec. 4.1 \\ \hline

RT-\allowbreak{}02.\allowbreak{}07 &
Los flujos de eventos garantizarán entrega al menos una vez, con deduplicación en el consumidor y orden garantizado dentro de la partición o del agregado cuando el proceso lo exija.\newline\textit{BT; Obligatorio.} &
Sí &
Spring Modulith; Arquitectura Modular &
SD4, sec. 4.1 \\ \hline

RT-\allowbreak{}02.\allowbreak{}08 &
La solución implementará patrones de resiliencia demostrables: reintento con retroceso exponencial y variación aleatoria, cortacircuitos, mamparos de aislamiento, límites de tasa y tiempo de espera explícito en toda llamada remota. No se admiten llamadas remotas sin tiempo de espera.\newline\textit{BT; Obligatorio.} &
Sí &
Spring Modulith; Arquitectura Modular &
SD4, sec. 4.1 \\ \hline

RT-\allowbreak{}02.\allowbreak{}09 &
La solución degradará de forma elegante: ante la indisponibilidad de un componente no crítico deberá continuar operando en modo reducido, informando la degradación a la persona usuaria, y nunca fallar de forma total.\newline\textit{BT; Obligatorio.} &
Sí &
Spring Modulith; Arquitectura Modular &
SD4, sec. 4.1 \\ \hline

RT-\allowbreak{}02.\allowbreak{}10 &
Las capas de aplicación e integración escalarán horizontalmente de forma automática, con umbrales, límites superiores y costo asociado declarados en la oferta.\newline\textit{BT; Obligatorio.} &
Sí &
Spring Modulith; Arquitectura Modular &
SD4, sec. 4.1 \\ \hline

RT-\allowbreak{}02.\allowbreak{}11 &
El PROPONENTE declarará explícitamente los puntos únicos de falla que subsistan en su arquitectura y justificará por qué son aceptables. Omitir esta declaración cuando existan puntos únicos de falla se evaluará como observación grave.\newline\textit{BT; Obligatorio.} &
Sí &
Spring Modulith; Arquitectura Modular &
SD4, sec. 4.1 \\ \hline

RT-\allowbreak{}02.\allowbreak{}12 &
La solución admitirá su replicación a nuevas unidades, sitios, sucursales o filiales del CLIENTE sin rediseño arquitectónico, mediante parametrización o multi-tenencia.\newline\textit{BT; Según caso.} &
Sí &
Spring Modulith; Arquitectura Modular &
SD4, sec. 4.1 \\ \hline

RT-\allowbreak{}02.\allowbreak{}13 &
El PROPONENTE presentará un modelo de dominio del negocio del caso, con las entidades principales, sus relaciones y los eventos de negocio que las modifican.\newline\textit{BT; Obligatorio.} &
Sí &
Spring Modulith; Arquitectura Modular &
SD4, sec. 4.1 \\ \hline

RT-\allowbreak{}02.\allowbreak{}14 &
Se valorará la aplicación documentada de patrones de arquitectura evolutiva que permitan sustituir un componente sin reescribir la solución: capa anticorrupción frente a sistemas heredados, estrangulamiento progresivo y abstracción de proveedores.\newline\textit{BT; Deseable.} &
No &
No aplica / Descartado: Arquitectura modular monolítica desacoplada mediante Spring Modulith; no requiere estrangulamiento progresivo de microservicios &
SD3, sec. 3.2.2; SD2, sec. 2.5.3 \\ \hline

\end{longtable}


\subsection*{BT, capítulo 3: DESPLIEGUE HÍBRIDO, NUBE Y BORDE OPERACIONAL}

\addcontentsline{toc}{subsection}{BT, capítulo 3: DESPLIEGUE HÍBRIDO, NUBE Y BORDE OPERACIONAL}

La Tabla~\ref{tab:t12-bt-cap-3} individualiza los requisitos técnicos de este ámbito y su localización en la propuesta.

\begin{longtable}{|L{0.09\AnchoMatrizTDoce}|L{0.385\AnchoMatrizTDoce}|L{0.065\AnchoMatrizTDoce}|L{0.18\AnchoMatrizTDoce}|L{0.28\AnchoMatrizTDoce}|}

\caption{Cumplimiento y trazabilidad: BT, capítulo 3: DESPLIEGUE HÍBRIDO, NUBE Y BORDE OPERACIONAL}\label{tab:t12-bt-cap-3}\\

\hline

\EncabezadoTabla{ID requerimiento} & \EncabezadoTabla{Descripción} & \EncabezadoTabla{Cumple\newline(Sí/No)} & \EncabezadoTabla{Componente que lo satisface} & \EncabezadoTabla{Sección de la propuesta} \\ \hline

\endfirsthead

\hline

\EncabezadoTabla{ID requerimiento} & \EncabezadoTabla{Descripción} & \EncabezadoTabla{Cumple\newline(Sí/No)} & \EncabezadoTabla{Componente que lo satisface} & \EncabezadoTabla{Sección de la propuesta} \\ \hline

\endhead

\multicolumn{5}{r}{\fontsize{9}{11}\selectfont Continúa en la página siguiente.}\\

\endfoot

\hline

\endlastfoot


RT-\allowbreak{}03.\allowbreak{}01 &
El PROPONENTE declarará el proveedor de nube pública, la región primaria y la región secundaria utilizadas. El proveedor deberá contar con presencia de región o zona en Chile o en Sudamérica.\newline\textit{BT; Obligatorio.} &
Sí &
AWS Santiago; Nodo de Borde Resiliente (NBR) &
SD4, sec. 4.2 \\ \hline

RT-\allowbreak{}03.\allowbreak{}02 &
Todos los componentes con requisito de alta disponibilidad se desplegarán en al menos dos zonas de disponibilidad. No se aceptará un diseño en una sola zona.\newline\textit{BT; Obligatorio.} &
Sí &
AWS Santiago; Nodo de Borde Resiliente (NBR) &
SD4, sec. 4.2 \\ \hline

RT-\allowbreak{}03.\allowbreak{}03 &
La totalidad de la infraestructura se definirá como código, versionada en el repositorio del CLIENTE, revisable y reproducible. No se admite infraestructura creada manualmente por consola, salvo la cuenta raíz inicial, cuya creación deberá documentarse.\newline\textit{BT; Obligatorio.} &
Sí &
AWS Santiago; Nodo de Borde Resiliente (NBR) &
SD4, sec. 4.2 \\ \hline

RT-\allowbreak{}03.\allowbreak{}04 &
La red se segmentará por capas, con subredes privadas para aplicación y datos, y exposición pública restringida a la capa de borde. Ningún componente de datos será alcanzable desde Internet.\newline\textit{BT; Obligatorio.} &
Sí &
AWS Santiago; Nodo de Borde Resiliente (NBR) &
SD4, sec. 4.2 \\ \hline

RT-\allowbreak{}03.\allowbreak{}05 &
El PROPONENTE privilegiará servicios administrados por sobre servicios autoadministrados cuando ello reduzca el riesgo operacional, y justificará cada excepción.\newline\textit{BT; Obligatorio.} &
Sí &
AWS Santiago; Nodo de Borde Resiliente (NBR) &
SD4, sec. 4.2 \\ \hline

RT-\allowbreak{}03.\allowbreak{}06 &
Se aplicarán prácticas FinOps: etiquetado obligatorio de todos los recursos por ambiente, módulo y centro de costo; presupuestos con alertas de desviación; y reporte mensual de consumo desglosado entregado al CLIENTE.\newline\textit{BT; Obligatorio.} &
Sí &
AWS Santiago; Nodo de Borde Resiliente (NBR) &
SD4, sec. 4.2 \\ \hline

RT-\allowbreak{}03.\allowbreak{}07 &
El PROPONENTE declarará su estrategia de reversibilidad y de mitigación del bloqueo por proveedor, identificando qué componentes son portables, cuáles no lo son y cuál sería el esfuerzo estimado de una migración.\newline\textit{BT; Obligatorio.} &
Sí &
AWS Santiago; Nodo de Borde Resiliente (NBR) &
SD4, sec. 4.2 \\ \hline

RT-\allowbreak{}03.\allowbreak{}08 &
La solución empleará instancias reservadas, planes de ahorro o capacidad comprometida cuando el perfil de carga lo justifique, y lo reflejará en la estructura de costos.\newline\textit{BT; Deseable.} &
No &
No aplica / Descartado: Modelo de cómputo en contenedores elásticos Fargate y Lambdas; no aplica compra anticipada de instancias reservadas EC2 &
SD3, sec. 3.2.2; SD4, sec. 4.2 \\ \hline

RT-\allowbreak{}03.\allowbreak{}09 &
Se valorará el uso de cómputo sin servidor o de contenedores administrados para las cargas de perfil variable, con el análisis comparativo de costo frente a instancias permanentes.\newline\textit{BT; Deseable.} &
No &
No aplica / Descartado: Esquema de dimensionamiento ya optimizado en AWS Fargate según volumetría acotada del puerto (296 naves) &
SD3, sec. 3.2.2; SUP-03 \\ \hline

RT-\allowbreak{}03.\allowbreak{}10 &
El componente on-premise operará de forma autónoma y degradada ante la pérdida total del enlace con la nube, durante un período mínimo de 24 horas continuas o el mayor que fije el caso.\newline\textit{BT; Obligatorio.} &
Sí &
AWS Santiago; Nodo de Borde Resiliente (NBR) &
SD4, sec. 4.2 \\ \hline

RT-\allowbreak{}03.\allowbreak{}11 &
Durante la operación desconectada, la solución continuará registrando las transacciones operacionales críticas de forma local, con integridad garantizada y sin pérdida de datos.\newline\textit{BT; Obligatorio.} &
Sí &
AWS Santiago; Nodo de Borde Resiliente (NBR) &
SD4, sec. 4.2 \\ \hline

RT-\allowbreak{}03.\allowbreak{}12 &
Restablecido el enlace, la sincronización será automática, con reconciliación determinista de conflictos, regla de resolución documentada y bitácora auditable de las decisiones aplicadas.\newline\textit{BT; Obligatorio.} &
Sí &
AWS Santiago; Nodo de Borde Resiliente (NBR) &
SD4, sec. 4.2 \\ \hline

RT-\allowbreak{}03.\allowbreak{}13 &
El PROPONENTE declarará qué funciones NO estarán disponibles en modo desconectado y qué procedimiento manual las suple. La ausencia de esta declaración se evaluará como observación grave.\newline\textit{BT; Obligatorio.} &
Sí &
AWS Santiago; Nodo de Borde Resiliente (NBR) &
SD4, sec. 4.2 \\ \hline

RT-\allowbreak{}03.\allowbreak{}14 &
Los equipos on-premise críticos serán redundantes. El almacenamiento local tolerará la falla de al menos un disco; el PROPONENTE declarará el nivel RAID escogido y lo justificará frente a las alternativas.\newline\textit{BT; Obligatorio.} &
Sí &
AWS Santiago; Nodo de Borde Resiliente (NBR) &
SD4, sec. 4.2 \\ \hline

RT-\allowbreak{}03.\allowbreak{}15 &
Los sistemas on-premise se endurecerán conforme a los CIS Benchmarks aplicables, con gestión centralizada de parches y ventana de aplicación acordada con el CLIENTE.\newline\textit{BT; Obligatorio.} &
Sí &
AWS Santiago; Nodo de Borde Resiliente (NBR) &
SD4, sec. 4.2 \\ \hline

RT-\allowbreak{}03.\allowbreak{}16 &
El monitoreo del componente on-premise se integrará a la misma plataforma de observabilidad que la nube, con alertamiento unificado.\newline\textit{BT; Obligatorio.} &
Sí &
AWS Santiago; Nodo de Borde Resiliente (NBR) &
SD4, sec. 4.2 \\ \hline

RT-\allowbreak{}03.\allowbreak{}17 &
El enlace entre el sitio on-premise y la nube será redundante, con caminos físicos y proveedores distintos, y conmutación automática con tiempo de conmutación declarado.\newline\textit{BT; Obligatorio.} &
Sí &
AWS Santiago; Nodo de Borde Resiliente (NBR) &
SD4, sec. 4.2 \\ \hline

RT-\allowbreak{}03.\allowbreak{}18 &
Los dispositivos de borde y de terreno se administrarán de forma remota y centralizada: inventario, configuración, actualización de firmware y de aplicación, bloqueo y borrado remoto.\newline\textit{BT; Obligatorio.} &
Sí &
AWS Santiago; Nodo de Borde Resiliente (NBR) &
SD4, sec. 4.2 \\ \hline

RT-\allowbreak{}03.\allowbreak{}19 &
Se valorará el procesamiento en el borde de las cargas que lo admitan —filtrado, agregación previa, inferencia local— reduciendo el volumen transferido y la dependencia del enlace.\newline\textit{BT; Deseable.} &
No &
No aplica / Descartado: Inferencia de IA pesada se ejecuta centralizada en nube; borde NBR opera lógica transaccional y persistencia determinista local &
SD3, sec. 3.2.2; SD4, sec. 4.2 \\ \hline

RT-\allowbreak{}03.\allowbreak{}20 &
El PROPONENTE dimensionará el ancho de banda requerido por sitio, en régimen normal y en peak, y lo justificará con el cálculo de volumen de transacciones y de datos.\newline\textit{BT; Obligatorio.} &
Sí &
AWS Santiago; Nodo de Borde Resiliente (NBR) &
SD4, sec. 4.2 \\ \hline

RT-\allowbreak{}03.\allowbreak{}21 &
La conexión entre la red del CLIENTE y la nube se establecerá mediante enlace privado dedicado o red privada virtual con cifrado, según lo que el volumen y la criticidad justifiquen.\newline\textit{BT; Obligatorio.} &
Sí &
AWS Santiago; Nodo de Borde Resiliente (NBR) &
SD4, sec. 4.2 \\ \hline

RT-\allowbreak{}03.\allowbreak{}22 &
El acceso remoto de las personas trabajadoras del CLIENTE, incluido el trabajo desde el hogar, se resolverá con acceso a la red de confianza cero, con verificación de postura del dispositivo. No se admite exponer servicios internos directamente a Internet.\newline\textit{BT; Obligatorio.} &
Sí &
AWS Santiago; Nodo de Borde Resiliente (NBR) &
SD4, sec. 4.2 \\ \hline

RT-\allowbreak{}03.\allowbreak{}23 &
La red inalámbrica de los sitios operacionales, cuando el caso la requiera, contará con segmentación por tipo de dispositivo, autenticación por certificado o credencial de empresa y cobertura verificada mediante estudio de sitio.\newline\textit{BT; Según caso.} &
Sí &
AWS Santiago; Nodo de Borde Resiliente (NBR) &
SD4, sec. 4.2 \\ \hline

RT-\allowbreak{}03.\allowbreak{}24 &
El PROPONENTE declarará la calidad de servicio y la priorización de tráfico aplicada a las transacciones operacionales críticas frente al tráfico administrativo.\newline\textit{BT; Deseable.} &
No &
No aplica / Descartado: Red de área local del recinto opera con switches administrados estándar sin modelado complejo de QoS por clases &
SD3, sec. 3.2.2; SUP-04 \\ \hline

\end{longtable}


\subsection*{BT, capítulo 4: AMBIENTES, ENTREGA CONTINUA Y GESTIÓN DE LA CONFIGURACIÓN}

\addcontentsline{toc}{subsection}{BT, capítulo 4: AMBIENTES, ENTREGA CONTINUA Y GESTIÓN DE LA CONFIGURACIÓN}

La Tabla~\ref{tab:t12-bt-cap-4} individualiza los requisitos técnicos de este ámbito y su localización en la propuesta.

\begin{longtable}{|L{0.09\AnchoMatrizTDoce}|L{0.385\AnchoMatrizTDoce}|L{0.065\AnchoMatrizTDoce}|L{0.18\AnchoMatrizTDoce}|L{0.28\AnchoMatrizTDoce}|}

\caption{Cumplimiento y trazabilidad: BT, capítulo 4: AMBIENTES, ENTREGA CONTINUA Y GESTIÓN DE LA CONFIGURACIÓN}\label{tab:t12-bt-cap-4}\\

\hline

\EncabezadoTabla{ID requerimiento} & \EncabezadoTabla{Descripción} & \EncabezadoTabla{Cumple\newline(Sí/No)} & \EncabezadoTabla{Componente que lo satisface} & \EncabezadoTabla{Sección de la propuesta} \\ \hline

\endfirsthead

\hline

\EncabezadoTabla{ID requerimiento} & \EncabezadoTabla{Descripción} & \EncabezadoTabla{Cumple\newline(Sí/No)} & \EncabezadoTabla{Componente que lo satisface} & \EncabezadoTabla{Sección de la propuesta} \\ \hline

\endhead

\multicolumn{5}{r}{\fontsize{9}{11}\selectfont Continúa en la página siguiente.}\\

\endfoot

\hline

\endlastfoot


RT-\allowbreak{}04.\allowbreak{}01 &
Los cinco ambientes del numeral 4.1 estarán habilitados y operativos como condición del hito H3 del Formulario E-25.\newline\textit{BT; Obligatorio.} &
Sí &
GitHub Actions; CI/CD; OpenTofu; SonarQube &
SD6, sec. 6.2 \\ \hline

RT-\allowbreak{}04.\allowbreak{}02 &
Preproducción será equivalente a producción en topología, versiones de componentes y configuración. Las diferencias que subsistan por costo se declararán expresamente y se justificarán.\newline\textit{BT; Obligatorio.} &
Sí &
GitHub Actions; CI/CD; OpenTofu; SonarQube &
SD6, sec. 6.2 \\ \hline

RT-\allowbreak{}04.\allowbreak{}03 &
El código residirá en un sistema de control de versiones con ramas protegidas, revisión obligatoria por pares y prohibición de escritura directa sobre la rama principal.\newline\textit{BT; Obligatorio.} &
Sí &
GitHub Actions; CI/CD; OpenTofu; SonarQube &
SD6, sec. 6.2 \\ \hline

RT-\allowbreak{}04.\allowbreak{}04 &
Existirá trazabilidad completa entre requerimiento, incidencia, cambio de código, prueba ejecutada y despliegue realizado.\newline\textit{BT; Obligatorio.} &
Sí &
GitHub Actions; CI/CD; OpenTofu; SonarQube &
SD6, sec. 6.2 \\ \hline

RT-\allowbreak{}04.\allowbreak{}05 &
El flujo de integración continua ejecutará, como mínimo: compilación, pruebas unitarias, análisis estático de código, análisis de composición de software, escaneo de secretos y escaneo de imágenes de contenedor, con criterios de bloqueo automático del despliegue.\newline\textit{BT; Obligatorio.} &
Sí &
GitHub Actions; CI/CD; OpenTofu; SonarQube &
SD6, sec. 6.2 \\ \hline

RT-\allowbreak{}04.\allowbreak{}06 &
Los despliegues serán automatizados y reproducibles, con reversión automatizada y sin intervención manual en el paso a producción.\newline\textit{BT; Obligatorio.} &
Sí &
GitHub Actions; CI/CD; OpenTofu; SonarQube &
SD6, sec. 6.2 \\ \hline

RT-\allowbreak{}04.\allowbreak{}07 &
La estrategia de despliegue permitirá liberar sin interrupción del servicio: azul-verde, canario o despliegue progresivo. Se declarará cuál se emplea y se demostrará en Preproducción antes de cada paso a producción.\newline\textit{BT; Obligatorio.} &
Sí &
GitHub Actions; CI/CD; OpenTofu; SonarQube &
SD6, sec. 6.2 \\ \hline

RT-\allowbreak{}04.\allowbreak{}08 &
Toda configuración estará externalizada del artefacto y gestionada por ambiente. Un mismo artefacto deberá poder promoverse de QA a Preproducción y a Producción sin recompilación.\newline\textit{BT; Obligatorio.} &
Sí &
GitHub Actions; CI/CD; OpenTofu; SonarQube &
SD6, sec. 6.2 \\ \hline

RT-\allowbreak{}04.\allowbreak{}09 &
Los secretos residirán en un gestor de secretos con rotación automática y auditoría de acceso. Queda prohibida toda credencial embebida en código, imágenes o archivos de configuración.\newline\textit{BT; Obligatorio.} &
Sí &
GitHub Actions; CI/CD; OpenTofu; SonarQube &
SD6, sec. 6.2 \\ \hline

RT-\allowbreak{}04.\allowbreak{}10 &
Las migraciones de esquema de base de datos serán versionadas, reversibles y ejecutadas de forma automatizada, con estrategia de compatibilidad que permita convivir dos versiones de la aplicación durante el despliegue.\newline\textit{BT; Obligatorio.} &
Sí &
GitHub Actions; CI/CD; OpenTofu; SonarQube &
SD6, sec. 6.2 \\ \hline

RT-\allowbreak{}04.\allowbreak{}11 &
La cobertura de pruebas automatizadas del código de lógica de negocio será de al menos 70\%, con umbral bloqueante en el flujo de integración continua.\newline\textit{BT; Obligatorio.} &
Sí &
GitHub Actions; CI/CD; OpenTofu; SonarQube &
SD6, sec. 6.2 \\ \hline

RT-\allowbreak{}04.\allowbreak{}12 &
El PROPONENTE declarará su frecuencia de despliegue objetivo, su tiempo desde el compromiso de código hasta producción, su tasa de cambios fallidos y su tiempo de restauración, y los medirá durante la Operación.\newline\textit{BT; Obligatorio.} &
Sí &
GitHub Actions; CI/CD; OpenTofu; SonarQube &
SD6, sec. 6.2 \\ \hline

RT-\allowbreak{}04.\allowbreak{}13 &
Los ambientes no productivos se apagarán o reducirán fuera del horario de uso, con el ahorro reflejado en la estructura de costos.\newline\textit{BT; Deseable.} &
Sí &
Automatización con AWS EventBridge y Lambda para apagado nocturno DEV/QA &
SD4, sec. 4.2; SD11, sec. 11.2 \\ \hline

RT-\allowbreak{}04.\allowbreak{}14 &
Se valorará la existencia de ambientes efímeros por rama o por incidencia, creados y destruidos automáticamente.\newline\textit{BT; Deseable.} &
No &
No aplica / Descartado: Se priorizan ambientes estables DEV, QA y PREPROD por presupuesto acotado y cliente sin área DevOps interna &
SD3, sec. 3.2.2; DEC-22 \\ \hline

\end{longtable}


\subsection*{BT, capítulo 5: DATOS, INTEGRACIÓN E INTEROPERABILIDAD}

\addcontentsline{toc}{subsection}{BT, capítulo 5: DATOS, INTEGRACIÓN E INTEROPERABILIDAD}

La Tabla~\ref{tab:t12-bt-cap-5} individualiza los requisitos técnicos de este ámbito y su localización en la propuesta.

\begin{longtable}{|L{0.09\AnchoMatrizTDoce}|L{0.385\AnchoMatrizTDoce}|L{0.065\AnchoMatrizTDoce}|L{0.18\AnchoMatrizTDoce}|L{0.28\AnchoMatrizTDoce}|}

\caption{Cumplimiento y trazabilidad: BT, capítulo 5: DATOS, INTEGRACIÓN E INTEROPERABILIDAD}\label{tab:t12-bt-cap-5}\\

\hline

\EncabezadoTabla{ID requerimiento} & \EncabezadoTabla{Descripción} & \EncabezadoTabla{Cumple\newline(Sí/No)} & \EncabezadoTabla{Componente que lo satisface} & \EncabezadoTabla{Sección de la propuesta} \\ \hline

\endfirsthead

\hline

\EncabezadoTabla{ID requerimiento} & \EncabezadoTabla{Descripción} & \EncabezadoTabla{Cumple\newline(Sí/No)} & \EncabezadoTabla{Componente que lo satisface} & \EncabezadoTabla{Sección de la propuesta} \\ \hline

\endhead

\multicolumn{5}{r}{\fontsize{9}{11}\selectfont Continúa en la página siguiente.}\\

\endfoot

\hline

\endlastfoot


RT-\allowbreak{}05.\allowbreak{}01 &
El PROPONENTE entregará el modelo de datos documentado y un diccionario de datos con el nombre, el tipo, el dominio de valores, la obligatoriedad, el propietario y la sensibilidad de cada atributo.\newline\textit{BT; Obligatorio.} &
Sí &
PostgreSQL 18; OpenAPI 3.1; AsyncAPI 2.6 &
SD5, sec. 5.1 y sec. 5.2 \\ \hline

RT-\allowbreak{}05.\allowbreak{}02 &
El PROPONENTE justificará la selección del paradigma y del motor de persistencia: relacional o no relacional, garantías transaccionales, y la posición escogida entre consistencia y disponibilidad conforme al teorema CAP, para cada dominio de datos.\newline\textit{BT; Obligatorio.} &
Sí &
PostgreSQL 18; OpenAPI 3.1; AsyncAPI 2.6 &
SD5, sec. 5.1 y sec. 5.2 \\ \hline

RT-\allowbreak{}05.\allowbreak{}03 &
Toda operación de negocio será trazable: la solución permitirá reconstruir quién, qué, cuándo, desde qué dispositivo y con qué valores anteriores y posteriores, para cualquier registro y en cualquier momento del período de retención.\newline\textit{BT; Obligatorio.} &
Sí &
PostgreSQL 18; OpenAPI 3.1; AsyncAPI 2.6 &
SD5, sec. 5.1 y sec. 5.2 \\ \hline

RT-\allowbreak{}05.\allowbreak{}04 &
La calidad de datos se gestionará conforme a ISO/IEC 25012, con validación en el punto de captura, indicadores de completitud, exactitud y consistencia, y tablero de calidad disponible para el CLIENTE.\newline\textit{BT; Obligatorio.} &
Sí &
PostgreSQL 18; OpenAPI 3.1; AsyncAPI 2.6 &
SD5, sec. 5.1 y sec. 5.2 \\ \hline

RT-\allowbreak{}05.\allowbreak{}05 &
El almacenamiento transaccional y el analítico estarán separados. Ninguna consulta analítica podrá degradar el desempeño de la operación.\newline\textit{BT; Obligatorio.} &
Sí &
PostgreSQL 18; OpenAPI 3.1; AsyncAPI 2.6 &
SD5, sec. 5.1 y sec. 5.2 \\ \hline

RT-\allowbreak{}05.\allowbreak{}06 &
La solución permitirá exportar la totalidad de la información del CLIENTE en formatos abiertos y documentados, en cualquier momento del Contrato, sin costo adicional y sin intervención del ADJUDICATARIO.\newline\textit{BT; Obligatorio.} &
Sí &
PostgreSQL 18; OpenAPI 3.1; AsyncAPI 2.6 &
SD5, sec. 5.1 y sec. 5.2 \\ \hline

RT-\allowbreak{}05.\allowbreak{}07 &
El PROPONENTE declarará la política de retención, archivado y eliminación por cada dominio de datos, coherente con la normativa aplicable al caso, e implementará un procedimiento verificable de eliminación segura.\newline\textit{BT; Obligatorio.} &
Sí &
PostgreSQL 18; OpenAPI 3.1; AsyncAPI 2.6 &
SD5, sec. 5.1 y sec. 5.2 \\ \hline

RT-\allowbreak{}05.\allowbreak{}08 &
Los datos personales se tratarán conforme al Artículo 85° de las Bases Administrativas, con seudonimización o cifrado a nivel de campo para las categorías sensibles que el caso identifique.\newline\textit{BT; Obligatorio.} &
Sí &
PostgreSQL 18; OpenAPI 3.1; AsyncAPI 2.6 &
SD5, sec. 5.1 y sec. 5.2 \\ \hline

RT-\allowbreak{}05.\allowbreak{}09 &
El PROPONENTE presentará una estrategia de gestión de datos maestros que evite la duplicación de entidades compartidas entre módulos y con sistemas externos.\newline\textit{BT; Obligatorio.} &
Sí &
PostgreSQL 18; OpenAPI 3.1; AsyncAPI 2.6 &
SD5, sec. 5.1 y sec. 5.2 \\ \hline

RT-\allowbreak{}05.\allowbreak{}10 &
Se valorará la implementación de un catálogo de datos con linaje automatizado, que permita rastrear el origen de cada indicador de negocio hasta su fuente.\newline\textit{BT; Deseable.} &
No &
No aplica / Descartado: Catálogo de datos manual y diccionario centralizado en PostgreSQL; linaje dinámico automatizado excede escala del caso &
SD3, sec. 3.2.2; SD5, sec. 5.1 \\ \hline

RT-\allowbreak{}05.\allowbreak{}11 &
El PROPONENTE presentará un plan de migración con alcance, origen, volumen, reglas de transformación, criterios de calidad, estrategia de ejecución y plan de reversión.\newline\textit{BT; Obligatorio.} &
Sí &
PostgreSQL 18; OpenAPI 3.1; AsyncAPI 2.6 &
SD5, sec. 5.1 y sec. 5.2 \\ \hline

RT-\allowbreak{}05.\allowbreak{}12 &
La migración incluirá una etapa de perfilado y saneamiento previo, con informe de los defectos detectados en los datos de origen y la decisión adoptada sobre cada uno.\newline\textit{BT; Obligatorio.} &
Sí &
PostgreSQL 18; OpenAPI 3.1; AsyncAPI 2.6 &
SD5, sec. 5.1 y sec. 5.2 \\ \hline

RT-\allowbreak{}05.\allowbreak{}13 &
Se ejecutarán al menos dos ensayos completos de migración sobre Preproducción antes de la migración definitiva, con medición del tiempo total y del resultado de la conciliación.\newline\textit{BT; Obligatorio.} &
Sí &
PostgreSQL 18; OpenAPI 3.1; AsyncAPI 2.6 &
SD5, sec. 5.1 y sec. 5.2 \\ \hline

RT-\allowbreak{}05.\allowbreak{}14 &
La conciliación posterior a la migración será cuantitativa y verificable: recuentos, sumas de control y muestreo dirigido. Toda diferencia deberá quedar explicada.\newline\textit{BT; Obligatorio.} &
Sí &
PostgreSQL 18; OpenAPI 3.1; AsyncAPI 2.6 &
SD5, sec. 5.1 y sec. 5.2 \\ \hline

RT-\allowbreak{}05.\allowbreak{}15 &
Los datos históricos que no se migren quedarán accesibles en un repositorio de consulta durante el período de retención que fije el caso.\newline\textit{BT; Según caso.} &
Sí &
PostgreSQL 18; OpenAPI 3.1; AsyncAPI 2.6 &
SD5, sec. 5.1 y sec. 5.2 \\ \hline

RT-\allowbreak{}05.\allowbreak{}16 &
Los servicios síncronos se documentarán en OpenAPI 3.1 y los flujos dirigidos por eventos en AsyncAPI 2.6 o superior. La documentación se generará desde el código y se mantendrá actualizada automáticamente.\newline\textit{BT; Obligatorio.} &
Sí &
PostgreSQL 18; OpenAPI 3.1; AsyncAPI 2.6 &
SD5, sec. 5.1 y sec. 5.2 \\ \hline

RT-\allowbreak{}05.\allowbreak{}17 &
Los contratos de interfaz se versionarán semánticamente, con compatibilidad hacia atrás y política de obsolescencia con preaviso mínimo de seis meses.\newline\textit{BT; Obligatorio.} &
Sí &
PostgreSQL 18; OpenAPI 3.1; AsyncAPI 2.6 &
SD5, sec. 5.1 y sec. 5.2 \\ \hline

RT-\allowbreak{}05.\allowbreak{}18 &
La autenticación entre sistemas empleará OAuth 2.1 con credenciales de cliente o autenticación mutua TLS. Queda prohibida la autenticación por clave estática en la ruta de la dirección web.\newline\textit{BT; Obligatorio.} &
Sí &
PostgreSQL 18; OpenAPI 3.1; AsyncAPI 2.6 &
SD5, sec. 5.1 y sec. 5.2 \\ \hline

RT-\allowbreak{}05.\allowbreak{}19 &
Toda integración registrará la transacción de entrada y de salida, con identificador de correlación común que permita seguir una operación de negocio a través de todos los sistemas involucrados.\newline\textit{BT; Obligatorio.} &
Sí &
PostgreSQL 18; OpenAPI 3.1; AsyncAPI 2.6 &
SD5, sec. 5.1 y sec. 5.2 \\ \hline

RT-\allowbreak{}05.\allowbreak{}20 &
Las integraciones con sistemas heredados o de terceros se aislarán mediante una capa anticorrupción, de modo que un cambio en el sistema externo no propague su modelo al núcleo de la solución.\newline\textit{BT; Obligatorio.} &
Sí &
PostgreSQL 18; OpenAPI 3.1; AsyncAPI 2.6 &
SD5, sec. 5.1 y sec. 5.2 \\ \hline

RT-\allowbreak{}05.\allowbreak{}21 &
El PROPONENTE declarará, por cada integración, el modo (síncrono o asíncrono), el volumen esperado, la ventana de disponibilidad del sistema contraparte y el comportamiento de la solución cuando ese sistema no responde.\newline\textit{BT; Obligatorio.} &
Sí &
PostgreSQL 18; OpenAPI 3.1; AsyncAPI 2.6 &
SD5, sec. 5.1 y sec. 5.2 \\ \hline

RT-\allowbreak{}05.\allowbreak{}22 &
La solución soportará la carga y descarga masiva de información en formatos abiertos, con validación previa, informe de errores por registro y procesamiento parcial.\newline\textit{BT; Obligatorio.} &
Sí &
PostgreSQL 18; OpenAPI 3.1; AsyncAPI 2.6 &
SD5, sec. 5.1 y sec. 5.2 \\ \hline

RT-\allowbreak{}05.\allowbreak{}23 &
Se emplearán los estándares sectoriales de intercambio que las Bases Técnicas del caso identifiquen.\newline\textit{BT; Según caso.} &
Sí &
PostgreSQL 18; OpenAPI 3.1; AsyncAPI 2.6 &
SD5, sec. 5.1 y sec. 5.2 \\ \hline

RT-\allowbreak{}05.\allowbreak{}24 &
Se valorará la publicación de un portal de servicios para desarrolladores con documentación navegable, ambiente de pruebas y credenciales de prueba autoservidas.\newline\textit{BT; Deseable.} &
No &
No aplica / Descartado: Modelo de integración B2B cerrado del puerto (armadores, contratistas y autoridades); no expone APIs a terceros abiertos &
SD3, sec. 3.2.2; sec. 2.5.3 \\ \hline

RT-\allowbreak{}05.\allowbreak{}25 &
La solución proveerá una capa analítica con tableros operacionales y de gestión, construidos sobre los indicadores que las Bases Técnicas del caso definan.\newline\textit{BT; Obligatorio.} &
Sí &
PostgreSQL 18; OpenAPI 3.1; AsyncAPI 2.6 &
SD5, sec. 5.1 y sec. 5.2 \\ \hline

RT-\allowbreak{}05.\allowbreak{}26 &
Los tableros permitirán filtrar por período, unidad organizacional y dimensiones propias del caso, y profundizar desde el indicador agregado hasta la transacción de origen.\newline\textit{BT; Obligatorio.} &
Sí &
PostgreSQL 18; OpenAPI 3.1; AsyncAPI 2.6 &
SD5, sec. 5.1 y sec. 5.2 \\ \hline

RT-\allowbreak{}05.\allowbreak{}27 &
El CLIENTE podrá construir sus propios informes sin intervención del ADJUDICATARIO, mediante una herramienta de autoservicio con modelo semántico documentado.\newline\textit{BT; Obligatorio.} &
Sí &
PostgreSQL 18; OpenAPI 3.1; AsyncAPI 2.6 &
SD5, sec. 5.1 y sec. 5.2 \\ \hline

RT-\allowbreak{}05.\allowbreak{}28 &
Todo informe será exportable en formatos abiertos y programable para envío automático por calendario.\newline\textit{BT; Obligatorio.} &
Sí &
PostgreSQL 18; OpenAPI 3.1; AsyncAPI 2.6 &
SD5, sec. 5.1 y sec. 5.2 \\ \hline

RT-\allowbreak{}05.\allowbreak{}29 &
La latencia máxima entre la ocurrencia de una transacción y su disponibilidad en la capa analítica será la que fije el caso y, en su defecto, no superará las 4 horas.\newline\textit{BT; Según caso.} &
Sí &
PostgreSQL 18; OpenAPI 3.1; AsyncAPI 2.6 &
SD5, sec. 5.1 y sec. 5.2 \\ \hline

RT-\allowbreak{}05.\allowbreak{}30 &
Se valorará la incorporación de analítica predictiva pertinente al proceso del caso, con el modelo, sus variables, su métrica de desempeño y su plan de reentrenamiento documentados.\newline\textit{BT; Deseable.} &
No &
No aplica / Descartado: Analítica predictiva se centraliza en modelos del Módulo de Innovación sin plataforma MLOps dedicada separada &
SD3, sec. 3.2.2; SD13, sec. 13.1 \\ \hline

\end{longtable}


\subsection*{BT, capítulo 6: SITE PRINCIPAL ON-PREMISE Y EMPLAZAMIENTO}

\addcontentsline{toc}{subsection}{BT, capítulo 6: SITE PRINCIPAL ON-PREMISE Y EMPLAZAMIENTO}

La Tabla~\ref{tab:t12-bt-cap-6} individualiza los requisitos técnicos de este ámbito y su localización en la propuesta.

\begin{longtable}{|L{0.09\AnchoMatrizTDoce}|L{0.385\AnchoMatrizTDoce}|L{0.065\AnchoMatrizTDoce}|L{0.18\AnchoMatrizTDoce}|L{0.28\AnchoMatrizTDoce}|}

\caption{Cumplimiento y trazabilidad: BT, capítulo 6: SITE PRINCIPAL ON-PREMISE Y EMPLAZAMIENTO}\label{tab:t12-bt-cap-6}\\

\hline

\EncabezadoTabla{ID requerimiento} & \EncabezadoTabla{Descripción} & \EncabezadoTabla{Cumple\newline(Sí/No)} & \EncabezadoTabla{Componente que lo satisface} & \EncabezadoTabla{Sección de la propuesta} \\ \hline

\endfirsthead

\hline

\EncabezadoTabla{ID requerimiento} & \EncabezadoTabla{Descripción} & \EncabezadoTabla{Cumple\newline(Sí/No)} & \EncabezadoTabla{Componente que lo satisface} & \EncabezadoTabla{Sección de la propuesta} \\ \hline

\endhead

\multicolumn{5}{r}{\fontsize{9}{11}\selectfont Continúa en la página siguiente.}\\

\endfoot

\hline

\endlastfoot


RT-\allowbreak{}06.\allowbreak{}01 &
El espacio asignado será de uso exclusivo de la solución y estará aislado de otras dependencias del CLIENTE, con acceso independiente.\newline\textit{BT; Obligatorio.} &
Sí &
Gabinete IP55; UPS Industrial; Nodo de Borde &
SD4, sec. 4.2 \\ \hline

RT-\allowbreak{}06.\allowbreak{}02 &
Los muros no estructurales del recinto contarán con blindaje perimetral; el PROPONENTE especificará el material y la resistencia.\newline\textit{BT; Obligatorio.} &
Sí &
Gabinete IP55; UPS Industrial; Nodo de Borde &
SD4, sec. 4.2 \\ \hline

RT-\allowbreak{}06.\allowbreak{}03 &
El PROPONENTE entregará el plano de distribución interna del recinto, con la separación de las zonas de generadores, baterías, climatización, servidores, comunicaciones, trabajo y respaldo.\newline\textit{BT; Obligatorio.} &
Sí &
Gabinete IP55; UPS Industrial; Nodo de Borde &
SD4, sec. 4.2 \\ \hline

RT-\allowbreak{}06.\allowbreak{}04 &
El piso técnico, la canalización, el cableado estructurado y el etiquetado se ejecutarán conforme a norma, con documentación de la certificación de cada enlace.\newline\textit{BT; Obligatorio.} &
Sí &
Gabinete IP55; UPS Industrial; Nodo de Borde &
SD4, sec. 4.2 \\ \hline

RT-\allowbreak{}06.\allowbreak{}05 &
Los racks de servidores serán independientes de los racks de equipos de comunicación. Se declarará la ocupación proyectada de cada rack y su margen de crecimiento.\newline\textit{BT; Obligatorio.} &
Sí &
Gabinete IP55; UPS Industrial; Nodo de Borde &
SD4, sec. 4.2 \\ \hline

RT-\allowbreak{}06.\allowbreak{}06 &
La obra civil de separación de las instalaciones es de cargo del CLIENTE; su especificación técnica y su coordinación son de cargo del PROPONENTE.\newline\textit{BT; Obligatorio.} &
Sí &
Gabinete IP55; UPS Industrial; Nodo de Borde &
SD4, sec. 4.2 \\ \hline

RT-\allowbreak{}06.\allowbreak{}07 &
El suministro eléctrico de los equipos será ininterrumpido, con sistema de alimentación ininterrumpida dimensionado para una autonomía mínima de 30 minutos a plena carga.\newline\textit{BT; Obligatorio.} &
Sí &
Gabinete IP55; UPS Industrial; Nodo de Borde &
SD4, sec. 4.2 \\ \hline

RT-\allowbreak{}06.\allowbreak{}08 &
La capacidad de generación autónoma asegurará un rango mínimo de 24 horas continuas de operación, con estanque de combustible dimensionado y contrato de reabastecimiento declarado.\newline\textit{BT; Obligatorio.} &
Sí &
Gabinete IP55; UPS Industrial; Nodo de Borde &
SD4, sec. 4.2 \\ \hline

RT-\allowbreak{}06.\allowbreak{}09 &
La instalación eléctrica del recinto será independiente de la del resto del edificio y cumplirá la normativa eléctrica chilena vigente, incluida la NCh Elec. 2777 sobre sistemas de puesta a tierra.\newline\textit{BT; Obligatorio.} &
Sí &
Gabinete IP55; UPS Industrial; Nodo de Borde &
SD4, sec. 4.2 \\ \hline

RT-\allowbreak{}06.\allowbreak{}10 &
Se efectuará revisión y medición semestral de las instalaciones eléctricas del recinto, con informe entregable al CLIENTE.\newline\textit{BT; Obligatorio.} &
Sí &
Gabinete IP55; UPS Industrial; Nodo de Borde &
SD4, sec. 4.2 \\ \hline

RT-\allowbreak{}06.\allowbreak{}11 &
El PROPONENTE declarará la carga eléctrica proyectada en kW, el factor de potencia y la eficiencia en el uso de la energía (PUE) estimada del recinto.\newline\textit{BT; Obligatorio.} &
Sí &
Gabinete IP55; UPS Industrial; Nodo de Borde &
SD4, sec. 4.2 \\ \hline

RT-\allowbreak{}06.\allowbreak{}12 &
Se valorará la redundancia de alimentación en configuración 2N o N+1 con doble acometida y transferencia automática.\newline\textit{BT; Deseable.} &
No &
No aplica / Descartado: Red eléctrica pública no dispone de doble acometida independiente en Bahía Ilque; se respalda con UPS industrial de 24 h y generador &
SD3, sec. 3.2.2; SUP-04 \\ \hline

RT-\allowbreak{}06.\allowbreak{}13 &
El recinto contará con climatización de precisión para operación continua, redundante en configuración N+1, con control de temperatura y de humedad relativa dentro de los rangos que recomienda el fabricante del equipamiento.\newline\textit{BT; Obligatorio.} &
Sí &
Gabinete IP55; UPS Industrial; Nodo de Borde &
SD4, sec. 4.2 \\ \hline

RT-\allowbreak{}06.\allowbreak{}14 &
Se monitorearán en línea la temperatura, la humedad y la presencia de agua, con alertamiento integrado a la plataforma de observabilidad.\newline\textit{BT; Obligatorio.} &
Sí &
Gabinete IP55; UPS Industrial; Nodo de Borde &
SD4, sec. 4.2 \\ \hline

RT-\allowbreak{}06.\allowbreak{}15 &
El PROPONENTE declarará la estrategia de contención de pasillo frío o caliente y su efecto en la eficiencia energética.\newline\textit{BT; Deseable.} &
No &
No aplica / Descartado: Emplazamiento en gabinete mural IP55 estanco en oficina técnica; no aplica confinamiento de pasillo frío de datacenter &
SD3, sec. 3.2.2; SD4, sec. 4.2 \\ \hline

RT-\allowbreak{}06.\allowbreak{}16 &
El recinto contará con detección temprana por aspiración de aire con tecnología láser, tipo AnaLASER o equivalente de prestaciones iguales o superiores.\newline\textit{BT; Obligatorio.} &
Sí &
Gabinete IP55; UPS Industrial; Nodo de Borde &
SD4, sec. 4.2 \\ \hline

RT-\allowbreak{}06.\allowbreak{}17 &
La extinción será automática mediante agente limpio tipo FM-200 o equivalente, con aprobación UL e instalación conforme a norma NFPA.\newline\textit{BT; Obligatorio.} &
Sí &
Gabinete IP55; UPS Industrial; Nodo de Borde &
SD4, sec. 4.2 \\ \hline

RT-\allowbreak{}06.\allowbreak{}18 &
Se proveerá un sistema secundario de extintores portátiles habilitados, con mantención y certificación vigentes.\newline\textit{BT; Obligatorio.} &
Sí &
Gabinete IP55; UPS Industrial; Nodo de Borde &
SD4, sec. 4.2 \\ \hline

RT-\allowbreak{}06.\allowbreak{}19 &
El sistema de detección y extinción se integrará al monitoreo en línea y notificará al NOC y a la contraparte del CLIENTE.\newline\textit{BT; Obligatorio.} &
Sí &
Gabinete IP55; UPS Industrial; Nodo de Borde &
SD4, sec. 4.2 \\ \hline

RT-\allowbreak{}06.\allowbreak{}20 &
El ingreso al recinto se controlará mediante seguridad física y control de acceso biométrico basado principalmente en biometría facial, con AFIS como respaldo. Se admite proponer sistemas de mayor seguridad.\newline\textit{BT; Obligatorio.} &
Sí &
Gabinete IP55; UPS Industrial; Nodo de Borde &
SD4, sec. 4.2 \\ \hline

RT-\allowbreak{}06.\allowbreak{}21 &
Todo ingreso y egreso quedará registrado en una bitácora auditable, con identificación de la persona, fecha, hora y motivo, conservada por el período de retención declarado.\newline\textit{BT; Obligatorio.} &
Sí &
Gabinete IP55; UPS Industrial; Nodo de Borde &
SD4, sec. 4.2 \\ \hline

RT-\allowbreak{}06.\allowbreak{}22 &
Entre el acceso principal y el término del pasillo de la zona de control se dispondrá un espacio para la atención de personas en proceso de enrolamiento. Se evaluará mejor la existencia de una estación de enrolamiento fuera de las instalaciones del recinto técnico.\newline\textit{BT; Obligatorio.} &
Sí &
Gabinete IP55; UPS Industrial; Nodo de Borde &
SD4, sec. 4.2 \\ \hline

RT-\allowbreak{}06.\allowbreak{}23 &
Al término del pasillo se instalará un acceso que impida el paso de más de una persona a la vez, con nueva verificación de identidad previa al ingreso.\newline\textit{BT; Obligatorio.} &
Sí &
Gabinete IP55; UPS Industrial; Nodo de Borde &
SD4, sec. 4.2 \\ \hline

RT-\allowbreak{}06.\allowbreak{}24 &
El recinto contará con videovigilancia y monitoreo IP, con imágenes en línea y disponibles para visualización de al menos los últimos 30 días. Las grabaciones anteriores se respaldarán en un medio secundario recuperable y auditable.\newline\textit{BT; Obligatorio.} &
Sí &
Gabinete IP55; UPS Industrial; Nodo de Borde &
SD4, sec. 4.2 \\ \hline

RT-\allowbreak{}06.\allowbreak{}25 &
El PROPONENTE declarará el procedimiento de acceso de terceros —fabricantes, mantenedores, auditores— con acompañamiento obligatorio y registro.\newline\textit{BT; Obligatorio.} &
Sí &
Gabinete IP55; UPS Industrial; Nodo de Borde &
SD4, sec. 4.2 \\ \hline

RT-\allowbreak{}06.\allowbreak{}26 &
Se habilitará un servicio de custodia de medios de respaldo para el sitio primario, en un medio físico transportable a otro lugar cuando el CLIENTE lo determine. Se admite proponer una solución más segura y eficiente, debidamente justificada.\newline\textit{BT; Obligatorio.} &
Sí &
Gabinete IP55; UPS Industrial; Nodo de Borde &
SD4, sec. 4.2 \\ \hline

RT-\allowbreak{}06.\allowbreak{}27 &
El recinto de custodia cumplirá exigencias de luminosidad, humedad, ventilación y cualquier otro factor que pueda afectar la calidad y la disponibilidad de los medios.\newline\textit{BT; Obligatorio.} &
Sí &
Gabinete IP55; UPS Industrial; Nodo de Borde &
SD4, sec. 4.2 \\ \hline

RT-\allowbreak{}06.\allowbreak{}28 &
Se llevará un inventario de medios con rotación, verificación periódica de legibilidad y registro de todo movimiento de entrada y de salida.\newline\textit{BT; Obligatorio.} &
Sí &
Gabinete IP55; UPS Industrial; Nodo de Borde &
SD4, sec. 4.2 \\ \hline

RT-\allowbreak{}06.\allowbreak{}29 &
El PROPONENTE habilitará el espacio físico necesario para el personal encargado de la operación y administración de la plataforma, con estaciones de trabajo, telefonía, conexión a Internet y todo elemento que permita realizar la labor en condiciones adecuadas.\newline\textit{BT; Obligatorio.} &
Sí &
Gabinete IP55; UPS Industrial; Nodo de Borde &
SD4, sec. 4.2 \\ \hline

RT-\allowbreak{}06.\allowbreak{}30 &
El espacio de operación estará separado de la sala de equipos y no requerirá el ingreso al recinto técnico para las labores habituales de operación.\newline\textit{BT; Obligatorio.} &
Sí &
Gabinete IP55; UPS Industrial; Nodo de Borde &
SD4, sec. 4.2 \\ \hline

RT-\allowbreak{}06.\allowbreak{}31 &
Las instalaciones sanitarias, las zonas de seguridad ante emergencia y las áreas exteriores existentes en el edificio del CLIENTE podrán utilizarse y no deben implementarse nuevamente. El PROPONENTE declarará de cuáles hará uso.\newline\textit{BT; Obligatorio.} &
Sí &
Gabinete IP55; UPS Industrial; Nodo de Borde &
SD4, sec. 4.2 \\ \hline

RT-\allowbreak{}06.\allowbreak{}32 &
El acceso a las redes de comunicaciones estará provisto a través de rutas físicas distintas, con ingreso al edificio por puntos separados.\newline\textit{BT; Obligatorio.} &
Sí &
Gabinete IP55; UPS Industrial; Nodo de Borde &
SD4, sec. 4.2 \\ \hline

RT-\allowbreak{}06.\allowbreak{}33 &
El PROPONENTE proveerá toda la conectividad, la seguridad y las canalizaciones, cañerías o ductos requeridos para cumplir los niveles de servicio comprometidos.\newline\textit{BT; Obligatorio.} &
Sí &
Gabinete IP55; UPS Industrial; Nodo de Borde &
SD4, sec. 4.2 \\ \hline

RT-\allowbreak{}06.\allowbreak{}34 &
Se privilegiará al PROPONENTE que ofrezca o provea especificaciones nuevas o mejores que las aquí establecidas, debidamente fundamentadas.\newline\textit{BT; Deseable.} &
No &
No aplica / Descartado: Equipamiento de cómputo en rack mural estanco sin requerimiento de supresión de incendios por gas inerte de sala blanca &
SD3, sec. 3.2.2; SD4, sec. 4.2 \\ \hline

\end{longtable}


\subsection*{BT, capítulo 7: SITE SECUNDARIO Y RECUPERACIÓN ANTE DESASTRES}

\addcontentsline{toc}{subsection}{BT, capítulo 7: SITE SECUNDARIO Y RECUPERACIÓN ANTE DESASTRES}

La Tabla~\ref{tab:t12-bt-cap-7} individualiza los requisitos técnicos de este ámbito y su localización en la propuesta.

\begin{longtable}{|L{0.09\AnchoMatrizTDoce}|L{0.385\AnchoMatrizTDoce}|L{0.065\AnchoMatrizTDoce}|L{0.18\AnchoMatrizTDoce}|L{0.28\AnchoMatrizTDoce}|}

\caption{Cumplimiento y trazabilidad: BT, capítulo 7: SITE SECUNDARIO Y RECUPERACIÓN ANTE DESASTRES}\label{tab:t12-bt-cap-7}\\

\hline

\EncabezadoTabla{ID requerimiento} & \EncabezadoTabla{Descripción} & \EncabezadoTabla{Cumple\newline(Sí/No)} & \EncabezadoTabla{Componente que lo satisface} & \EncabezadoTabla{Sección de la propuesta} \\ \hline

\endfirsthead

\hline

\EncabezadoTabla{ID requerimiento} & \EncabezadoTabla{Descripción} & \EncabezadoTabla{Cumple\newline(Sí/No)} & \EncabezadoTabla{Componente que lo satisface} & \EncabezadoTabla{Sección de la propuesta} \\ \hline

\endhead

\multicolumn{5}{r}{\fontsize{9}{11}\selectfont Continúa en la página siguiente.}\\

\endfoot

\hline

\endlastfoot


RT-\allowbreak{}07.\allowbreak{}01 &
El PROPONENTE declarará la modalidad escogida —activo-activo o activo-pasivo— y la justificará frente al costo, al RTO comprometido y a la complejidad operacional que introduce.\newline\textit{BT; Obligatorio.} &
Sí &
AWS Multi-AZ; Replicación DRP São Paulo; RTO 2 h &
SD4, sec. 4.3 \\ \hline

RT-\allowbreak{}07.\allowbreak{}02 &
El sitio secundario estará emplazado a una distancia suficiente del principal para no verse afectado por el mismo evento de fuerza mayor. El PROPONENTE declarará la distancia y el análisis de amenazas comunes considerado.\newline\textit{BT; Obligatorio.} &
Sí &
AWS Multi-AZ; Replicación DRP São Paulo; RTO 2 h &
SD4, sec. 4.3 \\ \hline

RT-\allowbreak{}07.\allowbreak{}03 &
La replicación de datos será continua, con medición y alertamiento del retraso de replicación.\newline\textit{BT; Obligatorio.} &
Sí &
AWS Multi-AZ; Replicación DRP São Paulo; RTO 2 h &
SD4, sec. 4.3 \\ \hline

RT-\allowbreak{}07.\allowbreak{}04 &
El objetivo de tiempo de recuperación (RTO) no superará 4 horas y el objetivo de punto de recuperación (RPO) no superará 15 minutos para los servicios críticos, salvo exigencia superior del caso.\newline\textit{BT; Obligatorio.} &
Sí &
AWS Multi-AZ; Replicación DRP São Paulo; RTO 2 h &
SD4, sec. 4.3 \\ \hline

RT-\allowbreak{}07.\allowbreak{}05 &
El procedimiento de conmutación estará documentado, automatizado en la mayor medida posible y ejecutable por el personal del CLIENTE tras la transferencia de conocimiento.\newline\textit{BT; Obligatorio.} &
Sí &
AWS Multi-AZ; Replicación DRP São Paulo; RTO 2 h &
SD4, sec. 4.3 \\ \hline

RT-\allowbreak{}07.\allowbreak{}06 &
Existirá un procedimiento de retorno al sitio principal igualmente documentado y probado, con reconciliación de los datos generados durante la contingencia.\newline\textit{BT; Obligatorio.} &
Sí &
AWS Multi-AZ; Replicación DRP São Paulo; RTO 2 h &
SD4, sec. 4.3 \\ \hline

RT-\allowbreak{}07.\allowbreak{}07 &
El plan de recuperación ante desastres se probará al menos dos veces al año mediante conmutación real, con informe de resultados, medición del RTO y del RPO efectivamente alcanzados y plan de corrección de las brechas detectadas.\newline\textit{BT; Obligatorio.} &
Sí &
AWS Multi-AZ; Replicación DRP São Paulo; RTO 2 h &
SD4, sec. 4.3 \\ \hline

RT-\allowbreak{}07.\allowbreak{}08 &
Se valorará que la conmutación sea automática ante la detección de indisponibilidad, con criterio de disparo declarado y protección contra conmutación innecesaria.\newline\textit{BT; Deseable.} &
Sí &
Políticas de conmutación en AWS Route 53 con chequeos de salud y umbrales de histéresis &
SD4, sec. 4.3.2 \\ \hline

RT-\allowbreak{}07.\allowbreak{}09 &
La política de respaldo seguirá el esquema 3-2-1-1-0: tres copias, en dos medios distintos, una fuera de sitio, una inmutable o fuera de línea y cero errores de verificación de restauración.\newline\textit{BT; Obligatorio.} &
Sí &
AWS Multi-AZ; Replicación DRP São Paulo; RTO 2 h &
SD4, sec. 4.3 \\ \hline

RT-\allowbreak{}07.\allowbreak{}10 &
Los respaldos estarán cifrados en reposo y en tránsito, con clave gestionada de forma independiente de la infraestructura respaldada.\newline\textit{BT; Obligatorio.} &
Sí &
AWS Multi-AZ; Replicación DRP São Paulo; RTO 2 h &
SD4, sec. 4.3 \\ \hline

RT-\allowbreak{}07.\allowbreak{}11 &
Las copias inmutables estarán protegidas contra borrado y contra modificación durante su período de retención, incluso frente a credenciales administrativas comprometidas.\newline\textit{BT; Obligatorio.} &
Sí &
AWS Multi-AZ; Replicación DRP São Paulo; RTO 2 h &
SD4, sec. 4.3 \\ \hline

RT-\allowbreak{}07.\allowbreak{}12 &
Se ejecutará y documentará una prueba de restauración al menos mensual, sobre una muestra representativa, con medición del tiempo efectivo de restauración.\newline\textit{BT; Obligatorio.} &
Sí &
AWS Multi-AZ; Replicación DRP São Paulo; RTO 2 h &
SD4, sec. 4.3 \\ \hline

RT-\allowbreak{}07.\allowbreak{}13 &
El PROPONENTE declarará, por cada dominio de datos, la frecuencia de respaldo, el período de retención y el tiempo estimado de restauración completa.\newline\textit{BT; Obligatorio.} &
Sí &
AWS Multi-AZ; Replicación DRP São Paulo; RTO 2 h &
SD4, sec. 4.3 \\ \hline

RT-\allowbreak{}07.\allowbreak{}14 &
Los respaldos permitirán la restauración granular: un registro, una tabla, un módulo o el sistema completo.\newline\textit{BT; Deseable.} &
Sí &
Restauración granular mediante Point-in-Time Recovery (PITR) y volcados lógicos por esquema &
SD4, sec. 4.2; SD5, sec. 5.2 \\ \hline

\end{longtable}


\subsection*{BT, capítulo 8: HARDWARE, PUESTOS DE TRABAJO Y EQUIPAMIENTO DE TERRENO}

\addcontentsline{toc}{subsection}{BT, capítulo 8: HARDWARE, PUESTOS DE TRABAJO Y EQUIPAMIENTO DE TERRENO}

La Tabla~\ref{tab:t12-bt-cap-8} individualiza los requisitos técnicos de este ámbito y su localización en la propuesta.

\begin{longtable}{|L{0.09\AnchoMatrizTDoce}|L{0.385\AnchoMatrizTDoce}|L{0.065\AnchoMatrizTDoce}|L{0.18\AnchoMatrizTDoce}|L{0.28\AnchoMatrizTDoce}|}

\caption{Cumplimiento y trazabilidad: BT, capítulo 8: HARDWARE, PUESTOS DE TRABAJO Y EQUIPAMIENTO DE TERRENO}\label{tab:t12-bt-cap-8}\\

\hline

\EncabezadoTabla{ID requerimiento} & \EncabezadoTabla{Descripción} & \EncabezadoTabla{Cumple\newline(Sí/No)} & \EncabezadoTabla{Componente que lo satisface} & \EncabezadoTabla{Sección de la propuesta} \\ \hline

\endfirsthead

\hline

\EncabezadoTabla{ID requerimiento} & \EncabezadoTabla{Descripción} & \EncabezadoTabla{Cumple\newline(Sí/No)} & \EncabezadoTabla{Componente que lo satisface} & \EncabezadoTabla{Sección de la propuesta} \\ \hline

\endhead

\multicolumn{5}{r}{\fontsize{9}{11}\selectfont Continúa en la página siguiente.}\\

\endfoot

\hline

\endlastfoot


RT-\allowbreak{}08.\allowbreak{}01 &
El PROPONENTE especificará el equipamiento de cómputo con su marca, modelo de referencia, procesador, memoria, almacenamiento local, interfaces y consumo, junto con el cálculo de dimensionamiento que lo sustenta.\newline\textit{BT; Obligatorio.} &
Sí &
Tablets IP68 Rugerizadas; NBR Industrial &
SD4, sec. 4.2.1; Formulario T-11 \\ \hline

RT-\allowbreak{}08.\allowbreak{}02 &
El almacenamiento será redundante, con tolerancia declarada a la falla de discos, control de errores y monitoreo predictivo de salud de los medios.\newline\textit{BT; Obligatorio.} &
Sí &
Tablets IP68 Rugerizadas; NBR Industrial &
SD4, sec. 4.2.1; Formulario T-11 \\ \hline

RT-\allowbreak{}08.\allowbreak{}03 &
Los conmutadores de núcleo, los cortafuegos y los balanceadores de carga estarán en configuración de alta disponibilidad, sin punto único de falla.\newline\textit{BT; Obligatorio.} &
Sí &
Tablets IP68 Rugerizadas; NBR Industrial &
SD4, sec. 4.2.1; Formulario T-11 \\ \hline

RT-\allowbreak{}08.\allowbreak{}04 &
Todo el equipamiento contará con fuentes de poder redundantes y conexión a circuitos eléctricos distintos.\newline\textit{BT; Obligatorio.} &
Sí &
Tablets IP68 Rugerizadas; NBR Industrial &
SD4, sec. 4.2.1; Formulario T-11 \\ \hline

RT-\allowbreak{}08.\allowbreak{}05 &
El PROPONENTE declarará el margen de crecimiento del dimensionamiento propuesto, expresado como porcentaje sobre la carga proyectada del caso, y el procedimiento de ampliación.\newline\textit{BT; Obligatorio.} &
Sí &
Tablets IP68 Rugerizadas; NBR Industrial &
SD4, sec. 4.2.1; Formulario T-11 \\ \hline

RT-\allowbreak{}08.\allowbreak{}06 &
El equipamiento será nuevo, sin uso previo, con garantía de fábrica vigente desde la recepción conforme.\newline\textit{BT; Obligatorio.} &
Sí &
Tablets IP68 Rugerizadas; NBR Industrial &
SD4, sec. 4.2.1; Formulario T-11 \\ \hline

RT-\allowbreak{}08.\allowbreak{}07 &
El PROPONENTE especificará las estaciones de trabajo requeridas para la operación y la administración de la plataforma, en la cantidad que determine el dimensionamiento del caso, con monitores duales.\newline\textit{BT; Obligatorio.} &
Sí &
Tablets IP68 Rugerizadas; NBR Industrial &
SD4, sec. 4.2.1; Formulario T-11 \\ \hline

RT-\allowbreak{}08.\allowbreak{}08 &
Los puestos de trabajo cumplirán las condiciones ergonómicas de la NCh 2527 y los equipos contarán con certificación de eficiencia energética.\newline\textit{BT; Obligatorio.} &
Sí &
Tablets IP68 Rugerizadas; NBR Industrial &
SD4, sec. 4.2.1; Formulario T-11 \\ \hline

RT-\allowbreak{}08.\allowbreak{}09 &
Las estaciones estarán gestionadas de forma centralizada, con cifrado de disco, control de dispositivos extraíbles, antivirus con detección y respuesta y actualización automatizada.\newline\textit{BT; Obligatorio.} &
Sí &
Tablets IP68 Rugerizadas; NBR Industrial &
SD4, sec. 4.2.1; Formulario T-11 \\ \hline

RT-\allowbreak{}08.\allowbreak{}10 &
El PROPONENTE especificará cada dispositivo de terreno con marca, modelo de referencia, cantidad, características mínimas, accesorios, consumibles y costo unitario estimado, aun cuando su compra sea de cargo del CLIENTE.\newline\textit{BT; Obligatorio.} &
Sí &
Tablets IP68 Rugerizadas; NBR Industrial &
SD4, sec. 4.2.1; Formulario T-11 \\ \hline

RT-\allowbreak{}08.\allowbreak{}11 &
La especificación considerará las condiciones reales de uso del caso: intemperie, humedad, polvo, vibración, temperatura, uso con guantes, luminosidad y autonomía de batería requerida por turno.\newline\textit{BT; Obligatorio.} &
Sí &
Tablets IP68 Rugerizadas; NBR Industrial &
SD4, sec. 4.2.1; Formulario T-11 \\ \hline

RT-\allowbreak{}08.\allowbreak{}12 &
Los dispositivos declararán su grado de protección contra polvo y agua y su resistencia a caídas, coherentes con el entorno de operación.\newline\textit{BT; Obligatorio.} &
Sí &
Tablets IP68 Rugerizadas; NBR Industrial &
SD4, sec. 4.2.1; Formulario T-11 \\ \hline

RT-\allowbreak{}08.\allowbreak{}13 &
El PROPONENTE indicará el ciclo de vida esperado de cada dispositivo, la disponibilidad de repuestos y el plan de reposición durante los 56 meses del Contrato.\newline\textit{BT; Obligatorio.} &
Sí &
Tablets IP68 Rugerizadas; NBR Industrial &
SD4, sec. 4.2.1; Formulario T-11 \\ \hline

RT-\allowbreak{}08.\allowbreak{}14 &
Los dispositivos se integrarán a la gestión centralizada de flota exigida en RT-03.18.\newline\textit{BT; Obligatorio.} &
Sí &
Tablets IP68 Rugerizadas; NBR Industrial &
SD4, sec. 4.2.1; Formulario T-11 \\ \hline

RT-\allowbreak{}08.\allowbreak{}15 &
El PROPONENTE proveerá una unidad de cada tipo de dispositivo especificado para pruebas de aceptación por parte del CLIENTE, antes de la compra masiva.\newline\textit{BT; Deseable.} &
No &
No aplica / Descartado: La adquisición y reemplazo de equipamiento de terreno corresponde al mandante según bases (SUP-04); Synaptix especifica y homologa &
SD3, sec. 3.2.2; SUP-04 \\ \hline

RT-\allowbreak{}08.\allowbreak{}16 &
El PROPONENTE presentará el plan de ciclo de vida del equipamiento: recepción, puesta en servicio, mantención, actualización, retiro y disposición final.\newline\textit{BT; Obligatorio.} &
Sí &
Tablets IP68 Rugerizadas; NBR Industrial &
SD4, sec. 4.2.1; Formulario T-11 \\ \hline

RT-\allowbreak{}08.\allowbreak{}17 &
Todo medio de almacenamiento que salga de servicio será borrado de forma segura y verificable, con certificado de destrucción o de sanitización entregado al CLIENTE.\newline\textit{BT; Obligatorio.} &
Sí &
Tablets IP68 Rugerizadas; NBR Industrial &
SD4, sec. 4.2.1; Formulario T-11 \\ \hline

RT-\allowbreak{}08.\allowbreak{}18 &
La disposición final de equipamiento electrónico se realizará con gestor autorizado, conforme a la normativa de residuos aplicable, con certificado de disposición.\newline\textit{BT; Obligatorio.} &
Sí &
Tablets IP68 Rugerizadas; NBR Industrial &
SD4, sec. 4.2.1; Formulario T-11 \\ \hline

RT-\allowbreak{}08.\allowbreak{}19 &
Se valorará una estrategia de reacondicionamiento o de extensión de vida útil que reduzca el impacto ambiental, cuantificada en la propuesta.\newline\textit{BT; Deseable.} &
No &
No aplica / Descartado: Todo el hardware homologado provisto por el mandante es 100\% nuevo de fábrica con garantía de 36 meses &
SD3, sec. 3.2.2; Formulario T-11 \\ \hline

\end{longtable}


\subsection*{BT, capítulo 9: DESEMPEÑO, CAPACIDAD Y ESCALABILIDAD}

\addcontentsline{toc}{subsection}{BT, capítulo 9: DESEMPEÑO, CAPACIDAD Y ESCALABILIDAD}

La Tabla~\ref{tab:t12-bt-cap-9} individualiza los requisitos técnicos de este ámbito y su localización en la propuesta.

\begin{longtable}{|L{0.09\AnchoMatrizTDoce}|L{0.385\AnchoMatrizTDoce}|L{0.065\AnchoMatrizTDoce}|L{0.18\AnchoMatrizTDoce}|L{0.28\AnchoMatrizTDoce}|}

\caption{Cumplimiento y trazabilidad: BT, capítulo 9: DESEMPEÑO, CAPACIDAD Y ESCALABILIDAD}\label{tab:t12-bt-cap-9}\\

\hline

\EncabezadoTabla{ID requerimiento} & \EncabezadoTabla{Descripción} & \EncabezadoTabla{Cumple\newline(Sí/No)} & \EncabezadoTabla{Componente que lo satisface} & \EncabezadoTabla{Sección de la propuesta} \\ \hline

\endfirsthead

\hline

\EncabezadoTabla{ID requerimiento} & \EncabezadoTabla{Descripción} & \EncabezadoTabla{Cumple\newline(Sí/No)} & \EncabezadoTabla{Componente que lo satisface} & \EncabezadoTabla{Sección de la propuesta} \\ \hline

\endhead

\multicolumn{5}{r}{\fontsize{9}{11}\selectfont Continúa en la página siguiente.}\\

\endfoot

\hline

\endlastfoot


RT-\allowbreak{}09.\allowbreak{}01 &
El PROPONENTE presentará el cálculo de capacidad que sustenta su dimensionamiento, con los supuestos de usuarios concurrentes, transacciones por segundo, volumen de datos y crecimiento anual, tomados de la volumetría del caso.\newline\textit{BT; Obligatorio.} &
Sí &
Suite k6/JMeter; JaCoCo; Pruebas de Carga &
SD6, sec. 6.2; SD9, sec. 9.1 \\ \hline

RT-\allowbreak{}09.\allowbreak{}02 &
La solución soportará la concurrencia y el volumen de transacciones que fije el caso, y mantendrá los umbrales del numeral 9.1 bajo esa carga.\newline\textit{BT; Según caso.} &
Sí &
Suite k6/JMeter; JaCoCo; Pruebas de Carga &
SD6, sec. 6.2; SD9, sec. 9.1 \\ \hline

RT-\allowbreak{}09.\allowbreak{}03 &
La solución soportará, sin rediseño, un crecimiento de al menos tres veces la volumetría inicial del caso en un horizonte de tres años.\newline\textit{BT; Obligatorio.} &
Sí &
Suite k6/JMeter; JaCoCo; Pruebas de Carga &
SD6, sec. 6.2; SD9, sec. 9.1 \\ \hline

RT-\allowbreak{}09.\allowbreak{}04 &
El escalamiento de las capas de aplicación e integración será horizontal y automático, con tiempo de reacción declarado y sin pérdida de transacciones en curso.\newline\textit{BT; Obligatorio.} &
Sí &
Suite k6/JMeter; JaCoCo; Pruebas de Carga &
SD6, sec. 6.2; SD9, sec. 9.1 \\ \hline

RT-\allowbreak{}09.\allowbreak{}05 &
El PROPONENTE identificará el componente que primero se convertirá en cuello de botella al crecer la carga y explicará cómo lo detectará y cómo lo resolverá.\newline\textit{BT; Obligatorio.} &
Sí &
Suite k6/JMeter; JaCoCo; Pruebas de Carga &
SD6, sec. 6.2; SD9, sec. 9.1 \\ \hline

RT-\allowbreak{}09.\allowbreak{}06 &
Se ejecutarán pruebas de carga sobre Preproducción con un volumen equivalente a 1,5 veces el peak declarado, y pruebas de estrés hasta identificar el punto de quiebre de la solución.\newline\textit{BT; Obligatorio.} &
Sí &
Suite k6/JMeter; JaCoCo; Pruebas de Carga &
SD6, sec. 6.2; SD9, sec. 9.1 \\ \hline

RT-\allowbreak{}09.\allowbreak{}07 &
El informe de pruebas de carga incluirá la curva de tiempo de respuesta frente a carga, el punto de saturación, el consumo de recursos y el comportamiento durante y después del peak.\newline\textit{BT; Obligatorio.} &
Sí &
Suite k6/JMeter; JaCoCo; Pruebas de Carga &
SD6, sec. 6.2; SD9, sec. 9.1 \\ \hline

RT-\allowbreak{}09.\allowbreak{}08 &
La solución degradará de forma controlada al superarse la capacidad: encolamiento, limitación de tasa y mensaje explícito a la persona usuaria, nunca error genérico ni pérdida silenciosa de transacciones.\newline\textit{BT; Obligatorio.} &
Sí &
Suite k6/JMeter; JaCoCo; Pruebas de Carga &
SD6, sec. 6.2; SD9, sec. 9.1 \\ \hline

RT-\allowbreak{}09.\allowbreak{}09 &
Se gestionará la capacidad durante la Operación con proyección trimestral de crecimiento, alertas anticipadas de agotamiento y propuesta de ajuste de dimensionamiento y de costo.\newline\textit{BT; Obligatorio.} &
Sí &
Suite k6/JMeter; JaCoCo; Pruebas de Carga &
SD6, sec. 6.2; SD9, sec. 9.1 \\ \hline

RT-\allowbreak{}09.\allowbreak{}10 &
Se valorará la existencia de pruebas de carga automatizadas ejecutadas de forma periódica en el flujo de integración continua, con detección de regresiones de desempeño.\newline\textit{BT; Deseable.} &
No &
No aplica / Descartado: Pruebas de carga se ejecutan como hito de validación en Quality Gates de Preproducción; no integradas en cada commit de CI &
SD3, sec. 3.2.2; SD6, sec. 6.2 \\ \hline

\end{longtable}


\subsection*{BT, capítulo 10: DISPONIBILIDAD, CONTINUIDAD Y RESILIENCIA}

\addcontentsline{toc}{subsection}{BT, capítulo 10: DISPONIBILIDAD, CONTINUIDAD Y RESILIENCIA}

La Tabla~\ref{tab:t12-bt-cap-10} individualiza los requisitos técnicos de este ámbito y su localización en la propuesta.

\begin{longtable}{|L{0.09\AnchoMatrizTDoce}|L{0.385\AnchoMatrizTDoce}|L{0.065\AnchoMatrizTDoce}|L{0.18\AnchoMatrizTDoce}|L{0.28\AnchoMatrizTDoce}|}

\caption{Cumplimiento y trazabilidad: BT, capítulo 10: DISPONIBILIDAD, CONTINUIDAD Y RESILIENCIA}\label{tab:t12-bt-cap-10}\\

\hline

\EncabezadoTabla{ID requerimiento} & \EncabezadoTabla{Descripción} & \EncabezadoTabla{Cumple\newline(Sí/No)} & \EncabezadoTabla{Componente que lo satisface} & \EncabezadoTabla{Sección de la propuesta} \\ \hline

\endfirsthead

\hline

\EncabezadoTabla{ID requerimiento} & \EncabezadoTabla{Descripción} & \EncabezadoTabla{Cumple\newline(Sí/No)} & \EncabezadoTabla{Componente que lo satisface} & \EncabezadoTabla{Sección de la propuesta} \\ \hline

\endhead

\multicolumn{5}{r}{\fontsize{9}{11}\selectfont Continúa en la página siguiente.}\\

\endfoot

\hline

\endlastfoot


RT-\allowbreak{}10.\allowbreak{}01 &
La solución alcanzará una disponibilidad mensual mínima de 99,9\% para los servicios clasificados como críticos, medida sobre la transacción de negocio de extremo a extremo.\newline\textit{BT; Obligatorio.} &
Sí &
SLO 99.9\%; Monitoreo Sintético y CloudWatch &
SD4, sec. 4.3; SD9, sec. 9.1 \\ \hline

RT-\allowbreak{}10.\allowbreak{}02 &
El PROPONENTE clasificará cada servicio de la solución en crítico, alto, medio o bajo, justificando la clasificación con el impacto operacional de su indisponibilidad, y aplicará a cada uno el nivel de servicio correspondiente del Artículo 78° de las Bases Administrativas.\newline\textit{BT; Obligatorio.} &
Sí &
SLO 99.9\%; Monitoreo Sintético y CloudWatch &
SD4, sec. 4.3; SD9, sec. 9.1 \\ \hline

RT-\allowbreak{}10.\allowbreak{}03 &
El plan de continuidad del negocio se elaborará conforme a ISO 22301, con análisis de impacto en el negocio, escenarios de contingencia, procedimientos manuales de respaldo y criterios de activación.\newline\textit{BT; Obligatorio.} &
Sí &
SLO 99.9\%; Monitoreo Sintético y CloudWatch &
SD4, sec. 4.3; SD9, sec. 9.1 \\ \hline

RT-\allowbreak{}10.\allowbreak{}04 &
La continuidad TIC se estructurará conforme a ISO/IEC 27031, articulada con el plan de recuperación ante desastres del Capítulo 7.\newline\textit{BT; Obligatorio.} &
Sí &
SLO 99.9\%; Monitoreo Sintético y CloudWatch &
SD4, sec. 4.3; SD9, sec. 9.1 \\ \hline

RT-\allowbreak{}10.\allowbreak{}05 &
Los mantenimientos programados se ejecutarán fuera de la ventana operacional crítica que defina el caso, con aviso previo mínimo de diez días hábiles.\newline\textit{BT; Según caso.} &
Sí &
SLO 99.9\%; Monitoreo Sintético y CloudWatch &
SD4, sec. 4.3; SD9, sec. 9.1 \\ \hline

RT-\allowbreak{}10.\allowbreak{}06 &
La solución permitirá desplegar cambios sin interrupción del servicio. Las ventanas de indisponibilidad programada serán excepcionales y deberán justificarse caso a caso.\newline\textit{BT; Obligatorio.} &
Sí &
SLO 99.9\%; Monitoreo Sintético y CloudWatch &
SD4, sec. 4.3; SD9, sec. 9.1 \\ \hline

RT-\allowbreak{}10.\allowbreak{}07 &
Se ejecutarán pruebas de resiliencia mediante inyección controlada de fallas —caída de instancia, de zona, de dependencia externa, latencia elevada, saturación de disco— antes de cada paso a producción y al menos una vez por semestre durante la Operación.\newline\textit{BT; Obligatorio.} &
Sí &
SLO 99.9\%; Monitoreo Sintético y CloudWatch &
SD4, sec. 4.3; SD9, sec. 9.1 \\ \hline

RT-\allowbreak{}10.\allowbreak{}08 &
El PROPONENTE documentará, por cada dependencia externa, el comportamiento de la solución cuando esa dependencia no responde, responde con error o responde con lentitud.\newline\textit{BT; Obligatorio.} &
Sí &
SLO 99.9\%; Monitoreo Sintético y CloudWatch &
SD4, sec. 4.3; SD9, sec. 9.1 \\ \hline

RT-\allowbreak{}10.\allowbreak{}09 &
Se declarará un presupuesto de error por servicio crítico y su vinculación con el ritmo de despliegue de cambios.\newline\textit{BT; Deseable.} &
Sí &
Presupuesto de error fijado en 0,1\% mensual (SLO 99,9\%); congelamiento automático ante agotamiento &
SD6, sec. 6.2; SD11, sec. 11.2 \\ \hline

\end{longtable}


\subsection*{BT, capítulo 11: SEGURIDAD DE LA INFORMACIÓN Y CIBERSEGURIDAD}

\addcontentsline{toc}{subsection}{BT, capítulo 11: SEGURIDAD DE LA INFORMACIÓN Y CIBERSEGURIDAD}

La Tabla~\ref{tab:t12-bt-cap-11} individualiza los requisitos técnicos de este ámbito y su localización en la propuesta.

\begin{longtable}{|L{0.09\AnchoMatrizTDoce}|L{0.385\AnchoMatrizTDoce}|L{0.065\AnchoMatrizTDoce}|L{0.18\AnchoMatrizTDoce}|L{0.28\AnchoMatrizTDoce}|}

\caption{Cumplimiento y trazabilidad: BT, capítulo 11: SEGURIDAD DE LA INFORMACIÓN Y CIBERSEGURIDAD}\label{tab:t12-bt-cap-11}\\

\hline

\EncabezadoTabla{ID requerimiento} & \EncabezadoTabla{Descripción} & \EncabezadoTabla{Cumple\newline(Sí/No)} & \EncabezadoTabla{Componente que lo satisface} & \EncabezadoTabla{Sección de la propuesta} \\ \hline

\endfirsthead

\hline

\EncabezadoTabla{ID requerimiento} & \EncabezadoTabla{Descripción} & \EncabezadoTabla{Cumple\newline(Sí/No)} & \EncabezadoTabla{Componente que lo satisface} & \EncabezadoTabla{Sección de la propuesta} \\ \hline

\endhead

\multicolumn{5}{r}{\fontsize{9}{11}\selectfont Continúa en la página siguiente.}\\

\endfoot

\hline

\endlastfoot


RT-\allowbreak{}11.\allowbreak{}01 &
La arquitectura de seguridad se basará en el modelo Zero Trust conforme a NIST SP 800-207: verificación explícita de cada solicitud, privilegio mínimo y presunción de compromiso.\newline\textit{BT; Obligatorio.} &
Sí &
Zero Trust; AWS KMS AES-GCM-256; AWS WAF v2 &
SD4, sec. 4.1.3; SD6, sec. 6.1 \\ \hline

RT-\allowbreak{}11.\allowbreak{}02 &
El PROPONENTE entregará un modelado de amenazas documentado por cada componente y por cada integración externa, con metodología declarada (STRIDE u otra), y lo actualizará ante cada cambio arquitectónico relevante.\newline\textit{BT; Obligatorio.} &
Sí &
Zero Trust; AWS KMS AES-GCM-256; AWS WAF v2 &
SD4, sec. 4.1.3; SD6, sec. 6.1 \\ \hline

RT-\allowbreak{}11.\allowbreak{}03 &
La información del CLIENTE se clasificará por nivel de sensibilidad, con controles diferenciados por nivel documentados en una matriz.\newline\textit{BT; Obligatorio.} &
Sí &
Zero Trust; AWS KMS AES-GCM-256; AWS WAF v2 &
SD4, sec. 4.1.3; SD6, sec. 6.1 \\ \hline

RT-\allowbreak{}11.\allowbreak{}04 &
Existirá un programa de gestión de vulnerabilidades con escaneo continuo y plazos máximos de remediación de 7 días corridos para vulnerabilidades críticas, 15 días para altas y 30 días para medias, contados desde su publicación o detección.\newline\textit{BT; Obligatorio.} &
Sí &
Zero Trust; AWS KMS AES-GCM-256; AWS WAF v2 &
SD4, sec. 4.1.3; SD6, sec. 6.1 \\ \hline

RT-\allowbreak{}11.\allowbreak{}05 &
El PROPONENTE mantendrá una matriz de controles de seguridad trazable a ISO/IEC 27001 e ISO/IEC 27002, indicando el control, su implementación concreta en la solución y la evidencia que lo acredita.\newline\textit{BT; Obligatorio.} &
Sí &
Zero Trust; AWS KMS AES-GCM-256; AWS WAF v2 &
SD4, sec. 4.1.3; SD6, sec. 6.1 \\ \hline

RT-\allowbreak{}11.\allowbreak{}06 &
Se aplicarán los controles de ISO/IEC 27017 para los servicios en nube y de ISO/IEC 27018 para el tratamiento de datos personales en nube.\newline\textit{BT; Obligatorio.} &
Sí &
Zero Trust; AWS KMS AES-GCM-256; AWS WAF v2 &
SD4, sec. 4.1.3; SD6, sec. 6.1 \\ \hline

RT-\allowbreak{}11.\allowbreak{}07 &
La publicación de servicios se realizará exclusivamente a través de la capa de borde, con red de distribución de contenidos, cortafuegos de aplicaciones web con reglas gestionadas y personalizadas, y protección contra denegación de servicio distribuida en capas 3, 4 y 7.\newline\textit{BT; Obligatorio.} &
Sí &
Zero Trust; AWS KMS AES-GCM-256; AWS WAF v2 &
SD4, sec. 4.1.3; SD6, sec. 6.1 \\ \hline

RT-\allowbreak{}11.\allowbreak{}08 &
El cifrado en tránsito empleará TLS 1.3, con prohibición expresa de TLS 1.0 y 1.1, conjuntos de cifrado modernos, HSTS con precarga y gestión automatizada de certificados con rotación y alerta anticipada de vencimiento.\newline\textit{BT; Obligatorio.} &
Sí &
Zero Trust; AWS KMS AES-GCM-256; AWS WAF v2 &
SD4, sec. 4.1.3; SD6, sec. 6.1 \\ \hline

RT-\allowbreak{}11.\allowbreak{}09 &
La totalidad de los datos en reposo estará cifrada, con claves gestionadas en un servicio de gestión de claves o en un módulo de seguridad de hardware, política de rotación declarada y separación de funciones en la custodia de claves.\newline\textit{BT; Obligatorio.} &
Sí &
Zero Trust; AWS KMS AES-GCM-256; AWS WAF v2 &
SD4, sec. 4.1.3; SD6, sec. 6.1 \\ \hline

RT-\allowbreak{}11.\allowbreak{}10 &
Los datos de categoría sensible que el caso identifique se cifrarán adicionalmente a nivel de campo, de modo que el acceso a la base de datos no revele su contenido.\newline\textit{BT; Según caso.} &
Sí &
Zero Trust; AWS KMS AES-GCM-256; AWS WAF v2 &
SD4, sec. 4.1.3; SD6, sec. 6.1 \\ \hline

RT-\allowbreak{}11.\allowbreak{}11 &
La puerta de enlace de servicios aplicará autenticación, autorización, cuotas, límites de tasa, validación de esquema e inspección de carga útil.\newline\textit{BT; Obligatorio.} &
Sí &
Zero Trust; AWS KMS AES-GCM-256; AWS WAF v2 &
SD4, sec. 4.1.3; SD6, sec. 6.1 \\ \hline

RT-\allowbreak{}11.\allowbreak{}12 &
Los puntos de entrada públicos contarán con protección contra bots y abuso automatizado, con reto progresivo que no degrade la accesibilidad ni bloquee a personas usuarias legítimas.\newline\textit{BT; Obligatorio.} &
Sí &
Zero Trust; AWS KMS AES-GCM-256; AWS WAF v2 &
SD4, sec. 4.1.3; SD6, sec. 6.1 \\ \hline

RT-\allowbreak{}11.\allowbreak{}13 &
El PROPONENTE declarará la superficie de exposición completa de la solución: cada nombre de dominio, puerto y servicio alcanzable desde fuera de la red del CLIENTE.\newline\textit{BT; Obligatorio.} &
Sí &
Zero Trust; AWS KMS AES-GCM-256; AWS WAF v2 &
SD4, sec. 4.1.3; SD6, sec. 6.1 \\ \hline

RT-\allowbreak{}11.\allowbreak{}14 &
Los eventos de seguridad se registrarán de forma centralizada e inalterable, con retención mínima de doce meses en línea y veinticuatro meses adicionales en archivo recuperable.\newline\textit{BT; Obligatorio.} &
Sí &
Zero Trust; AWS KMS AES-GCM-256; AWS WAF v2 &
SD4, sec. 4.1.3; SD6, sec. 6.1 \\ \hline

RT-\allowbreak{}11.\allowbreak{}15 &
Los eventos se correlacionarán en una plataforma SIEM, con casos de uso de detección definidos específicamente para el proceso de negocio del caso, y no sólo genéricos de infraestructura.\newline\textit{BT; Obligatorio.} &
Sí &
Zero Trust; AWS KMS AES-GCM-256; AWS WAF v2 &
SD4, sec. 4.1.3; SD6, sec. 6.1 \\ \hline

RT-\allowbreak{}11.\allowbreak{}16 &
Se implementará detección y respuesta en puntos finales y en cargas de trabajo, tanto en nube como on-premise.\newline\textit{BT; Obligatorio.} &
Sí &
Zero Trust; AWS KMS AES-GCM-256; AWS WAF v2 &
SD4, sec. 4.1.3; SD6, sec. 6.1 \\ \hline

RT-\allowbreak{}11.\allowbreak{}17 &
El PROPONENTE dispondrá de un centro de operaciones de seguridad con cobertura 24×7, propio o subcontratado, y declarará su ubicación, dotación y procedimientos.\newline\textit{BT; Obligatorio.} &
Sí &
Zero Trust; AWS KMS AES-GCM-256; AWS WAF v2 &
SD4, sec. 4.1.3; SD6, sec. 6.1 \\ \hline

RT-\allowbreak{}11.\allowbreak{}18 &
El plan de respuesta a incidentes de seguridad definirá clasificación, cadena de escalamiento, plazos, responsables y protocolo de comunicación al CLIENTE dentro de las dos horas de detectado un incidente de severidad crítica.\newline\textit{BT; Obligatorio.} &
Sí &
Zero Trust; AWS KMS AES-GCM-256; AWS WAF v2 &
SD4, sec. 4.1.3; SD6, sec. 6.1 \\ \hline

RT-\allowbreak{}11.\allowbreak{}19 &
Toda brecha de seguridad o de datos personales se notificará al CLIENTE dentro de las 24 horas de su detección, con informe preliminar, y el análisis de causa raíz se entregará dentro de los cinco días hábiles siguientes.\newline\textit{BT; Obligatorio.} &
Sí &
Zero Trust; AWS KMS AES-GCM-256; AWS WAF v2 &
SD4, sec. 4.1.3; SD6, sec. 6.1 \\ \hline

RT-\allowbreak{}11.\allowbreak{}20 &
Se ejecutarán pruebas de intrusión por un tercero independiente del ADJUDICATARIO, anualmente y antes de cada paso a producción, con entrega íntegra del informe al CLIENTE y plan de remediación con plazos.\newline\textit{BT; Obligatorio.} &
Sí &
Zero Trust; AWS KMS AES-GCM-256; AWS WAF v2 &
SD4, sec. 4.1.3; SD6, sec. 6.1 \\ \hline

RT-\allowbreak{}11.\allowbreak{}21 &
Se realizarán ejercicios de simulación de incidente con participación del CLIENTE al menos una vez al año durante la Operación.\newline\textit{BT; Deseable.} &
No &
No aplica / Descartado: Plan de respuesta a incidentes bajo servicio gestionado Synaptix; simulacros con comité de crisis externo exceden dotación del club &
SD3, sec. 3.2.2; DEC-22 \\ \hline

RT-\allowbreak{}11.\allowbreak{}22 &
El flujo de integración continua incorporará análisis estático de código, análisis de composición de software, análisis dinámico y escaneo de imágenes de contenedor, con criterios de bloqueo automático del despliegue ante hallazgos críticos.\newline\textit{BT; Obligatorio.} &
Sí &
Zero Trust; AWS KMS AES-GCM-256; AWS WAF v2 &
SD4, sec. 4.1.3; SD6, sec. 6.1 \\ \hline

RT-\allowbreak{}11.\allowbreak{}23 &
Cada versión liberada se acompañará de su inventario de componentes de software en formato CycloneDX o SPDX, entregado al CLIENTE.\newline\textit{BT; Obligatorio.} &
Sí &
Zero Trust; AWS KMS AES-GCM-256; AWS WAF v2 &
SD4, sec. 4.1.3; SD6, sec. 6.1 \\ \hline

RT-\allowbreak{}11.\allowbreak{}24 &
Los artefactos se firmarán y su procedencia se verificará conforme a SLSA nivel 3 o superior.\newline\textit{BT; Obligatorio.} &
Sí &
Zero Trust; AWS KMS AES-GCM-256; AWS WAF v2 &
SD4, sec. 4.1.3; SD6, sec. 6.1 \\ \hline

RT-\allowbreak{}11.\allowbreak{}25 &
Queda prohibido el uso de datos productivos reales en ambientes no productivos sin anonimización o seudonimización verificable.\newline\textit{BT; Obligatorio.} &
Sí &
Zero Trust; AWS KMS AES-GCM-256; AWS WAF v2 &
SD4, sec. 4.1.3; SD6, sec. 6.1 \\ \hline

RT-\allowbreak{}11.\allowbreak{}26 &
El PROPONENTE declarará el proceso de aprobación de nuevas dependencias de terceros, incluyendo criterios de licencia, mantención activa y ausencia de vulnerabilidades conocidas.\newline\textit{BT; Obligatorio.} &
Sí &
Zero Trust; AWS KMS AES-GCM-256; AWS WAF v2 &
SD4, sec. 4.1.3; SD6, sec. 6.1 \\ \hline

RT-\allowbreak{}11.\allowbreak{}27 &
Las personas desarrolladoras no tendrán acceso interactivo directo al ambiente de producción. Todo acceso excepcional será temporal, aprobado, registrado y con sesión grabada.\newline\textit{BT; Obligatorio.} &
Sí &
Zero Trust; AWS KMS AES-GCM-256; AWS WAF v2 &
SD4, sec. 4.1.3; SD6, sec. 6.1 \\ \hline

RT-\allowbreak{}11.\allowbreak{}28 &
Se aplicará el marco OWASP SAMM o equivalente para medir y mejorar la madurez del proceso de desarrollo seguro, con evaluación inicial y reevaluación anual.\newline\textit{BT; Deseable.} &
Sí &
Evaluación OWASP SAMM inicial en mes 3 y reevaluación anual comprometida para el gobierno del ciclo seguro &
SD6, sec. 6.2 \\ \hline

\end{longtable}


\subsection*{BT, capítulo 12: IDENTIDAD, ACCESO Y GESTIÓN DE SESIONES}

\addcontentsline{toc}{subsection}{BT, capítulo 12: IDENTIDAD, ACCESO Y GESTIÓN DE SESIONES}

La Tabla~\ref{tab:t12-bt-cap-12} individualiza los requisitos técnicos de este ámbito y su localización en la propuesta.

\begin{longtable}{|L{0.09\AnchoMatrizTDoce}|L{0.385\AnchoMatrizTDoce}|L{0.065\AnchoMatrizTDoce}|L{0.18\AnchoMatrizTDoce}|L{0.28\AnchoMatrizTDoce}|}

\caption{Cumplimiento y trazabilidad: BT, capítulo 12: IDENTIDAD, ACCESO Y GESTIÓN DE SESIONES}\label{tab:t12-bt-cap-12}\\

\hline

\EncabezadoTabla{ID requerimiento} & \EncabezadoTabla{Descripción} & \EncabezadoTabla{Cumple\newline(Sí/No)} & \EncabezadoTabla{Componente que lo satisface} & \EncabezadoTabla{Sección de la propuesta} \\ \hline

\endfirsthead

\hline

\EncabezadoTabla{ID requerimiento} & \EncabezadoTabla{Descripción} & \EncabezadoTabla{Cumple\newline(Sí/No)} & \EncabezadoTabla{Componente que lo satisface} & \EncabezadoTabla{Sección de la propuesta} \\ \hline

\endhead

\multicolumn{5}{r}{\fontsize{9}{11}\selectfont Continúa en la página siguiente.}\\

\endfoot

\hline

\endlastfoot


RT-\allowbreak{}12.\allowbreak{}01 &
La gestión de identidad será centralizada, con federación mediante OpenID Connect y OAuth 2.1, o SAML 2.0 cuando la integración con el CLIENTE lo requiera, e integración con el directorio corporativo del CLIENTE por LDAP o su equivalente en la nube.\newline\textit{BT; Obligatorio.} &
Sí &
AWS Cognito; OIDC; OAuth 2.1 con PKCE; RBAC &
SD4, sec. 4.1.3 \\ \hline

RT-\allowbreak{}12.\allowbreak{}02 &
Existirá inicio de sesión único para todos los módulos de la solución, con cierre de sesión propagado a todos ellos.\newline\textit{BT; Obligatorio.} &
Sí &
AWS Cognito; OIDC; OAuth 2.1 con PKCE; RBAC &
SD4, sec. 4.1.3 \\ \hline

RT-\allowbreak{}12.\allowbreak{}03 &
La autenticación multifactor será obligatoria para personas usuarias administradoras, para todo acceso privilegiado y para todo acceso originado fuera de la red corporativa.\newline\textit{BT; Obligatorio.} &
Sí &
AWS Cognito; OIDC; OAuth 2.1 con PKCE; RBAC &
SD4, sec. 4.1.3 \\ \hline

RT-\allowbreak{}12.\allowbreak{}04 &
Se soportarán factores resistentes a la suplantación de identidad, tipo FIDO2 o claves de acceso, al menos para los perfiles administradores.\newline\textit{BT; Deseable.} &
Sí &
Soporte nativo de Passkeys (WebAuthn/FIDO2) en AWS Cognito para perfiles de administración y torre &
SD4, sec. 4.1.3 \\ \hline

RT-\allowbreak{}12.\allowbreak{}05 &
El control de acceso será basado en roles, complementado con control basado en atributos donde el proceso lo exija, con matriz de segregación de funciones documentada y verificable.\newline\textit{BT; Obligatorio.} &
Sí &
AWS Cognito; OIDC; OAuth 2.1 con PKCE; RBAC &
SD4, sec. 4.1.3 \\ \hline

RT-\allowbreak{}12.\allowbreak{}06 &
Los accesos privilegiados se gestionarán con elevación temporal a demanda, aprobación previa y grabación de sesión para las operaciones de mayor riesgo.\newline\textit{BT; Obligatorio.} &
Sí &
AWS Cognito; OIDC; OAuth 2.1 con PKCE; RBAC &
SD4, sec. 4.1.3 \\ \hline

RT-\allowbreak{}12.\allowbreak{}07 &
La política de sesión declarará duración máxima, caducidad por inactividad, renovación de la credencial de sesión tras la autenticación, revocación inmediata y control de sesiones concurrentes.\newline\textit{BT; Obligatorio.} &
Sí &
AWS Cognito; OIDC; OAuth 2.1 con PKCE; RBAC &
SD4, sec. 4.1.3 \\ \hline

RT-\allowbreak{}12.\allowbreak{}08 &
Las credenciales de sesión serán firmadas y de vida breve, con credencial de refresco rotatoria. Queda prohibido transportar identificadores de sesión en la ruta de la dirección web.\newline\textit{BT; Obligatorio.} &
Sí &
AWS Cognito; OIDC; OAuth 2.1 con PKCE; RBAC &
SD4, sec. 4.1.3 \\ \hline

RT-\allowbreak{}12.\allowbreak{}09 &
Se registrará auditoría completa del ciclo de vida de la identidad: creación, modificación, elevación, bloqueo y baja de cuentas, con no repudio y retención declarada.\newline\textit{BT; Obligatorio.} &
Sí &
AWS Cognito; OIDC; OAuth 2.1 con PKCE; RBAC &
SD4, sec. 4.1.3 \\ \hline

RT-\allowbreak{}12.\allowbreak{}10 &
El aprovisionamiento y el desaprovisionamiento estarán automatizados y ligados al ciclo de vida laboral, con baja efectiva en un plazo no superior a 24 horas desde la desvinculación.\newline\textit{BT; Obligatorio.} &
Sí &
AWS Cognito; OIDC; OAuth 2.1 con PKCE; RBAC &
SD4, sec. 4.1.3 \\ \hline

RT-\allowbreak{}12.\allowbreak{}11 &
El mecanismo de autenticación se adecuará al perfil operacional real descrito en el caso: entornos de terreno, uso con guantes, baja alfabetización digital, dispositivos compartidos por turno y ausencia de correo electrónico personal.\newline\textit{BT; Según caso.} &
Sí &
AWS Cognito; OIDC; OAuth 2.1 con PKCE; RBAC &
SD4, sec. 4.1.3 \\ \hline

RT-\allowbreak{}12.\allowbreak{}12 &
Las personas usuarias externas del CLIENTE —clientes, proveedores, pacientes, productores, según el caso— dispondrán de un mecanismo de registro, verificación de identidad y recuperación de acceso autoservido y seguro.\newline\textit{BT; Según caso.} &
Sí &
AWS Cognito; OIDC; OAuth 2.1 con PKCE; RBAC &
SD4, sec. 4.1.3 \\ \hline

RT-\allowbreak{}12.\allowbreak{}13 &
El PROPONENTE declarará el procedimiento de acceso de emergencia (cuenta de último recurso), su custodia, su control y su auditoría.\newline\textit{BT; Obligatorio.} &
Sí &
AWS Cognito; OIDC; OAuth 2.1 con PKCE; RBAC &
SD4, sec. 4.1.3 \\ \hline

\end{longtable}


\subsection*{BT, capítulo 13: USABILIDAD, ACCESIBILIDAD Y EXPERIENCIA DE USUARIO}

\addcontentsline{toc}{subsection}{BT, capítulo 13: USABILIDAD, ACCESIBILIDAD Y EXPERIENCIA DE USUARIO}

La Tabla~\ref{tab:t12-bt-cap-13} individualiza los requisitos técnicos de este ámbito y su localización en la propuesta.

\begin{longtable}{|L{0.09\AnchoMatrizTDoce}|L{0.385\AnchoMatrizTDoce}|L{0.065\AnchoMatrizTDoce}|L{0.18\AnchoMatrizTDoce}|L{0.28\AnchoMatrizTDoce}|}

\caption{Cumplimiento y trazabilidad: BT, capítulo 13: USABILIDAD, ACCESIBILIDAD Y EXPERIENCIA DE USUARIO}\label{tab:t12-bt-cap-13}\\

\hline

\EncabezadoTabla{ID requerimiento} & \EncabezadoTabla{Descripción} & \EncabezadoTabla{Cumple\newline(Sí/No)} & \EncabezadoTabla{Componente que lo satisface} & \EncabezadoTabla{Sección de la propuesta} \\ \hline

\endfirsthead

\hline

\EncabezadoTabla{ID requerimiento} & \EncabezadoTabla{Descripción} & \EncabezadoTabla{Cumple\newline(Sí/No)} & \EncabezadoTabla{Componente que lo satisface} & \EncabezadoTabla{Sección de la propuesta} \\ \hline

\endhead

\multicolumn{5}{r}{\fontsize{9}{11}\selectfont Continúa en la página siguiente.}\\

\endfoot

\hline

\endlastfoot


RT-\allowbreak{}13.\allowbreak{}01 &
Todas las interfaces destinadas a personas usuarias cumplirán WCAG 2.2 nivel AA, verificado con herramientas automatizadas y con pruebas manuales, con informe de conformidad entregable.\newline\textit{BT; Obligatorio.} &
Sí &
Panitao Design System; WCAG 2.2 AA; Modo Cubierta &
SD4, sec. 4.1; SD9, sec. 9.2 \\ \hline

RT-\allowbreak{}13.\allowbreak{}02 &
El diseño será responsivo y funcionará correctamente en escritorio, tableta y teléfono, con puntos de quiebre coherentes y adaptación inteligente del contenido, no simple reducción.\newline\textit{BT; Obligatorio.} &
Sí &
Panitao Design System; WCAG 2.2 AA; Modo Cubierta &
SD4, sec. 4.1; SD9, sec. 9.2 \\ \hline

RT-\allowbreak{}13.\allowbreak{}03 &
El PROPONENTE ejecutará investigación con personas usuarias reales del CLIENTE, prototipado y pruebas de usabilidad antes de la construcción definitiva, y documentará los hallazgos y los cambios de diseño que produjeron.\newline\textit{BT; Obligatorio.} &
Sí &
Panitao Design System; WCAG 2.2 AA; Modo Cubierta &
SD4, sec. 4.1; SD9, sec. 9.2 \\ \hline

RT-\allowbreak{}13.\allowbreak{}04 &
Se comprometerán indicadores de usabilidad medibles: tiempo máximo de la transacción operacional crítica, número máximo de pasos, tasa de error tolerada y tiempo de aprendizaje esperado por perfil.\newline\textit{BT; Obligatorio.} &
Sí &
Panitao Design System; WCAG 2.2 AA; Modo Cubierta &
SD4, sec. 4.1; SD9, sec. 9.2 \\ \hline

RT-\allowbreak{}13.\allowbreak{}05 &
Ninguna funcionalidad principal requerirá más de tres interacciones desde la pantalla de inicio del perfil correspondiente.\newline\textit{BT; Obligatorio.} &
Sí &
Panitao Design System; WCAG 2.2 AA; Modo Cubierta &
SD4, sec. 4.1; SD9, sec. 9.2 \\ \hline

RT-\allowbreak{}13.\allowbreak{}06 &
La solución entregará retroalimentación visual clara ante cada acción, manejará los errores de forma comprensible —indicando qué ocurrió y qué hacer— y evitará mensajes técnicos dirigidos a la persona usuaria final.\newline\textit{BT; Obligatorio.} &
Sí &
Panitao Design System; WCAG 2.2 AA; Modo Cubierta &
SD4, sec. 4.1; SD9, sec. 9.2 \\ \hline

RT-\allowbreak{}13.\allowbreak{}07 &
La solución soportará a personas usuarias con baja alfabetización digital: alto contraste, objetivos táctiles de al menos 44 × 44 píxeles, iconografía acompañada de texto, y flujos guiados paso a paso.\newline\textit{BT; Obligatorio.} &
Sí &
Panitao Design System; WCAG 2.2 AA; Modo Cubierta &
SD4, sec. 4.1; SD9, sec. 9.2 \\ \hline

RT-\allowbreak{}13.\allowbreak{}08 &
Cuando el caso lo requiera, las interfaces de terreno operarán con guantes, a la intemperie, con luminosidad variable y sin conexión.\newline\textit{BT; Según caso.} &
Sí &
Panitao Design System; WCAG 2.2 AA; Modo Cubierta &
SD4, sec. 4.1; SD9, sec. 9.2 \\ \hline

RT-\allowbreak{}13.\allowbreak{}09 &
Existirá un sistema de diseño documentado con paleta de colores acotada, jerarquía tipográfica de a lo más dos familias, iconografía coherente, retícula y componentes reutilizables.\newline\textit{BT; Obligatorio.} &
Sí &
Panitao Design System; WCAG 2.2 AA; Modo Cubierta &
SD4, sec. 4.1; SD9, sec. 9.2 \\ \hline

RT-\allowbreak{}13.\allowbreak{}10 &
El PROPONENTE declarará la matriz de navegadores y de versiones soportadas y su política de actualización.\newline\textit{BT; Obligatorio.} &
Sí &
Panitao Design System; WCAG 2.2 AA; Modo Cubierta &
SD4, sec. 4.1; SD9, sec. 9.2 \\ \hline

RT-\allowbreak{}13.\allowbreak{}11 &
La solución será navegable íntegramente por teclado, con orden de foco lógico y atajos para las operaciones frecuentes.\newline\textit{BT; Obligatorio.} &
Sí &
Panitao Design System; WCAG 2.2 AA; Modo Cubierta &
SD4, sec. 4.1; SD9, sec. 9.2 \\ \hline

RT-\allowbreak{}13.\allowbreak{}12 &
Se valorará el soporte de modo oscuro, de personalización de la interfaz por persona usuaria y de múltiples idiomas cuando el caso lo justifique.\newline\textit{BT; Deseable.} &
Sí &
Modo oscuro automático/manual para visión nocturna y bilingüismo integral en comunicaciones marítimas &
SD3, sec. 3.4.1; SD4, sec. 4.1 \\ \hline

\end{longtable}


\subsection*{BT, capítulo 14: OBSERVABILIDAD Y GESTIÓN DEL SERVICIO}

\addcontentsline{toc}{subsection}{BT, capítulo 14: OBSERVABILIDAD Y GESTIÓN DEL SERVICIO}

La Tabla~\ref{tab:t12-bt-cap-14} individualiza los requisitos técnicos de este ámbito y su localización en la propuesta.

\begin{longtable}{|L{0.09\AnchoMatrizTDoce}|L{0.385\AnchoMatrizTDoce}|L{0.065\AnchoMatrizTDoce}|L{0.18\AnchoMatrizTDoce}|L{0.28\AnchoMatrizTDoce}|}

\caption{Cumplimiento y trazabilidad: BT, capítulo 14: OBSERVABILIDAD Y GESTIÓN DEL SERVICIO}\label{tab:t12-bt-cap-14}\\

\hline

\EncabezadoTabla{ID requerimiento} & \EncabezadoTabla{Descripción} & \EncabezadoTabla{Cumple\newline(Sí/No)} & \EncabezadoTabla{Componente que lo satisface} & \EncabezadoTabla{Sección de la propuesta} \\ \hline

\endfirsthead

\hline

\EncabezadoTabla{ID requerimiento} & \EncabezadoTabla{Descripción} & \EncabezadoTabla{Cumple\newline(Sí/No)} & \EncabezadoTabla{Componente que lo satisface} & \EncabezadoTabla{Sección de la propuesta} \\ \hline

\endhead

\multicolumn{5}{r}{\fontsize{9}{11}\selectfont Continúa en la página siguiente.}\\

\endfoot

\hline

\endlastfoot


RT-\allowbreak{}14.\allowbreak{}01 &
La observabilidad será unificada para nube y on-premise, con métricas, registros y trazas distribuidas correlacionadas por un identificador único de transacción, instrumentadas conforme a OpenTelemetry.\newline\textit{BT; Obligatorio.} &
Sí &
OpenTelemetry; CloudWatch; Prometheus/Grafana &
SD4, sec. 4.1.2; SD4, sec. 4.2 \\ \hline

RT-\allowbreak{}14.\allowbreak{}02 &
El CLIENTE dispondrá de acceso propio y permanente a los tableros operacionales y de negocio, con datos en tiempo real y capacidad de exportación.\newline\textit{BT; Obligatorio.} &
Sí &
OpenTelemetry; CloudWatch; Prometheus/Grafana &
SD4, sec. 4.1.2; SD4, sec. 4.2 \\ \hline

RT-\allowbreak{}14.\allowbreak{}03 &
Los indicadores de nivel de servicio se medirán sobre la experiencia real de la persona usuaria y no sobre pruebas sintéticas, sin perjuicio de que estas se empleen como complemento.\newline\textit{BT; Obligatorio.} &
Sí &
OpenTelemetry; CloudWatch; Prometheus/Grafana &
SD4, sec. 4.1.2; SD4, sec. 4.2 \\ \hline

RT-\allowbreak{}14.\allowbreak{}04 &
El alertamiento se basará en síntomas de negocio y no sólo en umbrales de infraestructura, con supresión de ruido, agrupación, escalamiento automático y turnos de disponibilidad declarados.\newline\textit{BT; Obligatorio.} &
Sí &
OpenTelemetry; CloudWatch; Prometheus/Grafana &
SD4, sec. 4.1.2; SD4, sec. 4.2 \\ \hline

RT-\allowbreak{}14.\allowbreak{}05 &
Existirá un libro de operación y una guía de resolución documentados para cada escenario de falla previsible, con automatización progresiva de las tareas repetitivas.\newline\textit{BT; Obligatorio.} &
Sí &
OpenTelemetry; CloudWatch; Prometheus/Grafana &
SD4, sec. 4.1.2; SD4, sec. 4.2 \\ \hline

RT-\allowbreak{}14.\allowbreak{}06 &
Todo incidente crítico dará lugar a un análisis de causa raíz obligatorio, con informe entregable dentro de cinco días hábiles y seguimiento de las acciones correctivas hasta su cierre.\newline\textit{BT; Obligatorio.} &
Sí &
OpenTelemetry; CloudWatch; Prometheus/Grafana &
SD4, sec. 4.1.2; SD4, sec. 4.2 \\ \hline

RT-\allowbreak{}14.\allowbreak{}07 &
Los registros de la solución no contendrán datos personales sensibles ni credenciales, y su acceso estará controlado y auditado.\newline\textit{BT; Obligatorio.} &
Sí &
OpenTelemetry; CloudWatch; Prometheus/Grafana &
SD4, sec. 4.1.2; SD4, sec. 4.2 \\ \hline

RT-\allowbreak{}14.\allowbreak{}08 &
El PROPONENTE declarará la retención de métricas, registros y trazas, y su costo asociado, distinguiendo el almacenamiento en línea del archivado.\newline\textit{BT; Obligatorio.} &
Sí &
OpenTelemetry; CloudWatch; Prometheus/Grafana &
SD4, sec. 4.1.2; SD4, sec. 4.2 \\ \hline

RT-\allowbreak{}14.\allowbreak{}09 &
Se valorará la detección proactiva de anomalías mediante análisis del comportamiento histórico, con alerta antes de que el incidente afecte a la operación.\newline\textit{BT; Deseable.} &
No &
No aplica / Descartado: Detección de anomalías mediante umbrales y alarmas en CloudWatch/OpenTelemetry; sin motor no supervisado de series temporales &
SD3, sec. 3.2.2; SD4, sec. 4.2 \\ \hline

\end{longtable}


\subsection*{BT, capítulo 15: SOSTENIBILIDAD, EFICIENCIA Y CERTIFICACIONES}

\addcontentsline{toc}{subsection}{BT, capítulo 15: SOSTENIBILIDAD, EFICIENCIA Y CERTIFICACIONES}

La Tabla~\ref{tab:t12-bt-cap-15} individualiza los requisitos técnicos de este ámbito y su localización en la propuesta.

\begin{longtable}{|L{0.09\AnchoMatrizTDoce}|L{0.385\AnchoMatrizTDoce}|L{0.065\AnchoMatrizTDoce}|L{0.18\AnchoMatrizTDoce}|L{0.28\AnchoMatrizTDoce}|}

\caption{Cumplimiento y trazabilidad: BT, capítulo 15: SOSTENIBILIDAD, EFICIENCIA Y CERTIFICACIONES}\label{tab:t12-bt-cap-15}\\

\hline

\EncabezadoTabla{ID requerimiento} & \EncabezadoTabla{Descripción} & \EncabezadoTabla{Cumple\newline(Sí/No)} & \EncabezadoTabla{Componente que lo satisface} & \EncabezadoTabla{Sección de la propuesta} \\ \hline

\endfirsthead

\hline

\EncabezadoTabla{ID requerimiento} & \EncabezadoTabla{Descripción} & \EncabezadoTabla{Cumple\newline(Sí/No)} & \EncabezadoTabla{Componente que lo satisface} & \EncabezadoTabla{Sección de la propuesta} \\ \hline

\endhead

\multicolumn{5}{r}{\fontsize{9}{11}\selectfont Continúa en la página siguiente.}\\

\endfoot

\hline

\endlastfoot


RT-\allowbreak{}15.\allowbreak{}01 &
El PROPONENTE dimensionará la infraestructura ajustada a la demanda real, evitando capacidad ociosa permanente, y declarará el factor de utilización proyectado.\newline\textit{BT; Obligatorio.} &
Sí &
Servicio Gestionado Synaptix; Parches y Soporte &
SD3, sec. 3.1.2; SD10, sec. 10.1 \\ \hline

RT-\allowbreak{}15.\allowbreak{}02 &
Los ambientes no productivos se apagarán o reducirán fuera del horario de uso.\newline\textit{BT; Obligatorio.} &
Sí &
Servicio Gestionado Synaptix; Parches y Soporte &
SD3, sec. 3.1.2; SD10, sec. 10.1 \\ \hline

RT-\allowbreak{}15.\allowbreak{}03 &
El PROPONENTE estimará la huella de carbono anual de la operación de la solución y declarará la metodología empleada.\newline\textit{BT; Obligatorio.} &
Sí &
Servicio Gestionado Synaptix; Parches y Soporte &
SD3, sec. 3.1.2; SD10, sec. 10.1 \\ \hline

RT-\allowbreak{}15.\allowbreak{}04 &
Se declarará la eficiencia en el uso de la energía (PUE) del recinto on-premise y la intensidad de carbono de la región de nube escogida.\newline\textit{BT; Obligatorio.} &
Sí &
Servicio Gestionado Synaptix; Parches y Soporte &
SD3, sec. 3.1.2; SD10, sec. 10.1 \\ \hline

RT-\allowbreak{}15.\allowbreak{}05 &
Se valorará la elección de regiones de nube con menor intensidad de carbono cuando la latencia y la regulación lo permitan, con el análisis comparativo correspondiente.\newline\textit{BT; Deseable.} &
No &
No aplica / Descartado: Cómputo primario alojado obligatoriamente en AWS Santiago (Región Chile) por exigencia de soberanía y baja latencia &
SD3, sec. 3.2.2; SD4, sec. 4.2 \\ \hline

RT-\allowbreak{}15.\allowbreak{}06 &
Se valorará la definición de metas de reducción del consumo durante la Operación, con medición y reporte anual.\newline\textit{BT; Deseable.} &
No &
No aplica / Descartado: Monitoreo de eficiencia ambiental se acota a los consumos reales de agua y energía de naves en muelle (CA-18 y CA-19) &
SD3, sec. 3.2.2; SD4, sec. 4.1 \\ \hline

RT-\allowbreak{}15.\allowbreak{}07 &
Las certificaciones se acreditarán con copia del certificado vigente y con el código de verificación del organismo emisor cuando exista.\newline\textit{BT; Obligatorio.} &
Sí &
Servicio Gestionado Synaptix; Parches y Soporte &
SD3, sec. 3.1.2; SD10, sec. 10.1 \\ \hline

RT-\allowbreak{}15.\allowbreak{}08 &
Las personas certificadas formarán parte del equipo efectivamente asignado al PROYECTO, con dedicación declarada. No se aceptará acreditar personal que no participe.\newline\textit{BT; Obligatorio.} &
Sí &
Servicio Gestionado Synaptix; Parches y Soporte &
SD3, sec. 3.1.2; SD10, sec. 10.1 \\ \hline

RT-\allowbreak{}15.\allowbreak{}09 &
El ADJUDICATARIO mantendrá vigentes estas certificaciones durante todo el período contractual y las repondrá ante la salida de una persona certificada, conforme al Artículo 76° de las Bases Administrativas.\newline\textit{BT; Obligatorio.} &
Sí &
Servicio Gestionado Synaptix; Parches y Soporte &
SD3, sec. 3.1.2; SD10, sec. 10.1 \\ \hline

\end{longtable}


\subsection*{BT, capítulo 16: MÓDULOS TRANSVERSALES OBLIGATORIOS Y REGLAS}

\addcontentsline{toc}{subsection}{BT, capítulo 16: MÓDULOS TRANSVERSALES OBLIGATORIOS Y REGLAS}

La Tabla~\ref{tab:t12-bt-cap-16} individualiza los requisitos técnicos de este ámbito y su localización en la propuesta.

\begin{longtable}{|L{0.09\AnchoMatrizTDoce}|L{0.385\AnchoMatrizTDoce}|L{0.065\AnchoMatrizTDoce}|L{0.18\AnchoMatrizTDoce}|L{0.28\AnchoMatrizTDoce}|}

\caption{Cumplimiento y trazabilidad: BT, capítulo 16: MÓDULOS TRANSVERSALES OBLIGATORIOS Y REGLAS}\label{tab:t12-bt-cap-16}\\

\hline

\EncabezadoTabla{ID requerimiento} & \EncabezadoTabla{Descripción} & \EncabezadoTabla{Cumple\newline(Sí/No)} & \EncabezadoTabla{Componente que lo satisface} & \EncabezadoTabla{Sección de la propuesta} \\ \hline

\endfirsthead

\hline

\EncabezadoTabla{ID requerimiento} & \EncabezadoTabla{Descripción} & \EncabezadoTabla{Cumple\newline(Sí/No)} & \EncabezadoTabla{Componente que lo satisface} & \EncabezadoTabla{Sección de la propuesta} \\ \hline

\endhead

\multicolumn{5}{r}{\fontsize{9}{11}\selectfont Continúa en la página siguiente.}\\

\endfoot

\hline

\endlastfoot


RT-\allowbreak{}16.\allowbreak{}01 &
La solución dispondrá de un módulo de administración que permita al CLIENTE, sin intervención del ADJUDICATARIO, gestionar personas usuarias, roles, permisos, unidades organizacionales y sus jerarquías.\newline\textit{BT; Obligatorio.} &
Sí &
Spring Modulith; Motor de Reglas Paramétricas &
SD4, sec. 4.1 \\ \hline

RT-\allowbreak{}16.\allowbreak{}02 &
Las reglas de negocio parametrizables —umbrales, plazos, montos, tolerancias, catálogos, listas de valores, textos de notificación— serán configurables desde la interfaz de administración, con control de versiones y registro de quién cambió qué y cuándo.\newline\textit{BT; Obligatorio.} &
Sí &
Spring Modulith; Motor de Reglas Paramétricas &
SD4, sec. 4.1 \\ \hline

RT-\allowbreak{}16.\allowbreak{}03 &
Todo cambio de parámetro con impacto operacional requerirá aprobación de un segundo perfil y quedará registrado con su justificación.\newline\textit{BT; Obligatorio.} &
Sí &
Spring Modulith; Motor de Reglas Paramétricas &
SD4, sec. 4.1 \\ \hline

RT-\allowbreak{}16.\allowbreak{}04 &
El PROPONENTE declarará expresamente qué elementos son parametrizables y cuáles requieren desarrollo. Presentar como parametrizable lo que exige desarrollo se evaluará como observación grave.\newline\textit{BT; Obligatorio.} &
Sí &
Spring Modulith; Motor de Reglas Paramétricas &
SD4, sec. 4.1 \\ \hline

RT-\allowbreak{}16.\allowbreak{}05 &
Existirá un ambiente de simulación que permita probar el efecto de un cambio de parámetro antes de aplicarlo a producción.\newline\textit{BT; Deseable.} &
No &
No aplica / Descartado: Pruebas de reglas se realizan en ambiente de homologación QA; no se requiere motor de simulación montecarlo en producción &
SD3, sec. 3.2.2; SD6, sec. 6.2 \\ \hline

RT-\allowbreak{}16.\allowbreak{}06 &
Toda operación que cree, modifique o elimine información quedará registrada con identificación de la persona o del sistema que la ejecutó, fecha y hora con zona horaria, origen, valores anteriores y valores posteriores.\newline\textit{BT; Obligatorio.} &
Sí &
Spring Modulith; Motor de Reglas Paramétricas &
SD4, sec. 4.1 \\ \hline

RT-\allowbreak{}16.\allowbreak{}07 &
El registro de auditoría será inalterable y no podrá ser modificado ni eliminado por ningún perfil, incluido el administrador de la plataforma.\newline\textit{BT; Obligatorio.} &
Sí &
Spring Modulith; Motor de Reglas Paramétricas &
SD4, sec. 4.1 \\ \hline

RT-\allowbreak{}16.\allowbreak{}08 &
El CLIENTE podrá consultar y exportar la auditoría desde la interfaz, con filtros por persona, período, entidad y tipo de operación, sin requerir acceso a la base de datos.\newline\textit{BT; Obligatorio.} &
Sí &
Spring Modulith; Motor de Reglas Paramétricas &
SD4, sec. 4.1 \\ \hline

RT-\allowbreak{}16.\allowbreak{}09 &
Las consultas a información sensible quedarán registradas, no sólo las modificaciones.\newline\textit{BT; Según caso.} &
Sí &
Spring Modulith; Motor de Reglas Paramétricas &
SD4, sec. 4.1 \\ \hline

RT-\allowbreak{}16.\allowbreak{}10 &
El período de retención de la auditoría será el que fije el caso y, en su defecto, no inferior a cinco años.\newline\textit{BT; Según caso.} &
Sí &
Spring Modulith; Motor de Reglas Paramétricas &
SD4, sec. 4.1 \\ \hline

RT-\allowbreak{}16.\allowbreak{}11 &
La solución soportará flujos de trabajo con estados, transiciones, responsables, plazos, escalamiento automático por vencimiento y delegación por ausencia.\newline\textit{BT; Obligatorio.} &
Sí &
Spring Modulith; Motor de Reglas Paramétricas &
SD4, sec. 4.1 \\ \hline

RT-\allowbreak{}16.\allowbreak{}12 &
Los flujos serán configurables por el CLIENTE sin desarrollo, al menos en lo relativo a responsables, plazos y niveles de aprobación.\newline\textit{BT; Obligatorio.} &
Sí &
Spring Modulith; Motor de Reglas Paramétricas &
SD4, sec. 4.1 \\ \hline

RT-\allowbreak{}16.\allowbreak{}13 &
Toda solicitud pendiente será visible para su responsable en una bandeja de tareas unificada, con priorización y alerta de vencimiento.\newline\textit{BT; Obligatorio.} &
Sí &
Spring Modulith; Motor de Reglas Paramétricas &
SD4, sec. 4.1 \\ \hline

RT-\allowbreak{}16.\allowbreak{}14 &
El motor de reglas permitirá definir condiciones de negocio evaluables sin recompilación, con trazabilidad de qué regla se aplicó a cada transacción.\newline\textit{BT; Deseable.} &
No &
No aplica / Descartado: Reglas de negocio gobernadas por motor paramétrico del Módulo ADM; no requiere motor Drools/DMN desacoplado externo &
SD3, sec. 3.2.2; SD4, sec. 4.1 \\ \hline

RT-\allowbreak{}16.\allowbreak{}15 &
La solución gestionará documentos con versionado, metadatos, control de acceso, búsqueda por contenido y por metadato, y previsualización sin descarga.\newline\textit{BT; Obligatorio.} &
Sí &
Spring Modulith; Motor de Reglas Paramétricas &
SD4, sec. 4.1 \\ \hline

RT-\allowbreak{}16.\allowbreak{}16 &
Los documentos se almacenarán cifrados, con verificación de integridad y retención conforme a la política del caso.\newline\textit{BT; Obligatorio.} &
Sí &
Spring Modulith; Motor de Reglas Paramétricas &
SD4, sec. 4.1 \\ \hline

RT-\allowbreak{}16.\allowbreak{}17 &
La solución soportará firma electrónica conforme a la Ley N° 19.799, con firma avanzada para los actos que el caso lo requiera, y verificación de validez del certificado al momento de la firma.\newline\textit{BT; Según caso.} &
Sí &
Spring Modulith; Motor de Reglas Paramétricas &
SD4, sec. 4.1 \\ \hline

RT-\allowbreak{}16.\allowbreak{}18 &
Se generará el sello de tiempo y se conservará la evidencia de firma que permita verificar el documento con posterioridad al vencimiento del certificado.\newline\textit{BT; Según caso.} &
Sí &
Spring Modulith; Motor de Reglas Paramétricas &
SD4, sec. 4.1 \\ \hline

RT-\allowbreak{}16.\allowbreak{}19 &
La solución generará documentos a partir de plantillas administrables por el CLIENTE, con datos de la transacción y salida en formato abierto.\newline\textit{BT; Obligatorio.} &
Sí &
Spring Modulith; Motor de Reglas Paramétricas &
SD4, sec. 4.1 \\ \hline

RT-\allowbreak{}16.\allowbreak{}20 &
La solución enviará notificaciones por al menos tres canales: correo electrónico, notificación en la aplicación y mensajería instantánea o SMS, según lo que el caso requiera.\newline\textit{BT; Obligatorio.} &
Sí &
Spring Modulith; Motor de Reglas Paramétricas &
SD4, sec. 4.1 \\ \hline

RT-\allowbreak{}16.\allowbreak{}21 &
Las plantillas de notificación serán administrables por el CLIENTE, con variables de la transacción, y versionadas.\newline\textit{BT; Obligatorio.} &
Sí &
Spring Modulith; Motor de Reglas Paramétricas &
SD4, sec. 4.1 \\ \hline

RT-\allowbreak{}16.\allowbreak{}22 &
Cada persona usuaria podrá configurar sus preferencias de canal y de frecuencia, respetando las notificaciones que el CLIENTE defina como obligatorias.\newline\textit{BT; Obligatorio.} &
Sí &
Spring Modulith; Motor de Reglas Paramétricas &
SD4, sec. 4.1 \\ \hline

RT-\allowbreak{}16.\allowbreak{}23 &
El envío será asíncrono, con reintento ante falla, control de duplicados y registro de entrega, apertura y error por cada mensaje.\newline\textit{BT; Obligatorio.} &
Sí &
Spring Modulith; Motor de Reglas Paramétricas &
SD4, sec. 4.1 \\ \hline

RT-\allowbreak{}16.\allowbreak{}24 &
El PROPONENTE declarará el proveedor de cada canal, su costo unitario, su volumen proyectado y el tratamiento del costo variable en la Oferta Económica.\newline\textit{BT; Obligatorio.} &
Sí &
Spring Modulith; Motor de Reglas Paramétricas &
SD4, sec. 4.1 \\ \hline

RT-\allowbreak{}16.\allowbreak{}25 &
Las notificaciones respetarán la normativa de comunicaciones comerciales y permitirán la baja cuando corresponda.\newline\textit{BT; Obligatorio.} &
Sí &
Spring Modulith; Motor de Reglas Paramétricas &
SD4, sec. 4.1 \\ \hline

RT-\allowbreak{}16.\allowbreak{}26 &
Se valorará la integración con canales conversacionales que permitan a la persona usuaria responder y ejecutar acciones desde el propio canal.\newline\textit{BT; Deseable.} &
No &
No aplica / Descartado: Notificaciones a armadores operan vía WhatsApp Business API, SMS y correo; no incluye chatbot conversacional desatendido &
SD3, sec. 3.2.2; SD4, sec. 4.1 \\ \hline

RT-\allowbreak{}16.\allowbreak{}27 &
La solución dispondrá de búsqueda global con indexación de texto completo, tolerancia a errores de escritura, filtros facetados y respeto del control de acceso de la persona que busca.\newline\textit{BT; Obligatorio.} &
Sí &
Spring Modulith; Motor de Reglas Paramétricas &
SD4, sec. 4.1 \\ \hline

RT-\allowbreak{}16.\allowbreak{}28 &
Los listados serán ordenables, filtrables, paginados y exportables en formatos abiertos, con el filtro aplicado reflejado en la exportación.\newline\textit{BT; Obligatorio.} &
Sí &
Spring Modulith; Motor de Reglas Paramétricas &
SD4, sec. 4.1 \\ \hline

RT-\allowbreak{}16.\allowbreak{}29 &
Las exportaciones de gran volumen se procesarán de forma asíncrona, con notificación al completarse y sin bloquear la sesión.\newline\textit{BT; Obligatorio.} &
Sí &
Spring Modulith; Motor de Reglas Paramétricas &
SD4, sec. 4.1 \\ \hline

RT-\allowbreak{}16.\allowbreak{}30 &
Toda exportación de información sensible quedará registrada en la auditoría, con identificación de quién exportó qué y cuándo.\newline\textit{BT; Obligatorio.} &
Sí &
Spring Modulith; Motor de Reglas Paramétricas &
SD4, sec. 4.1 \\ \hline

RT-\allowbreak{}16.\allowbreak{}31 &
La solución dispondrá de un portal público con la información que el caso determine, accesible sin autenticación, con los mismos estándares de accesibilidad y desempeño que el resto de la plataforma.\newline\textit{BT; Según caso.} &
Sí &
Spring Modulith; Motor de Reglas Paramétricas &
SD4, sec. 4.1 \\ \hline

RT-\allowbreak{}16.\allowbreak{}32 &
Las personas usuarias externas dispondrán de autoatención para las consultas de mayor frecuencia, evitando el contacto telefónico para operaciones simples.\newline\textit{BT; Obligatorio.} &
Sí &
Spring Modulith; Motor de Reglas Paramétricas &
SD4, sec. 4.1 \\ \hline

RT-\allowbreak{}16.\allowbreak{}33 &
El PROPONENTE estimará la reducción esperada del volumen de atención asistida por efecto de la autoatención y comprometerá el indicador.\newline\textit{BT; Obligatorio.} &
Sí &
Spring Modulith; Motor de Reglas Paramétricas &
SD4, sec. 4.1 \\ \hline

RT-\allowbreak{}16.\allowbreak{}34 &
El portal público resistirá picos de tráfico sin degradar los servicios transaccionales internos, mediante aislamiento de recursos y caché.\newline\textit{BT; Obligatorio.} &
Sí &
Spring Modulith; Motor de Reglas Paramétricas &
SD4, sec. 4.1 \\ \hline

\end{longtable}


\subsection*{BT, capítulo 17: CANALES DIGITALES Y MOVILIDAD}

\addcontentsline{toc}{subsection}{BT, capítulo 17: CANALES DIGITALES Y MOVILIDAD}

La Tabla~\ref{tab:t12-bt-cap-17} individualiza los requisitos técnicos de este ámbito y su localización en la propuesta.

\begin{longtable}{|L{0.09\AnchoMatrizTDoce}|L{0.385\AnchoMatrizTDoce}|L{0.065\AnchoMatrizTDoce}|L{0.18\AnchoMatrizTDoce}|L{0.28\AnchoMatrizTDoce}|}

\caption{Cumplimiento y trazabilidad: BT, capítulo 17: CANALES DIGITALES Y MOVILIDAD}\label{tab:t12-bt-cap-17}\\

\hline

\EncabezadoTabla{ID requerimiento} & \EncabezadoTabla{Descripción} & \EncabezadoTabla{Cumple\newline(Sí/No)} & \EncabezadoTabla{Componente que lo satisface} & \EncabezadoTabla{Sección de la propuesta} \\ \hline

\endfirsthead

\hline

\EncabezadoTabla{ID requerimiento} & \EncabezadoTabla{Descripción} & \EncabezadoTabla{Cumple\newline(Sí/No)} & \EncabezadoTabla{Componente que lo satisface} & \EncabezadoTabla{Sección de la propuesta} \\ \hline

\endhead

\multicolumn{5}{r}{\fontsize{9}{11}\selectfont Continúa en la página siguiente.}\\

\endfoot

\hline

\endlastfoot


RT-\allowbreak{}17.\allowbreak{}01 &
La solución proveerá una aplicación móvil para los perfiles operacionales que el caso identifique, con funcionamiento sin conexión y sincronización diferida cuando la operación en terreno lo requiera.\newline\textit{BT; Según caso.} &
Sí &
Portal Armador PWA; App Terreno Android &
SD4, sec. 4.1; SD4, sec. 4.2 \\ \hline

RT-\allowbreak{}17.\allowbreak{}02 &
El PROPONENTE declarará si la aplicación es nativa, híbrida o web progresiva, y justificará la decisión frente al requisito de operación desconectada, al acceso a periféricos y al costo de mantención.\newline\textit{BT; Obligatorio.} &
Sí &
Portal Armador PWA; App Terreno Android &
SD4, sec. 4.1; SD4, sec. 4.2 \\ \hline

RT-\allowbreak{}17.\allowbreak{}03 &
La aplicación soportará las versiones de sistema operativo móvil vigentes y las dos anteriores, con política de actualización declarada.\newline\textit{BT; Obligatorio.} &
Sí &
Portal Armador PWA; App Terreno Android &
SD4, sec. 4.1; SD4, sec. 4.2 \\ \hline

RT-\allowbreak{}17.\allowbreak{}04 &
La aplicación se distribuirá por las tiendas oficiales o mediante gestión de flota corporativa, con firma de la aplicación y verificación de integridad.\newline\textit{BT; Obligatorio.} &
Sí &
Portal Armador PWA; App Terreno Android &
SD4, sec. 4.1; SD4, sec. 4.2 \\ \hline

RT-\allowbreak{}17.\allowbreak{}05 &
La información almacenada en el dispositivo estará cifrada, con borrado remoto y bloqueo ante pérdida o desvinculación de la persona usuaria.\newline\textit{BT; Obligatorio.} &
Sí &
Portal Armador PWA; App Terreno Android &
SD4, sec. 4.1; SD4, sec. 4.2 \\ \hline

RT-\allowbreak{}17.\allowbreak{}06 &
La aplicación integrará los periféricos que el caso requiera: cámara, lector de códigos, NFC, GPS, impresora de etiquetas, balanza o báscula.\newline\textit{BT; Según caso.} &
Sí &
Portal Armador PWA; App Terreno Android &
SD4, sec. 4.1; SD4, sec. 4.2 \\ \hline

RT-\allowbreak{}17.\allowbreak{}07 &
El consumo de datos móviles y de batería se optimizará y se declarará el consumo estimado por turno de trabajo.\newline\textit{BT; Obligatorio.} &
Sí &
Portal Armador PWA; App Terreno Android &
SD4, sec. 4.1; SD4, sec. 4.2 \\ \hline

RT-\allowbreak{}17.\allowbreak{}08 &
Se valorará la disponibilidad de una versión de la interfaz para dispositivos de bajo costo o de generaciones anteriores, ampliando la cobertura de personas usuarias.\newline\textit{BT; Deseable.} &
No &
No aplica / Descartado: Aplicación móvil optimizada para tablets rugerizadas de 10 pulgadas en Modo Cubierta; no requiere variante específica para smartwatches &
SD3, sec. 3.2.2; SD4, sec. 4.1 \\ \hline

\end{longtable}


\subsection*{BT, capítulo 18: INTELIGENCIA ARTIFICIAL Y AUTOMATIZACIÓN}

\addcontentsline{toc}{subsection}{BT, capítulo 18: INTELIGENCIA ARTIFICIAL Y AUTOMATIZACIÓN}

La Tabla~\ref{tab:t12-bt-cap-18} individualiza los requisitos técnicos de este ámbito y su localización en la propuesta.

\begin{longtable}{|L{0.09\AnchoMatrizTDoce}|L{0.385\AnchoMatrizTDoce}|L{0.065\AnchoMatrizTDoce}|L{0.18\AnchoMatrizTDoce}|L{0.28\AnchoMatrizTDoce}|}

\caption{Cumplimiento y trazabilidad: BT, capítulo 18: INTELIGENCIA ARTIFICIAL Y AUTOMATIZACIÓN}\label{tab:t12-bt-cap-18}\\

\hline

\EncabezadoTabla{ID requerimiento} & \EncabezadoTabla{Descripción} & \EncabezadoTabla{Cumple\newline(Sí/No)} & \EncabezadoTabla{Componente que lo satisface} & \EncabezadoTabla{Sección de la propuesta} \\ \hline

\endfirsthead

\hline

\EncabezadoTabla{ID requerimiento} & \EncabezadoTabla{Descripción} & \EncabezadoTabla{Cumple\newline(Sí/No)} & \EncabezadoTabla{Componente que lo satisface} & \EncabezadoTabla{Sección de la propuesta} \\ \hline

\endhead

\multicolumn{5}{r}{\fontsize{9}{11}\selectfont Continúa en la página siguiente.}\\

\endfoot

\hline

\endlastfoot


RT-\allowbreak{}18.\allowbreak{}01 &
El PROPONENTE declarará cada componente de inteligencia artificial: propósito, modelo o servicio empleado, proveedor, versión y ubicación de procesamiento de los datos.\newline\textit{BT; Obligatorio.} &
Sí &
Motor Predictivo y Detección de Discrepancias &
SD4, sec. 4.1.2; SD13, sec. 13.1 \\ \hline

RT-\allowbreak{}18.\allowbreak{}02 &
Se garantizará contractualmente que los datos del CLIENTE no serán utilizados para entrenar modelos de terceros, salvo autorización expresa y escrita.\newline\textit{BT; Obligatorio.} &
Sí &
Motor Predictivo y Detección de Discrepancias &
SD4, sec. 4.1.2; SD13, sec. 13.1 \\ \hline

RT-\allowbreak{}18.\allowbreak{}03 &
Se documentarán los límites de uso, los casos en que el resultado requiere validación humana previa y el procedimiento de supervisión.\newline\textit{BT; Obligatorio.} &
Sí &
Motor Predictivo y Detección de Discrepancias &
SD4, sec. 4.1.2; SD13, sec. 13.1 \\ \hline

RT-\allowbreak{}18.\allowbreak{}04 &
Los riesgos de sesgo, alucinación, fuga de información y uso indebido se evaluarán y mitigarán conforme al NIST AI Risk Management Framework 1.0 y a la norma ISO/IEC 42001.\newline\textit{BT; Obligatorio.} &
Sí &
Motor Predictivo y Detección de Discrepancias &
SD4, sec. 4.1.2; SD13, sec. 13.1 \\ \hline

RT-\allowbreak{}18.\allowbreak{}05 &
Toda interacción relevante con un componente de inteligencia artificial quedará registrada para efectos de auditoría y trazabilidad, incluyendo la entrada, la salida y la decisión humana posterior.\newline\textit{BT; Obligatorio.} &
Sí &
Motor Predictivo y Detección de Discrepancias &
SD4, sec. 4.1.2; SD13, sec. 13.1 \\ \hline

RT-\allowbreak{}18.\allowbreak{}06 &
El componente podrá desactivarse sin comprometer la operación del resto de la solución, y existirá un procedimiento manual de respaldo para la función que automatiza.\newline\textit{BT; Obligatorio.} &
Sí &
Motor Predictivo y Detección de Discrepancias &
SD4, sec. 4.1.2; SD13, sec. 13.1 \\ \hline

RT-\allowbreak{}18.\allowbreak{}07 &
Cuando el resultado se presente a una persona usuaria, se indicará expresamente que fue generado o sugerido de forma automática, y su nivel de confianza cuando el modelo lo provea.\newline\textit{BT; Obligatorio.} &
Sí &
Motor Predictivo y Detección de Discrepancias &
SD4, sec. 4.1.2; SD13, sec. 13.1 \\ \hline

RT-\allowbreak{}18.\allowbreak{}08 &
Los modelos predictivos declararán sus variables de entrada, su métrica de desempeño, su línea base y su plan de reentrenamiento y de detección de deriva.\newline\textit{BT; Obligatorio.} &
Sí &
Motor Predictivo y Detección de Discrepancias &
SD4, sec. 4.1.2; SD13, sec. 13.1 \\ \hline

RT-\allowbreak{}18.\allowbreak{}09 &
La responsabilidad por los resultados de los componentes de inteligencia artificial recae íntegramente en el ADJUDICATARIO, conforme al Artículo 86° de las Bases Administrativas.\newline\textit{BT; Obligatorio.} &
Sí &
Motor Predictivo y Detección de Discrepancias &
SD4, sec. 4.1.2; SD13, sec. 13.1 \\ \hline

RT-\allowbreak{}18.\allowbreak{}10 &
Se valorará el uso de automatización robótica de procesos o de agentes para tareas repetitivas de back office, con el ahorro de horas cuantificado y reflejado en la Oferta Económica.\newline\textit{BT; Deseable.} &
No &
No aplica / Descartado: Automatización de procesos basada en eventos reactivos de Spring Modulith; no aplica software RPA de terceros &
SD3, sec. 3.2.2; SD4, sec. 4.1 \\ \hline

\end{longtable}


\subsection*{BT, capítulo 19: ESTRUCTURA Y GOBIERNO DEL PROYECTO}

\addcontentsline{toc}{subsection}{BT, capítulo 19: ESTRUCTURA Y GOBIERNO DEL PROYECTO}

La Tabla~\ref{tab:t12-bt-cap-19} individualiza los requisitos técnicos de este ámbito y su localización en la propuesta.

\begin{longtable}{|L{0.09\AnchoMatrizTDoce}|L{0.385\AnchoMatrizTDoce}|L{0.065\AnchoMatrizTDoce}|L{0.18\AnchoMatrizTDoce}|L{0.28\AnchoMatrizTDoce}|}

\caption{Cumplimiento y trazabilidad: BT, capítulo 19: ESTRUCTURA Y GOBIERNO DEL PROYECTO}\label{tab:t12-bt-cap-19}\\

\hline

\EncabezadoTabla{ID requerimiento} & \EncabezadoTabla{Descripción} & \EncabezadoTabla{Cumple\newline(Sí/No)} & \EncabezadoTabla{Componente que lo satisface} & \EncabezadoTabla{Sección de la propuesta} \\ \hline

\endfirsthead

\hline

\EncabezadoTabla{ID requerimiento} & \EncabezadoTabla{Descripción} & \EncabezadoTabla{Cumple\newline(Sí/No)} & \EncabezadoTabla{Componente que lo satisface} & \EncabezadoTabla{Sección de la propuesta} \\ \hline

\endhead

\multicolumn{5}{r}{\fontsize{9}{11}\selectfont Continúa en la página siguiente.}\\

\endfoot

\hline

\endlastfoot


RT-\allowbreak{}19.\allowbreak{}01 &
El ADJUDICATARIO constituirá una oficina de gestión del proyecto con metodología declarada, basada en el PMBOK del Project Management Institute e integrando prácticas ágiles donde el PROPONENTE lo justifique.\newline\textit{BT; Obligatorio.} &
Sí &
Metodología Ágil Synaptix; Gobierno y RACI &
SD1, sec. 1.2; SD7, sec. 7.1 \\ \hline

RT-\allowbreak{}19.\allowbreak{}02 &
El Jefe de Proyecto tendrá dedicación exclusiva durante toda la fase de implementación y facultades para comprometer al ADJUDICATARIO en materias de ejecución.\newline\textit{BT; Obligatorio.} &
Sí &
Metodología Ágil Synaptix; Gobierno y RACI &
SD1, sec. 1.2; SD7, sec. 7.1 \\ \hline

RT-\allowbreak{}19.\allowbreak{}03 &
Existirá un procedimiento formal de control de cambios conforme al Artículo 72° de las Bases Administrativas, con registro de cambios, análisis de impacto y aprobación previa a la ejecución.\newline\textit{BT; Obligatorio.} &
Sí &
Metodología Ágil Synaptix; Gobierno y RACI &
SD1, sec. 1.2; SD7, sec. 7.1 \\ \hline

RT-\allowbreak{}19.\allowbreak{}04 &
La gestión del riesgo seguirá la norma ISO 31000, con registro de riesgos vivo, revisión en cada Comité de Proyecto y planes de mitigación con responsable, plazo y disparador.\newline\textit{BT; Obligatorio.} &
Sí &
Metodología Ágil Synaptix; Gobierno y RACI &
SD1, sec. 1.2; SD7, sec. 7.1 \\ \hline

RT-\allowbreak{}19.\allowbreak{}05 &
El ADJUDICATARIO habilitará un espacio colaborativo accesible al CLIENTE, con la documentación del proyecto, los entregables, las actas, el registro de riesgos y el registro de cambios siempre actualizados.\newline\textit{BT; Obligatorio.} &
Sí &
Metodología Ágil Synaptix; Gobierno y RACI &
SD1, sec. 1.2; SD7, sec. 7.1 \\ \hline

RT-\allowbreak{}19.\allowbreak{}06 &
El ADJUDICATARIO entregará un informe mensual de avance con estado del cronograma, avance físico y financiero, entregables del período, desviaciones, riesgos, incidencias y compromisos del período siguiente.\newline\textit{BT; Obligatorio.} &
Sí &
Metodología Ágil Synaptix; Gobierno y RACI &
SD1, sec. 1.2; SD7, sec. 7.1 \\ \hline

RT-\allowbreak{}19.\allowbreak{}07 &
El avance se medirá con valor ganado, reportando el índice de desempeño del cronograma y del costo, y no mediante declaración cualitativa de porcentaje de avance.\newline\textit{BT; Obligatorio.} &
Sí &
Metodología Ágil Synaptix; Gobierno y RACI &
SD1, sec. 1.2; SD7, sec. 7.1 \\ \hline

RT-\allowbreak{}19.\allowbreak{}08 &
El CLIENTE dispondrá de un tablero de estado del proyecto actualizado, accesible en cualquier momento.\newline\textit{BT; Obligatorio.} &
Sí &
Metodología Ágil Synaptix; Gobierno y RACI &
SD1, sec. 1.2; SD7, sec. 7.1 \\ \hline

RT-\allowbreak{}19.\allowbreak{}09 &
Las actas de todos los comités del Artículo 71° de las Bases Administrativas se levantarán dentro de los dos días hábiles siguientes y registrarán acuerdos, responsables y plazos.\newline\textit{BT; Obligatorio.} &
Sí &
Metodología Ágil Synaptix; Gobierno y RACI &
SD1, sec. 1.2; SD7, sec. 7.1 \\ \hline

RT-\allowbreak{}19.\allowbreak{}10 &
Toda desviación superior al 10\% en un hito se comunicará dentro de los cinco días hábiles de detectada, con plan de recuperación.\newline\textit{BT; Obligatorio.} &
Sí &
Metodología Ágil Synaptix; Gobierno y RACI &
SD1, sec. 1.2; SD7, sec. 7.1 \\ \hline

\end{longtable}


\subsection*{BT, capítulo 20: IMPLANTACIÓN, PRUEBAS Y CRITERIOS DE ACEPTACIÓN}

\addcontentsline{toc}{subsection}{BT, capítulo 20: IMPLANTACIÓN, PRUEBAS Y CRITERIOS DE ACEPTACIÓN}

La Tabla~\ref{tab:t12-bt-cap-20} individualiza los requisitos técnicos de este ámbito y su localización en la propuesta.

\begin{longtable}{|L{0.09\AnchoMatrizTDoce}|L{0.385\AnchoMatrizTDoce}|L{0.065\AnchoMatrizTDoce}|L{0.18\AnchoMatrizTDoce}|L{0.28\AnchoMatrizTDoce}|}

\caption{Cumplimiento y trazabilidad: BT, capítulo 20: IMPLANTACIÓN, PRUEBAS Y CRITERIOS DE ACEPTACIÓN}\label{tab:t12-bt-cap-20}\\

\hline

\EncabezadoTabla{ID requerimiento} & \EncabezadoTabla{Descripción} & \EncabezadoTabla{Cumple\newline(Sí/No)} & \EncabezadoTabla{Componente que lo satisface} & \EncabezadoTabla{Sección de la propuesta} \\ \hline

\endfirsthead

\hline

\EncabezadoTabla{ID requerimiento} & \EncabezadoTabla{Descripción} & \EncabezadoTabla{Cumple\newline(Sí/No)} & \EncabezadoTabla{Componente que lo satisface} & \EncabezadoTabla{Sección de la propuesta} \\ \hline

\endhead

\multicolumn{5}{r}{\fontsize{9}{11}\selectfont Continúa en la página siguiente.}\\

\endfoot

\hline

\endlastfoot


RT-\allowbreak{}20.\allowbreak{}01 &
El PROPONENTE presentará un plan de implantación por sitio y por perfil, coherente con el cronograma obligatorio del Artículo 17° de las Bases Administrativas.\newline\textit{BT; Obligatorio.} &
Sí &
Plan de Aceptación; Ensayos de Migración y Marcha Blanca &
SD7, sec. 7.2; SD8, sec. 8.1 \\ \hline

RT-\allowbreak{}20.\allowbreak{}02 &
La estrategia de paso a producción será gradual y reversible. Se declarará el criterio de avance entre olas y el procedimiento de reversión, con su tiempo de ejecución.\newline\textit{BT; Obligatorio.} &
Sí &
Plan de Aceptación; Ensayos de Migración y Marcha Blanca &
SD7, sec. 7.2; SD8, sec. 8.1 \\ \hline

RT-\allowbreak{}20.\allowbreak{}03 &
Durante la marcha blanca, la solución convivirá con la operación vigente del CLIENTE, con conciliación diaria y sin doble digitación no declarada.\newline\textit{BT; Obligatorio.} &
Sí &
Plan de Aceptación; Ensayos de Migración y Marcha Blanca &
SD7, sec. 7.2; SD8, sec. 8.1 \\ \hline

RT-\allowbreak{}20.\allowbreak{}04 &
El PROPONENTE definirá los indicadores diarios que se medirán durante la marcha blanca y sus umbrales de avance, conforme al Artículo 17.3 de las Bases Administrativas.\newline\textit{BT; Obligatorio.} &
Sí &
Plan de Aceptación; Ensayos de Migración y Marcha Blanca &
SD7, sec. 7.2; SD8, sec. 8.1 \\ \hline

RT-\allowbreak{}20.\allowbreak{}05 &
Se dispondrá de acompañamiento en terreno durante las primeras semanas de cada ola, con dotación declarada y decreciente según la curva de adopción.\newline\textit{BT; Obligatorio.} &
Sí &
Plan de Aceptación; Ensayos de Migración y Marcha Blanca &
SD7, sec. 7.2; SD8, sec. 8.1 \\ \hline

RT-\allowbreak{}20.\allowbreak{}06 &
Se establecerá un período de estabilización con atención reforzada tras cada paso a producción, con dotación y duración declaradas y sin costo adicional.\newline\textit{BT; Obligatorio.} &
Sí &
Plan de Aceptación; Ensayos de Migración y Marcha Blanca &
SD7, sec. 7.2; SD8, sec. 8.1 \\ \hline

RT-\allowbreak{}20.\allowbreak{}07 &
Existirá una definición de terminado acordada con el CLIENTE, aplicable a cada entregable, que incluya código, pruebas, documentación, seguridad y despliegue.\newline\textit{BT; Obligatorio.} &
Sí &
Plan de Aceptación; Ensayos de Migración y Marcha Blanca &
SD7, sec. 7.2; SD8, sec. 8.1 \\ \hline

RT-\allowbreak{}20.\allowbreak{}08 &
El protocolo de aceptación de cada hito se formalizará conforme al Formulario T-17, con criterios objetivos y verificables.\newline\textit{BT; Obligatorio.} &
Sí &
Plan de Aceptación; Ensayos de Migración y Marcha Blanca &
SD7, sec. 7.2; SD8, sec. 8.1 \\ \hline

\end{longtable}


\subsection*{BT, capítulo 21: MODELO DE OPERACIÓN, MANTENCIÓN Y SOPORTE}

\addcontentsline{toc}{subsection}{BT, capítulo 21: MODELO DE OPERACIÓN, MANTENCIÓN Y SOPORTE}

La Tabla~\ref{tab:t12-bt-cap-21} individualiza los requisitos técnicos de este ámbito y su localización en la propuesta.

\begin{longtable}{|L{0.09\AnchoMatrizTDoce}|L{0.385\AnchoMatrizTDoce}|L{0.065\AnchoMatrizTDoce}|L{0.18\AnchoMatrizTDoce}|L{0.28\AnchoMatrizTDoce}|}

\caption{Cumplimiento y trazabilidad: BT, capítulo 21: MODELO DE OPERACIÓN, MANTENCIÓN Y SOPORTE}\label{tab:t12-bt-cap-21}\\

\hline

\EncabezadoTabla{ID requerimiento} & \EncabezadoTabla{Descripción} & \EncabezadoTabla{Cumple\newline(Sí/No)} & \EncabezadoTabla{Componente que lo satisface} & \EncabezadoTabla{Sección de la propuesta} \\ \hline

\endfirsthead

\hline

\EncabezadoTabla{ID requerimiento} & \EncabezadoTabla{Descripción} & \EncabezadoTabla{Cumple\newline(Sí/No)} & \EncabezadoTabla{Componente que lo satisface} & \EncabezadoTabla{Sección de la propuesta} \\ \hline

\endhead

\multicolumn{5}{r}{\fontsize{9}{11}\selectfont Continúa en la página siguiente.}\\

\endfoot

\hline

\endlastfoot


RT-\allowbreak{}21.\allowbreak{}01 &
El ADJUDICATARIO dispondrá de un centro de operaciones de red con cobertura 24×7×365, propio o subcontratado, y declarará su ubicación, dotación por turno y procedimientos.\newline\textit{BT; Obligatorio.} &
Sí &
Mesa de Ayuda Synaptix Niveles 1 y 2; SLA 24x7 &
SD3, sec. 3.1.2; SD10, sec. 10.1 \\ \hline

RT-\allowbreak{}21.\allowbreak{}02 &
Se designará un gerente de servicio dedicado como contraparte permanente del CLIENTE durante la fase de Operación.\newline\textit{BT; Obligatorio.} &
Sí &
Mesa de Ayuda Synaptix Niveles 1 y 2; SLA 24x7 &
SD3, sec. 3.1.2; SD10, sec. 10.1 \\ \hline

RT-\allowbreak{}21.\allowbreak{}03 &
El equipo de operación contará con especialistas por tecnología, nominados y con dedicación declarada.\newline\textit{BT; Obligatorio.} &
Sí &
Mesa de Ayuda Synaptix Niveles 1 y 2; SLA 24x7 &
SD3, sec. 3.1.2; SD10, sec. 10.1 \\ \hline

RT-\allowbreak{}21.\allowbreak{}04 &
Los procedimientos de operación estarán documentados y serán ejecutables por el personal del CLIENTE tras la transferencia de conocimiento.\newline\textit{BT; Obligatorio.} &
Sí &
Mesa de Ayuda Synaptix Niveles 1 y 2; SLA 24x7 &
SD3, sec. 3.1.2; SD10, sec. 10.1 \\ \hline

RT-\allowbreak{}21.\allowbreak{}05 &
El PROPONENTE dispondrá de un centro de atención adecuado para soportar integralmente las consultas, propio o subcontratado con un operador de nivel acreditado.\newline\textit{BT; Obligatorio.} &
Sí &
Mesa de Ayuda Synaptix Niveles 1 y 2; SLA 24x7 &
SD3, sec. 3.1.2; SD10, sec. 10.1 \\ \hline

RT-\allowbreak{}21.\allowbreak{}06 &
El centro de atención cumplirá un tiempo de atención al usuario final de 80\% antes de 20 segundos y una resolución al primer contacto de al menos 70\%. La tasa de abandono no superará el 5\%.\newline\textit{BT; Obligatorio.} &
Sí &
Mesa de Ayuda Synaptix Niveles 1 y 2; SLA 24x7 &
SD3, sec. 3.1.2; SD10, sec. 10.1 \\ \hline

RT-\allowbreak{}21.\allowbreak{}07 &
El horario mínimo de atención será de 8:00 a 20:00 en días hábiles, ampliado a 24×7 para los incidentes de severidad crítica y para la ventana operacional que defina el caso.\newline\textit{BT; Según caso.} &
Sí &
Mesa de Ayuda Synaptix Niveles 1 y 2; SLA 24x7 &
SD3, sec. 3.1.2; SD10, sec. 10.1 \\ \hline

RT-\allowbreak{}21.\allowbreak{}08 &
La atención cubrirá orientación funcional de distintos grados de complejidad, desde preguntas simples hasta situaciones complejas, y las preguntas técnicas más frecuentes sobre la aplicación y su entorno de uso.\newline\textit{BT; Obligatorio.} &
Sí &
Mesa de Ayuda Synaptix Niveles 1 y 2; SLA 24x7 &
SD3, sec. 3.1.2; SD10, sec. 10.1 \\ \hline

RT-\allowbreak{}21.\allowbreak{}09 &
Se mantendrá un registro histórico de las actividades de soporte que permita monitorear y gestionar los principales requerimientos, con análisis de tendencia mensual.\newline\textit{BT; Obligatorio.} &
Sí &
Mesa de Ayuda Synaptix Niveles 1 y 2; SLA 24x7 &
SD3, sec. 3.1.2; SD10, sec. 10.1 \\ \hline

RT-\allowbreak{}21.\allowbreak{}10 &
Existirá un proceso de aprendizaje del centro de soporte que acumule conocimiento sobre las complejidades de la operación y las necesidades de las personas usuarias, reflejado en la base de conocimiento.\newline\textit{BT; Obligatorio.} &
Sí &
Mesa de Ayuda Synaptix Niveles 1 y 2; SLA 24x7 &
SD3, sec. 3.1.2; SD10, sec. 10.1 \\ \hline

RT-\allowbreak{}21.\allowbreak{}11 &
Se proveerán servicios de capacitación en línea en modalidad de autoformación, con registro de avance para efectos de monitoreo.\newline\textit{BT; Obligatorio.} &
Sí &
Mesa de Ayuda Synaptix Niveles 1 y 2; SLA 24x7 &
SD3, sec. 3.1.2; SD10, sec. 10.1 \\ \hline

RT-\allowbreak{}21.\allowbreak{}12 &
Se habilitarán espacios de interacción entre las entidades y personas usuarias del sistema, que permitan compartir experiencias de uso.\newline\textit{BT; Deseable.} &
No &
No aplica / Descartado: Canal de soporte unificado mediante Mesa de Ayuda Synaptix e ITSM; sin foro comunitario abierto entre clientes &
SD3, sec. 3.2.2; SD10, sec. 10.1 \\ \hline

RT-\allowbreak{}21.\allowbreak{}13 &
Los indicadores clave del proceso de soporte estarán abiertos al CLIENTE, y los indicadores específicos serán accesibles según nivel de autorización.\newline\textit{BT; Obligatorio.} &
Sí &
Mesa de Ayuda Synaptix Niveles 1 y 2; SLA 24x7 &
SD3, sec. 3.1.2; SD10, sec. 10.1 \\ \hline

RT-\allowbreak{}21.\allowbreak{}14 &
El PROPONENTE dimensionará la dotación del centro de atención con fundamento cuantitativo, empleando teoría de colas o el modelo Erlang C, a partir del volumen de contactos proyectado del caso.\newline\textit{BT; Obligatorio.} &
Sí &
Mesa de Ayuda Synaptix Niveles 1 y 2; SLA 24x7 &
SD3, sec. 3.1.2; SD10, sec. 10.1 \\ \hline

RT-\allowbreak{}21.\allowbreak{}15 &
Existirá un canal único de registro de incidentes y solicitudes, con número de ticket, clasificación por severidad y seguimiento del ciclo de vida completo hasta el cierre conforme.\newline\textit{BT; Obligatorio.} &
Sí &
Mesa de Ayuda Synaptix Niveles 1 y 2; SLA 24x7 &
SD3, sec. 3.1.2; SD10, sec. 10.1 \\ \hline

RT-\allowbreak{}21.\allowbreak{}16 &
Cuando el caso comprenda sitios alejados, los especialistas de niveles 2 y 3 deberán trasladarse cuando la resolución lo requiera, por el medio más rápido disponible y con disponibilidad de tiempo suficiente, sin alterar la atención normal del resto de los sitios. El costo del traslado está incluido en la oferta.\newline\textit{BT; Según caso.} &
Sí &
Mesa de Ayuda Synaptix Niveles 1 y 2; SLA 24x7 &
SD3, sec. 3.1.2; SD10, sec. 10.1 \\ \hline

RT-\allowbreak{}21.\allowbreak{}17 &
El cierre de un ticket requerirá confirmación de la persona usuaria o transcurso del plazo de confirmación automática declarado.\newline\textit{BT; Obligatorio.} &
Sí &
Mesa de Ayuda Synaptix Niveles 1 y 2; SLA 24x7 &
SD3, sec. 3.1.2; SD10, sec. 10.1 \\ \hline

RT-\allowbreak{}21.\allowbreak{}18 &
El ADJUDICATARIO reportará mensualmente el cumplimiento de los niveles de servicio de atención, con el detalle por severidad y el análisis de los incumplimientos.\newline\textit{BT; Obligatorio.} &
Sí &
Mesa de Ayuda Synaptix Niveles 1 y 2; SLA 24x7 &
SD3, sec. 3.1.2; SD10, sec. 10.1 \\ \hline

RT-\allowbreak{}21.\allowbreak{}19 &
El PROPONENTE declarará el tamaño de la bolsa anual de horas de mantención evolutiva, su composición por perfil y el procedimiento de solicitud, estimación, aprobación y liquidación.\newline\textit{BT; Obligatorio.} &
Sí &
Mesa de Ayuda Synaptix Niveles 1 y 2; SLA 24x7 &
SD3, sec. 3.1.2; SD10, sec. 10.1 \\ \hline

RT-\allowbreak{}21.\allowbreak{}20 &
La mantención evolutiva se someterá al mismo estándar de calidad, pruebas y seguridad que el desarrollo original.\newline\textit{BT; Obligatorio.} &
Sí &
Mesa de Ayuda Synaptix Niveles 1 y 2; SLA 24x7 &
SD3, sec. 3.1.2; SD10, sec. 10.1 \\ \hline

RT-\allowbreak{}21.\allowbreak{}21 &
El PROPONENTE mantendrá un registro de deuda técnica cuantificado y destinará una fracción declarada de la capacidad de la Operación a reducirla.\newline\textit{BT; Obligatorio.} &
Sí &
Mesa de Ayuda Synaptix Niveles 1 y 2; SLA 24x7 &
SD3, sec. 3.1.2; SD10, sec. 10.1 \\ \hline

RT-\allowbreak{}21.\allowbreak{}22 &
Las actualizaciones de versión de los componentes de base se planificarán anualmente, con ventana acordada y plan de reversión.\newline\textit{BT; Obligatorio.} &
Sí &
Mesa de Ayuda Synaptix Niveles 1 y 2; SLA 24x7 &
SD3, sec. 3.1.2; SD10, sec. 10.1 \\ \hline

\end{longtable}


\subsection*{BT, capítulo 22: CAPACITACIÓN Y TRANSFERENCIA DE CONOCIMIENTO}

\addcontentsline{toc}{subsection}{BT, capítulo 22: CAPACITACIÓN Y TRANSFERENCIA DE CONOCIMIENTO}

La Tabla~\ref{tab:t12-bt-cap-22} individualiza los requisitos técnicos de este ámbito y su localización en la propuesta.

\begin{longtable}{|L{0.09\AnchoMatrizTDoce}|L{0.385\AnchoMatrizTDoce}|L{0.065\AnchoMatrizTDoce}|L{0.18\AnchoMatrizTDoce}|L{0.28\AnchoMatrizTDoce}|}

\caption{Cumplimiento y trazabilidad: BT, capítulo 22: CAPACITACIÓN Y TRANSFERENCIA DE CONOCIMIENTO}\label{tab:t12-bt-cap-22}\\

\hline

\EncabezadoTabla{ID requerimiento} & \EncabezadoTabla{Descripción} & \EncabezadoTabla{Cumple\newline(Sí/No)} & \EncabezadoTabla{Componente que lo satisface} & \EncabezadoTabla{Sección de la propuesta} \\ \hline

\endfirsthead

\hline

\EncabezadoTabla{ID requerimiento} & \EncabezadoTabla{Descripción} & \EncabezadoTabla{Cumple\newline(Sí/No)} & \EncabezadoTabla{Componente que lo satisface} & \EncabezadoTabla{Sección de la propuesta} \\ \hline

\endhead

\multicolumn{5}{r}{\fontsize{9}{11}\selectfont Continúa en la página siguiente.}\\

\endfoot

\hline

\endlastfoot


RT-\allowbreak{}22.\allowbreak{}01 &
El plan de capacitación se estructurará por perfil: personas usuarias finales, usuarias avanzadas, administradoras, equipo técnico y soporte de niveles 1 y 2.\newline\textit{BT; Obligatorio.} &
Sí &
Plan de Capacitación por Perfiles y Manuales &
SD10, sec. 10.2 \\ \hline

RT-\allowbreak{}22.\allowbreak{}02 &
Las modalidades incluirán capacitación presencial en cada sitio de operación, sesiones en línea sincrónicas, autoformación y acompañamiento en puesto de trabajo durante la marcha blanca.\newline\textit{BT; Obligatorio.} &
Sí &
Plan de Capacitación por Perfiles y Manuales &
SD10, sec. 10.2 \\ \hline

RT-\allowbreak{}22.\allowbreak{}03 &
Todo el material de capacitación se entregará en español, en formato editable y de propiedad del CLIENTE: manuales por perfil, guías rápidas, preguntas frecuentes, videos tutoriales y base de conocimiento consultable.\newline\textit{BT; Obligatorio.} &
Sí &
Plan de Capacitación por Perfiles y Manuales &
SD10, sec. 10.2 \\ \hline

RT-\allowbreak{}22.\allowbreak{}04 &
La capacitación no podrá afectar la operación del CLIENTE: se programará por turnos y en horarios acordados, considerando la estacionalidad y la ventana operacional del caso.\newline\textit{BT; Según caso.} &
Sí &
Plan de Capacitación por Perfiles y Manuales &
SD10, sec. 10.2 \\ \hline

RT-\allowbreak{}22.\allowbreak{}05 &
Las personas usuarias administradoras y el equipo técnico del CLIENTE serán evaluados y certificados como condición para el cierre de cada marcha blanca.\newline\textit{BT; Obligatorio.} &
Sí &
Plan de Capacitación por Perfiles y Manuales &
SD10, sec. 10.2 \\ \hline

RT-\allowbreak{}22.\allowbreak{}06 &
Durante la Operación se ejecutarán al menos dos jornadas anuales de actualización y se capacitará al personal nuevo del CLIENTE sin costo adicional.\newline\textit{BT; Obligatorio.} &
Sí &
Plan de Capacitación por Perfiles y Manuales &
SD10, sec. 10.2 \\ \hline

RT-\allowbreak{}22.\allowbreak{}07 &
El ADJUDICATARIO proveerá acompañamiento y mentoría al equipo técnico del CLIENTE durante al menos seis meses posteriores al paso a producción de la Etapa 2.\newline\textit{BT; Obligatorio.} &
Sí &
Plan de Capacitación por Perfiles y Manuales &
SD10, sec. 10.2 \\ \hline

RT-\allowbreak{}22.\allowbreak{}08 &
La base de conocimiento será mantenida y actualizada durante todo el Contrato, con métricas de uso y de utilidad percibida.\newline\textit{BT; Obligatorio.} &
Sí &
Plan de Capacitación por Perfiles y Manuales &
SD10, sec. 10.2 \\ \hline

RT-\allowbreak{}22.\allowbreak{}09 &
Se valorará la existencia de un ambiente permanente de entrenamiento, con datos ficticios, disponible para la práctica sin riesgo.\newline\textit{BT; Deseable.} &
No &
No aplica / Descartado: Capacitación presencial y guiada en noviembre antes del congelamiento estival (RST-TEM-01); no requiere plataforma LMS permanente &
SD3, sec. 3.2.2; SD10, sec. 10.2 \\ \hline

\end{longtable}


\subsection*{BT, capítulo 23: INFORMACIÓN CORPORATIVA Y PRESENCIA DIGITAL}

\addcontentsline{toc}{subsection}{BT, capítulo 23: INFORMACIÓN CORPORATIVA Y PRESENCIA DIGITAL}

La Tabla~\ref{tab:t12-bt-cap-23} individualiza los requisitos técnicos de este ámbito y su localización en la propuesta.

\begin{longtable}{|L{0.09\AnchoMatrizTDoce}|L{0.385\AnchoMatrizTDoce}|L{0.065\AnchoMatrizTDoce}|L{0.18\AnchoMatrizTDoce}|L{0.28\AnchoMatrizTDoce}|}

\caption{Cumplimiento y trazabilidad: BT, capítulo 23: INFORMACIÓN CORPORATIVA Y PRESENCIA DIGITAL}\label{tab:t12-bt-cap-23}\\

\hline

\EncabezadoTabla{ID requerimiento} & \EncabezadoTabla{Descripción} & \EncabezadoTabla{Cumple\newline(Sí/No)} & \EncabezadoTabla{Componente que lo satisface} & \EncabezadoTabla{Sección de la propuesta} \\ \hline

\endfirsthead

\hline

\EncabezadoTabla{ID requerimiento} & \EncabezadoTabla{Descripción} & \EncabezadoTabla{Cumple\newline(Sí/No)} & \EncabezadoTabla{Componente que lo satisface} & \EncabezadoTabla{Sección de la propuesta} \\ \hline

\endhead

\multicolumn{5}{r}{\fontsize{9}{11}\selectfont Continúa en la página siguiente.}\\

\endfoot

\hline

\endlastfoot


RT-\allowbreak{}23.\allowbreak{}01 &
Información institucional: descripción detallada del giro principal, historia y trayectoria de la empresa con al menos tres años de antigüedad comprobable, misión, visión y valores, certificaciones y acreditaciones vigentes, y presencia geográfica y oficinas.\newline\textit{BT; Obligatorio.} &
Sí &
Perfil Corporativo Synaptix; Experiencia Comprobada &
SD3, sec. 3.1; SD12, sec. 12.1 \\ \hline

RT-\allowbreak{}23.\allowbreak{}02 &
Experiencia y casos de éxito: portafolio de proyectos similares de los últimos cinco años, casos documentados con métricas verificables, testimonios de clientes con autorización de publicación, industrias atendidas con énfasis en la del caso asignado, y volumen de transacciones o de personas usuarias gestionadas.\newline\textit{BT; Obligatorio.} &
Sí &
Perfil Corporativo Synaptix; Experiencia Comprobada &
SD3, sec. 3.1; SD12, sec. 12.1 \\ \hline

RT-\allowbreak{}23.\allowbreak{}03 &
Equipo profesional: organigrama del equipo directivo, perfiles de socios y directores, currículos resumidos del equipo técnico clave y certificaciones profesionales del personal.\newline\textit{BT; Obligatorio.} &
Sí &
Perfil Corporativo Synaptix; Experiencia Comprobada &
SD3, sec. 3.1; SD12, sec. 12.1 \\ \hline

RT-\allowbreak{}23.\allowbreak{}04 &
Capacidades técnicas: servicios y soluciones ofrecidas, conjunto tecnológico dominado, alianzas tecnológicas con proveedores de nube y de plataforma, metodologías de trabajo certificadas, e infraestructura y capacidad instalada.\newline\textit{BT; Obligatorio.} &
Sí &
Perfil Corporativo Synaptix; Experiencia Comprobada &
SD3, sec. 3.1; SD12, sec. 12.1 \\ \hline

RT-\allowbreak{}23.\allowbreak{}05 &
El sitio será accesible conforme a WCAG 2.2 nivel AA, responsivo y con certificado TLS válido y vigente.\newline\textit{BT; Obligatorio.} &
Sí &
Perfil Corporativo Synaptix; Experiencia Comprobada &
SD3, sec. 3.1; SD12, sec. 12.1 \\ \hline

RT-\allowbreak{}23.\allowbreak{}06 &
El sitio se mantendrá activo y disponible durante todo el proceso de licitación, desde el registro del participante hasta la adjudicación.\newline\textit{BT; Obligatorio.} &
Sí &
Perfil Corporativo Synaptix; Experiencia Comprobada &
SD3, sec. 3.1; SD12, sec. 12.1 \\ \hline

RT-\allowbreak{}23.\allowbreak{}07 &
Material educativo digital: seminarios en línea, artículos técnicos o centro de recursos y documentación.\newline\textit{BT; Deseable.} &
No &
No aplica / Descartado: Material formativo enfocado exclusivamente en los roles operacionales del puerto; no incluye seminarios públicos corporativos &
SD3, sec. 3.2.2; SD10, sec. 10.2 \\ \hline

RT-\allowbreak{}23.\allowbreak{}08 &
Demostración en línea, recorrido virtual de la solución o portal de soporte para clientes.\newline\textit{BT; Deseable.} &
No &
No aplica / Descartado: Recorrido y validación provistos mediante Prototipo Interactivo Figma oficial (Subdocumento 12); no requiere showroom virtual 3D &
SD3, sec. 3.2.2; SD12, sec. 12.3 \\ \hline

RT-\allowbreak{}23.\allowbreak{}09 &
Métricas de disponibilidad y desempeño de los servicios del proponente publicadas en tiempo real.\newline\textit{BT; Deseable.} &
No &
No aplica / Descartado: Métricas de servicio expuestas en consola privada de administración para la gerencia; no en sitio web público de marketing &
SD3, sec. 3.2.2; SD4, sec. 4.1 \\ \hline

RT-\allowbreak{}23.\allowbreak{}10 &
Calculadora de retorno de la inversión o herramienta de estimación pertinente a la industria del caso.\newline\textit{BT; Deseable.} &
No &
No aplica / Descartado: Propuesta económica y modelo de amortización cerrados en Formulario T-2; no aplica calculadora de ROI pública &
SD3, sec. 3.2.2; SD1, sec. 1.2 \\ \hline

\end{longtable}


\subsection*{BT, capítulo 24: VIDEO DE PRESENTACIÓN DE LA PROPUESTA}

\addcontentsline{toc}{subsection}{BT, capítulo 24: VIDEO DE PRESENTACIÓN DE LA PROPUESTA}

La Tabla~\ref{tab:t12-bt-cap-24} individualiza los requisitos técnicos de este ámbito y su localización en la propuesta.

\begin{longtable}{|L{0.09\AnchoMatrizTDoce}|L{0.385\AnchoMatrizTDoce}|L{0.065\AnchoMatrizTDoce}|L{0.18\AnchoMatrizTDoce}|L{0.28\AnchoMatrizTDoce}|}

\caption{Cumplimiento y trazabilidad: BT, capítulo 24: VIDEO DE PRESENTACIÓN DE LA PROPUESTA}\label{tab:t12-bt-cap-24}\\

\hline

\EncabezadoTabla{ID requerimiento} & \EncabezadoTabla{Descripción} & \EncabezadoTabla{Cumple\newline(Sí/No)} & \EncabezadoTabla{Componente que lo satisface} & \EncabezadoTabla{Sección de la propuesta} \\ \hline

\endfirsthead

\hline

\EncabezadoTabla{ID requerimiento} & \EncabezadoTabla{Descripción} & \EncabezadoTabla{Cumple\newline(Sí/No)} & \EncabezadoTabla{Componente que lo satisface} & \EncabezadoTabla{Sección de la propuesta} \\ \hline

\endhead

\multicolumn{5}{r}{\fontsize{9}{11}\selectfont Continúa en la página siguiente.}\\

\endfoot

\hline

\endlastfoot


RT-\allowbreak{}24.\allowbreak{}01 &
Todos los integrantes clave del equipo deberán aparecer individualmente presentándose, en toma de medio cuerpo, mirando directamente a la cámara, con vestimenta formal o de negocio informal, y con una duración mínima de aparición de 10 segundos por persona.\newline\textit{BT; Obligatorio.} &
Sí &
Video de Presentación de la Propuesta Synaptix &
SD12, sec. 12.2 \\ \hline

RT-\allowbreak{}24.\allowbreak{}02 &
Cada vez que aparezca un integrante deberá mostrarse en pantalla su nombre completo, su cargo en el proyecto, sus años de experiencia y su especialización relevante.\newline\textit{BT; Obligatorio.} &
Sí &
Video de Presentación de la Propuesta Synaptix &
SD12, sec. 12.2 \\ \hline

RT-\allowbreak{}24.\allowbreak{}03 &
Las certificaciones principales de cada integrante podrán mostrarse en pantalla.\newline\textit{BT; Deseable.} &
No &
No aplica / Descartado: Equipo de trabajo y certificaciones profesionales acreditados formalmente en Subdocumento 1 y Anexo T-5; no en video promocional &
SD3, sec. 3.2.2; SD1, sec. 1.1 \\ \hline

RT-\allowbreak{}24.\allowbreak{}04 &
El fondo será el mismo para todas las tomas de personas, con iluminación consistente y encuadre similar para todos los participantes.\newline\textit{BT; Obligatorio.} &
Sí &
Video de Presentación de la Propuesta Synaptix &
SD12, sec. 12.2 \\ \hline

RT-\allowbreak{}24.\allowbreak{}05 &
La paleta de colores corporativa, la tipografía y el uso de la marca serán consistentes en todo el video, con logotipo visible en todas las escenas y datos de contacto en el encabezado o pie.\newline\textit{BT; Obligatorio.} &
Sí &
Video de Presentación de la Propuesta Synaptix &
SD12, sec. 12.2 \\ \hline

RT-\allowbreak{}24.\allowbreak{}06 &
Los niveles de volumen estarán normalizados, con la misma calidad de grabación y sin variaciones bruscas de sonido.\newline\textit{BT; Obligatorio.} &
Sí &
Video de Presentación de la Propuesta Synaptix &
SD12, sec. 12.2 \\ \hline

\end{longtable}


\subsection*{BT, capítulo 25: PROTOTIPO INTERACTIVO DE INTERFAZ Y DISEÑO UX/UI}

\addcontentsline{toc}{subsection}{BT, capítulo 25: PROTOTIPO INTERACTIVO DE INTERFAZ Y DISEÑO UX/UI}

La Tabla~\ref{tab:t12-bt-cap-25} individualiza los requisitos técnicos de este ámbito y su localización en la propuesta.

\begin{longtable}{|L{0.09\AnchoMatrizTDoce}|L{0.385\AnchoMatrizTDoce}|L{0.065\AnchoMatrizTDoce}|L{0.18\AnchoMatrizTDoce}|L{0.28\AnchoMatrizTDoce}|}

\caption{Cumplimiento y trazabilidad: BT, capítulo 25: PROTOTIPO INTERACTIVO DE INTERFAZ Y DISEÑO UX/UI}\label{tab:t12-bt-cap-25}\\

\hline

\EncabezadoTabla{ID requerimiento} & \EncabezadoTabla{Descripción} & \EncabezadoTabla{Cumple\newline(Sí/No)} & \EncabezadoTabla{Componente que lo satisface} & \EncabezadoTabla{Sección de la propuesta} \\ \hline

\endfirsthead

\hline

\EncabezadoTabla{ID requerimiento} & \EncabezadoTabla{Descripción} & \EncabezadoTabla{Cumple\newline(Sí/No)} & \EncabezadoTabla{Componente que lo satisface} & \EncabezadoTabla{Sección de la propuesta} \\ \hline

\endhead

\multicolumn{5}{r}{\fontsize{9}{11}\selectfont Continúa en la página siguiente.}\\

\endfoot

\hline

\endlastfoot


RT-\allowbreak{}25.\allowbreak{}01 &
Navegación intuitiva sin necesidad de manual, con un máximo de tres interacciones para alcanzar cualquier funcionalidad principal.\newline\textit{BT; Obligatorio.} &
Sí &
Prototipo Interactivo Figma y Panitao UI &
SD4, sec. 4.1; SD12, sec. 12.3 \\ \hline

RT-\allowbreak{}25.\allowbreak{}02 &
Retroalimentación visual clara ante cada acción, estados de carga y transición fluidos, y manejo elegante de los errores.\newline\textit{BT; Obligatorio.} &
Sí &
Prototipo Interactivo Figma y Panitao UI &
SD4, sec. 4.1; SD12, sec. 12.3 \\ \hline

RT-\allowbreak{}25.\allowbreak{}03 &
Cumplimiento de WCAG 2.2 nivel AA: contraste adecuado, tamaños de fuente legibles, objetivos táctiles de al menos 44 × 44 píxeles, navegación completa por teclado, y textos alternativos y etiquetas de accesibilidad.\newline\textit{BT; Obligatorio.} &
Sí &
Prototipo Interactivo Figma y Panitao UI &
SD4, sec. 4.1; SD12, sec. 12.3 \\ \hline

RT-\allowbreak{}25.\allowbreak{}04 &
Diseño responsivo demostrado en escritorio (1920 × 1080 y 1366 × 768), tableta (768 × 1024 en vertical y horizontal) y teléfono (375 × 812), con puntos de quiebre coherentes.\newline\textit{BT; Obligatorio.} &
Sí &
Prototipo Interactivo Figma y Panitao UI &
SD4, sec. 4.1; SD12, sec. 12.3 \\ \hline

RT-\allowbreak{}25.\allowbreak{}05 &
Sistema de diseño documentado: paleta de a lo más cinco colores principales, a lo más dos familias tipográficas jerarquizadas, iconografía coherente, retícula y espaciado consistentes, y componentes reutilizables.\newline\textit{BT; Obligatorio.} &
Sí &
Prototipo Interactivo Figma y Panitao UI &
SD4, sec. 4.1; SD12, sec. 12.3 \\ \hline

RT-\allowbreak{}25.\allowbreak{}06 &
El prototipo se entregará como enlace navegable en línea, sin instalación ni complementos, compatible con los navegadores modernos vigentes.\newline\textit{BT; Obligatorio.} &
Sí &
Prototipo Interactivo Figma y Panitao UI &
SD4, sec. 4.1; SD12, sec. 12.3 \\ \hline

RT-\allowbreak{}25.\allowbreak{}07 &
El acceso al prototipo estará garantizado por un mínimo de seis meses desde su entrega.\newline\textit{BT; Obligatorio.} &
Sí &
Prototipo Interactivo Figma y Panitao UI &
SD4, sec. 4.1; SD12, sec. 12.3 \\ \hline

RT-\allowbreak{}25.\allowbreak{}08 &
El prototipo permitirá navegación completa entre pantallas, simulación de ingreso de datos, transiciones y animaciones básicas, estados de posado y activo, simulación de carga de archivos, elementos desplegables funcionales y simulación de validaciones.\newline\textit{BT; Obligatorio.} &
Sí &
Prototipo Interactivo Figma y Panitao UI &
SD4, sec. 4.1; SD12, sec. 12.3 \\ \hline

RT-\allowbreak{}25.\allowbreak{}09 &
Se entregará la arquitectura de información del prototipo: mapa del sitio completo, diagrama de navegación, taxonomía y nomenclatura, estructura de menús y jerarquía de la información.\newline\textit{BT; Obligatorio.} &
Sí &
Prototipo Interactivo Figma y Panitao UI &
SD4, sec. 4.1; SD12, sec. 12.3 \\ \hline

RT-\allowbreak{}25.\allowbreak{}10 &
Se valorará la posibilidad de exportar el prototipo a un documento interactivo para revisión sin conexión.\newline\textit{BT; Deseable.} &
No &
No aplica / Descartado: Prototipo navegable entregado en enlace Figma activo e interactivo conforme a Bases Técnicas; sin exportación a PDF interactivo &
SD3, sec. 3.2.2; SD12, sec. 12.3 \\ \hline

\end{longtable}


\subsection*{BT, capítulo 26: INNOVACIONES}

\addcontentsline{toc}{subsection}{BT, capítulo 26: INNOVACIONES}

La Tabla~\ref{tab:t12-bt-cap-26} individualiza los requisitos técnicos de este ámbito y su localización en la propuesta.

\begin{longtable}{|L{0.09\AnchoMatrizTDoce}|L{0.385\AnchoMatrizTDoce}|L{0.065\AnchoMatrizTDoce}|L{0.18\AnchoMatrizTDoce}|L{0.28\AnchoMatrizTDoce}|}

\caption{Cumplimiento y trazabilidad: BT, capítulo 26: INNOVACIONES}\label{tab:t12-bt-cap-26}\\

\hline

\EncabezadoTabla{ID requerimiento} & \EncabezadoTabla{Descripción} & \EncabezadoTabla{Cumple\newline(Sí/No)} & \EncabezadoTabla{Componente que lo satisface} & \EncabezadoTabla{Sección de la propuesta} \\ \hline

\endfirsthead

\hline

\EncabezadoTabla{ID requerimiento} & \EncabezadoTabla{Descripción} & \EncabezadoTabla{Cumple\newline(Sí/No)} & \EncabezadoTabla{Componente que lo satisface} & \EncabezadoTabla{Sección de la propuesta} \\ \hline

\endhead

\multicolumn{5}{r}{\fontsize{9}{11}\selectfont Continúa en la página siguiente.}\\

\endfoot

\hline

\endlastfoot


RT-\allowbreak{}26.\allowbreak{}01 &
Cada innovación se ubicará explícitamente en la arquitectura: qué capa la contiene, qué componentes la implementan y qué interfaces consume o expone.\newline\textit{BT; Obligatorio.} &
Sí &
Cinco Innovaciones Tecnológicas Art. 28° &
SD13, sec. 13.1 a 13.5 \\ \hline

RT-\allowbreak{}26.\allowbreak{}02 &
Cada innovación identificará los paquetes de la estructura de descomposición del trabajo que la ejecutan y el mes del cronograma en que se materializa.\newline\textit{BT; Obligatorio.} &
Sí &
Cinco Innovaciones Tecnológicas Art. 28° &
SD13, sec. 13.1 a 13.5 \\ \hline

RT-\allowbreak{}26.\allowbreak{}03 &
Las innovaciones de base tecnológica declararán el nivel de madurez de la tecnología con la escala utilizada y citarán las fuentes en norma APA 7.ª edición.\newline\textit{BT; Obligatorio.} &
Sí &
Cinco Innovaciones Tecnológicas Art. 28° &
SD13, sec. 13.1 a 13.5 \\ \hline

RT-\allowbreak{}26.\allowbreak{}04 &
Cada innovación declarará su riesgo de adopción, su probabilidad, su impacto, la estrategia de mitigación y el plan de contingencia si no rinde lo esperado.\newline\textit{BT; Obligatorio.} &
Sí &
Cinco Innovaciones Tecnológicas Art. 28° &
SD13, sec. 13.1 a 13.5 \\ \hline

RT-\allowbreak{}26.\allowbreak{}05 &
Cada innovación declarará su indicador de verificación con línea base, meta y momento de medición, y su impacto en inversión, costo operacional y beneficio esperado.\newline\textit{BT; Obligatorio.} &
Sí &
Cinco Innovaciones Tecnológicas Art. 28° &
SD13, sec. 13.1 a 13.5 \\ \hline

RT-\allowbreak{}26.\allowbreak{}06 &
Las innovaciones que incorporen inteligencia artificial cumplirán íntegramente el Capítulo 18 de este documento.\newline\textit{BT; Obligatorio.} &
Sí &
Cinco Innovaciones Tecnológicas Art. 28° &
SD13, sec. 13.1 a 13.5 \\ \hline

RT-\allowbreak{}26.\allowbreak{}07 &
Las innovaciones que modifiquen la arquitectura de seguridad requerirán su propio modelado de amenazas.\newline\textit{BT; Obligatorio.} &
Sí &
Cinco Innovaciones Tecnológicas Art. 28° &
SD13, sec. 13.1 a 13.5 \\ \hline

RT-\allowbreak{}26.\allowbreak{}08 &
Se valorará que al menos una innovación sea verificable durante la marcha blanca de la Etapa 1, es decir, que su beneficio pueda medirse antes del mes 16.\newline\textit{BT; Deseable.} &
Sí &
Innovación 1 (Custodia Náutica Preventiva) operativa y medible en la marcha blanca del mes 13 al 15 &
SD13, sec. 13.1; SD7, sec. 7.3 \\ \hline

\end{longtable}


\end{paginaHorizontalSynaptix}
